%% Alternation, Invariants, and the Limits of Collapse -- revised version, September 30, 2026.
%% Single-file version generated from main_siam.tex and sections/*.tex.
\documentclass[final]{siamart0516}
% SIAM class siamart0516 (May 2016 macros; obtain the full distribution from SIAM).
% The class loads amsmath, ntheorem, hyperref, xcolor and cleveref itself.
\usepackage{amssymb,mathtools}
\usepackage{enumitem}
\usepackage{booktabs,longtable,tabularx,array}
\usepackage{tikz}
\usepackage{adjustbox}
\usetikzlibrary{positioning,arrows.meta,shapes.misc}

\setlist{itemsep=2pt,topsep=4pt}
\setlength{\emergencystretch}{2.5em}
\newcolumntype{L}[1]{>{\raggedright\arraybackslash}p{#1\linewidth}}
\setlength{\LTpre}{6pt}\setlength{\LTpost}{6pt}

% ------------------------------------------------------------------ theorem environments
\numberwithin{theorem}{section}
\numberwithin{equation}{section}
\newsiamthm{conjecture}{Conjecture}
\newsiamthm{claim}{Claim}
\newsiamremark{remark}{Remark}
\newsiamremark{example}{Example}
\newsiamremark{problem}{Open Problem}
% Unnumbered remark* and definition* (ntheorem's own starred variants are numbered by this class).
\makeatletter
\newtheoremstyle{siamnonum}%
 {\item[\hskip\labelsep\hskip\parindent\theorem@headerfont ##1\theorem@separator]}%
 {\item[\hskip\labelsep\hskip\parindent\theorem@headerfont ##1\ \textup{(##3)}\theorem@separator]}
\makeatother
\theoremstyle{siamnonum}
\theoremheaderfont{\normalfont\itshape}
\theorembodyfont{\normalfont}
\theoremseparator{.}
\theoremsymbol{}
\newtheorem{remarkstar}{Remark}
\theoremheaderfont{\normalfont\sc}
\theorembodyfont{\normalfont\itshape}
\newtheorem{definitionstar}{Definition}
\renewenvironment{remark*}{\remarkstar}{\endremarkstar}
\renewenvironment{definition*}{\definitionstar}{\enddefinitionstar}
\theoremstyle{plain}
\crefname{conjecture}{Conjecture}{Conjectures}
\crefname{claim}{Claim}{Claims}
\crefname{problem}{Open Problem}{Open Problems}
\crefname{remark}{Remark}{Remarks}
\crefname{example}{Example}{Examples}
\providecommand{\qed}{\hfill\proofbox}

% ------------------------------------------------------------------ macros
\newcommand{\code}[1]{{\ttfamily\nolinkurl{#1}}}
\makeatletter\g@addto@macro\UrlBreaks{\do\_\do\-}\makeatother
\newcommand{\NN}{\mathbb{N}}
\newcommand{\ZZ}{\mathbb{Z}}
\newcommand{\QQ}{\mathbb{Q}}
\newcommand{\RR}{\mathbb{R}}
\newcommand{\CC}{\mathbb{C}}
\newcommand{\FF}{\mathbb{F}}
\newcommand{\Qz}{\QQ(z)}
\newcommand{\Rz}{\RR(z)}
\newcommand{\Qzbar}{\overline{\QQ(z)}}
\newcommand{\Puis}{\mathcal{P}_{\RR}}
\newcommand{\val}{\mathcal{V}}            % the value
\newcommand{\pay}{\mathcal{U}}            % the payoff
\newcommand{\SG}{\mathcal{G}}             % the state graph
\newcommand{\arena}{\mathcal{A}}          % an arena
\newcommand{\Mpot}{\mathcal{M}}           % set-potential values
\newcommand{\cont}{\mathcal{K}}           % continuants
\newcommand{\game}{\mathfrak{G}}          % the game G(H;A)
\newcommand{\arcs}{\mathsf{A}}            % a set of one-way arcs
\newcommand{\mm}{\mathfrak{m}}            % number of maximum matchings
\newcommand{\Dv}{\mathsf{D}}              % divisor map
\newcommand{\ord}{\operatorname{ord}}
\newcommand{\Res}{\operatorname{Res}}
\newcommand{\back}{\operatorname{back}}
\newcommand{\tr}{\operatorname{tr}}
\renewcommand{\Im}{\operatorname{Im}}
\newcommand{\rk}{\operatorname{rk}}
\newcommand{\poly}{\operatorname{poly}}
\newcommand{\Sym}{\operatorname{Sym}}
\newcommand{\Rdown}{R^{\downarrow}}
\newcommand{\nplus}{N^{+}}
\newcommand{\cls}[1]{\mathsf{#1}}
\newcommand{\PSPACE}{\cls{PSPACE}}
\newcommand{\EXP}{\cls{EXP}}
\newcommand{\EXPTIME}{\cls{EXPTIME}}
\newcommand{\NP}{\cls{NP}}
\newcommand{\coNP}{\cls{coNP}}
\newcommand{\BPP}{\cls{BPP}}
\newcommand{\BQP}{\cls{BQP}}
\newcommand{\Pclass}{\cls{P}}
\newcommand{\NC}{\cls{NC}}
\newcommand{\CH}{\cls{CH}}
\newcommand{\MA}{\cls{MA}}
\newcommand{\IP}{\cls{IP}}
\newcommand{\sharpP}{\#\cls{P}}
\newcommand{\Nw}{\mathrm{N}}   % N-position
\newcommand{\Pw}{\mathrm{P}}   % P-position
\newcommand{\defeq}{\coloneqq}
\newcommand{\one}{\mathbf{1}}

\begin{document}
\newcommand{\TheTitle}{Alternation, Invariants, and the Limits of Collapse}
\newcommand{\TheAuthors}{Ruslan Baiandin}
\headers{Alternation, invariants, and the limits of collapse}{\TheAuthors}
\title{{\TheTitle}\thanks{Revised version of September 30, 2026. The verification archive \code{limits-of-collapse} accompanies this paper.}}
\author{Ruslan Baiandin}\thanks{Graduate School of Economics and Management, Ural Federal University, Yekaterinburg, 620002, Russia (ruslanbayandin47@gmail.com).}
\ifpdf
\hypersetup{pdftitle={\TheTitle},pdfauthor={\TheAuthors}}
\fi
\maketitle

\begin{abstract}
For succinctly described two-player games the minimax value has no invariant computable in deterministic polynomial time when play can be exponentially long; for play of polynomial length such an invariant exists exactly when $\Pclass=\PSPACE$; and an evaluator that sees payoffs only through queries needs exponentially many queries in the depth. We restate these classical limits in one model and then study what collapses exactly. The value is a signed spectral index of the unfolded game tree. For undirected vertex geography the index is a ratio of matching polynomials of the input graph, and for directed vertex geography it is a ratio of weighted matching polynomials of one fixed graph whenever every strong component is symmetric. We prove the converse for real weights: a digraph admits such a local weighted matching potential with weights in $\Qz$, in $\Rz$ or in real Puiseux series if and only if every strong component is symmetric. The proof applies Br\"and\'en's criterion for real stability to first order at $z=\infty$, and it also excludes Hermitian determinantal potentials and frozen environments. Over the algebraic closure of $\Qz$ the statement is false: a directed $m$-cycle together with $m+1$ isolated vertices carries a potential whose environment weights are the roots of an explicit polynomial, with full symmetric monodromy for $m\le6$. We also decide directed geography in polynomial time when the one-way arcs of every strong component share a tail and show that matching data decide no other two-arc geometry; prove that every two-periodic blow-up of a directed cycle carries a set potential; and show that exact resolvents on digraphs of small width need exponential degree and non-local memory, while circuits over component states have linear size.

\end{abstract}

\begin{keywords}
minimax, combinatorial games, geography, matching polynomial, real stable polynomials, Heilmann--Lieb theorem, Rayleigh differences, spectral collapse, alternation
\end{keywords}

\begin{AMS}
91A46, 68Q17, 68Q15, 05C50, 05C57, 05C70, 32A60
\end{AMS}

\begin{figure}[t]
\centering
\begin{adjustbox}{max width=\linewidth}
\begin{tikzpicture}[
  node distance=6mm and 4mm,
  q/.style={draw=black!45, rounded corners=3pt, minimum width=34mm, minimum height=8mm, align=center, font=\small\bfseries},
  yes/.style={draw=green!55!black, fill=green!8, rounded corners=3pt, line width=0.9pt, text width=46mm, align=center, font=\scriptsize},
  no/.style={draw=red!65!black, fill=red!6, rounded corners=3pt, line width=0.9pt, text width=38mm, align=center, font=\scriptsize},
  iff/.style={draw=orange!85!black, fill=orange!8, rounded corners=3pt, line width=0.9pt, text width=34mm, align=center, font=\scriptsize},
  lab/.style={font=\scriptsize, fill=white, inner sep=1pt},
  >={Stealth[length=2mm]}]
\node[q] (root) {How is the arena given?};
\node[no, below left=14mm and 12mm of root] (oracle) {\textbf{No when $h=\omega(\log N)$}\\ $2^{\Omega(h)}$ queries to $\pay$, classical or quantum\\ \textcolor{black!60}{\cref{thm:query}}};
\node[q, below=14mm of root] (rules) {Which class of rules?};
\node[yes, below right=14mm and 12mm of root, text width=36mm] (explicit) {\textbf{Yes: linear time}\\ $\Pclass$-complete; in $\NC$ only if $\NC=\Pclass$\\ \textcolor{black!60}{\cref{prop:explicit}}};
\node[yes, below left=10mm and -6mm of rules] (collapse) {\textbf{Yes: polynomial at any horizon}\\ Nim, Wythoff (\S\ref{ssec:phase}); UVG and SCC-symmetric DVG (\S\ref{sec:matching}); one-way arcs with a common tail (\S\ref{sec:oneway}); affine, 2-CNF, Horn QBF (\S\ref{ssec:qcsp}); matroids (\S\ref{ssec:matroid}); Maker--Breaker rank $\le 3$ (\S\ref{ssec:positional}); width along the prefix (\S\ref{ssec:width})};
\node[q, below right=10mm and -2mm of rules] (hor) {Horizon?};
\node[no, below left=9mm and -14mm of hor] (unb) {\textbf{No in deterministic time}\\ $\EXPTIME$-complete; randomized and quantum open\\ \textcolor{black!60}{\cref{cor:exp}}};
\node[iff, below right=9mm and -14mm of hor] (pol) {\textbf{Yes iff $\Pclass=\PSPACE$}\\ $\PSPACE$-complete\\ \textcolor{black!60}{\cref{cor:polyhorizon}}};
\draw[->] (root) -- node[lab, pos=0.45, above left]{$\pay$ by queries only} (oracle);
\draw[->] (root) -- node[lab, right]{succinct circuits} (rules);
\draw[->] (root) -- node[lab, pos=0.45, above right]{explicit graph} (explicit);
\draw[->] (rules) -- node[lab]{collapsing class} (collapse);
\draw[->] (rules) -- node[lab]{all arenas} (hor);
\draw[->] (hor) -- node[lab]{unbounded} (unb);
\draw[->] (hor) -- node[lab]{polynomial} (pol);
\end{tikzpicture}
\end{adjustbox}
\caption{Verdicts on the existence of a polynomial-time invariant $\Phi$ with $\val=f\circ\Phi$, by representation, rule class and horizon. Each ``no'' verdict is for the class of all arenas in the model its box names; the small print names the proof.}
\label{fig:glance}
\end{figure}

\section*{Contributions and status of results}

\Cref{fig:glance} summarizes the complexity verdicts. Sections \ref{sec:setting}--\ref{sec:extensions} restate classical results in the language of arenas. The new material is \cref{prop:blind}, \S\S\ref{sec:spectral}--\ref{sec:matching}, \cref{thm:sigma2} and \S\S\ref{sec:realstab}--\ref{sec:treewidth}. New results are proved here; cited results name their source in the statement. \Cref{app:concordance} maps the numbering of the earlier drafts to this version.

\begin{enumerate}[label=(\arabic*)]
\item \emph{Signed resolvent formula for $\val$} (\cref{cor:index}).
\item \emph{Spectral derivation of the Fraenkel--Scheinerman--Ullman theorem} from Godsil's path-tree theorem (\cref{thm:matching-collapse}).
\item \emph{Directed resolvent collapse} on SCC-symmetric digraphs (\cref{thm:directed-collapse,cor:dvg-poly}), and its \emph{converse}: a local weighted matching potential with real weights exists only on SCC-symmetric digraphs (\cref{thm:conj610}; proof in \S\ref{sec:realstab}).
\item \emph{Real-stable potentials force reciprocity} (\cref{thm:main-real}): the obstruction uses only real stability, so it also covers Hermitian determinantal potentials, nonnegative edge weights, and frozen environments (\cref{thm:frozen-real}); it is sharp for the directed triangle at every $|z|>2$ (\cref{prop:triangle-band}).
\item \emph{The complex form of \cref{thm:conj610} is false} (\cref{thm:complex-counterexample,cor:complex-false}): genuine digraphs realize the complex frozen environments that rescue every directed cycle (\cref{thm:every-cycle}); the obstruction to real potentials is reality, not monodromy (\cref{prop:Ym-arithmetic,rem:monodromy}). What remains open over $\Qzbar$ is \cref{prob:complex}.
\item \emph{One-way arcs}: polynomial time when the one-way arcs of every strong component share a tail (\cref{thm:common-tail}); no key-local rule for any other two-arc geometry (\cref{cor:dichotomy}); at most $p$ crossings with $p$ tails (\cref{lem:p-tails}).
\item \emph{Set potentials}: exact criteria for cycles, books, one blown-up vertex and cacti; a continuant criterion for all blow-ups of directed cycles (\cref{thm:blowup-criterion}), with which every two-periodic blow-up carries a set potential (\cref{lem:twin,cor:two-periodic}).
\item \emph{Treewidth}: the exponential degree of exact resolvents at pathwidth~$2$ (\cref{thm:ladders}) is intrinsic to every exact multiplicative potential, however indexed (\cref{thm:intrinsic-degree}); the resolvent recursion closes on canonical subgames, which are path-dependent and at treewidth~$3$ need $2^{\Omega(\sqrt{n/\log n})}$ states at a single vertex (\cref{prop:memory}); component-state circuits have linear size when strong components are small (\cref{thm:component-circuits}).
\item \emph{Other results}: the residue of the resolvent as a density of matching strategies (\cref{prop:residue}); the torsion obstruction (\cref{thm:torsion}); unconditional limits of landscape invariants (\cref{prop:blind}); $\Sigma^0_2$-completeness of collapsibility (\cref{thm:sigma2}).
\end{enumerate}

\begin{longtable}{@{}L{0.42}L{0.53}@{}}
\toprule
Result & Status\\ \midrule\endhead
\cref{lem:well-defined}, \cref{thm:local}, \cref{cor:certify} & classical: attractors and progress measures\\
\cref{lem:normal-form}, \S\S\ref{sec:barriers}--\ref{sec:extensions}, \cref{thm:strategies}, \cref{cor:chess} & classical, restated for arenas\\
\cref{prop:blind} & new, elementary\\
\cref{lem:threshold}, \cref{thm:resolvent-outcome}, \cref{cor:index} & new assembly of classical facts (Cvetkovi\'c--Gutman; Fraenkel--Scheinerman--Ullman)\\
\cref{lem:transfer}, \cref{thm:matching-collapse} & Godsil 1981 read as a game; the derivation of the FSU theorem appears new\\
\cref{thm:unweighted}, \cref{lem:entry} & new, elementary\\
\cref{thm:directed-collapse}, \cref{prop:no-cancel}, \cref{cor:dvg-poly}, \cref{prop:fails}, \cref{cor:edge-geo} & new as far as we know\\
\cref{thm:conj610} (formerly Conjecture 6.10), \cref{thm:main-real}, \cref{thm:frozen-real}, \cref{prop:triangle-band} & new; proofs in \S\ref{sec:realstab}\\
\S\ref{sec:fields}--\S\ref{sec:complex} (results over arbitrary fields, complex frontier) & new; \cref{thm:complex-counterexample} proved, and independently verified by computer; \cref{thm:exact5}, \cref{prop:env-barrier,prop:smallest-env,prop:c5c6,prop:c3-minimal} and \cref{prop:Ym-arithmetic}(3) are computer-assisted (exact Gr\"obner bases with saturation, PARI)\\
\S\ref{sec:oneway} (one-way arcs) & new; \cref{prop:two-arcs,prop:barrier,prop:ge-barrier} are explicit examples checked by exhaustive search\\
\S\ref{sec:setpot} (set potentials) & new; \cref{thm:cores5,prop:census67} computer-assisted; \cref{thm:blowup-criterion,lem:twin,cor:two-periodic} proved, and checked by computer\\
\S\ref{sec:treewidth} (treewidth) & new; \cref{prop:memory}(2) also checked by computer for $m\le6$\\
\cref{prop:residue} & strategy from FSU 1993; the residue as a density of strategies is new\\
\cref{thm:kasteleyn,thm:hardness}, \cref{prop:rank-game} & consequences of cited hardness results\\
\cref{thm:compositional,thm:torsion} & new, elementary\\
\cref{thm:sigma2} & index-set argument in the style of H\'ajek and Kozen; $\Sigma^0_2$-completeness proved here\\
\bottomrule
\end{longtable}

\section{Problem and setting}\label{sec:setting}

This section states the problem, fixes the model and makes precise what an invariant is. \Cref{prop:blind} then shows that the state graph alone, undirected or through its spectrum, cannot determine $\val$.

\subsection*{The problem as posed}
A finite oriented state graph has the states $\sigma=(S_0,B_0)\in\{0,1\}^N$ as vertices, $2^N$ of them, and two deterministic transition operators $\to^{+}$ and $\to^{-}$ generate its edges. The value of a start $\sigma_0$ is the alternating expression $\val(\sigma_0)=\sup\inf\sup\cdots\pay(\sigma_{\mathrm{terminal}})\in\{-1,0,1\}$ over the successive states. The problem asks for a global algebraic or topological invariant $\Phi\colon\{0,1\}^N\to\{-1,0,1\}$, defined on the connected components or on the global landscape of the state graph, with $\Phi(\sigma_0)=\val(\sigma_0)$, computable in time $O(N^k)$ relative to the input description and without running the alternating search. It suggests three directions: continuous extensions and spectral mapping (\S\S\ref{sec:extensions}--\ref{sec:spectral}), global invariant operators (\S\S\ref{sec:local} and \ref{sec:matching}), and the exact boundary of feasibility (\S\S\ref{sec:algebraic}--\ref{sec:boundary}).

\subsection{Arenas on the hypercube}\label{ssec:arenas}
Positions are the vertices $\sigma$ of the hypercube $\{0,1\}^N$. An \emph{arena} is a tuple $\arena=(N,p,T^{+},T^{-},\tau,\pay)$ where
\begin{itemize}
\item $p\colon\{0,1\}^N\to\{\mathrm{Max},\mathrm{Min}\}$ names the player to move;
\item $T^{+},T^{-}\colon\{0,1\}^N\to\{0,1\}^N$ are the transition operators, so a non-terminal $\sigma$ has the options $T^{+}\sigma$ and $T^{-}\sigma$;
\item $\tau\subseteq\{0,1\}^N$ is the set of terminal positions and $\pay\colon\tau\to\{-1,0,+1\}$ is the payoff from Min to Max.
\end{itemize}
Players alternate: $p(T^{\pm}\sigma)\neq p(\sigma)$ for every non-terminal $\sigma$ (the operators are irrelevant at terminals). The \emph{state graph} $\SG=(\{0,1\}^N,E)$ has an edge from each non-terminal $\sigma$ to $T^{+}\sigma$ and to $T^{-}\sigma$. We do not use the split $\sigma=(S_0,B_0)$ of the problem statement; if $S_0$ and $B_0$ are the sets of elements already claimed by two players, the arena is a positional game, and \S\ref{ssec:positional} gives the sharp threshold for that reading.

\begin{remark*}[conventions]
$\val$ is the value, $\pay$ the payoff, $\SG$ the state graph and $\arena$ an arena. $N$ is the dimension of the cube and $n$ the input size in \S\S\ref{sec:setting}--\ref{sec:extensions} and \S\ref{sec:boundary}; from \S\ref{sec:spectral} on, $n$ is the number of vertices of the graph or digraph at hand and $V=V(D)$ its vertex set. N- and P-positions are outcomes of normal play (\S\ref{sec:spectral}). \Cref{app:notation} lists all symbols.
\end{remark*}

The answer depends on how the arena is presented.

\begin{center}\small
\begin{tabularx}{\linewidth}{@{}l>{\raggedright\arraybackslash}X>{\raggedright\arraybackslash}p{0.27\linewidth}@{}}
\toprule
Representation & What is given & Input size $n$\\ \midrule
Explicit & adjacency lists of $\SG$ & $\Theta(2^N)$\\
Succinct & Boolean circuits for $p$, $T^{+}$, $T^{-}$, $\tau$, $\pay$ & total circuit size $s\ge N$\\
Landscape (oracle) & the move structure; $\pay$ only through queries at terminals & queries counted separately\\
\bottomrule
\end{tabularx}
\end{center}

``Polynomial time relative to the input description size'' is read in the succinct model, where $n=\Theta(s)$. The \emph{horizon} $h(\arena,\sigma_0)$ is the supremum of play lengths from $\sigma_0$, possibly $\infty$; only positions reachable from $\sigma_0$ matter for $\val(\sigma_0)$. A family has \emph{polynomial horizon} if $h(\arena,\sigma_0)\le q(n)$ for a fixed polynomial $q$. This is a promise on the family, not a property one can read off the circuits; \emph{clocked} arenas make it syntactic: a step counter in the position ends play after $q(n)$ moves with payoff $0$.

\subsection{The value}
For finite horizon, $\val$ is defined by backward induction, with $\sup$ at Max positions and $\inf$ at Min positions. The expression in the problem statement is the case where Max moves at $\sigma_0$; when Min moves there it begins with $\inf$:
\[
\val(\sigma_0)=\sup_{\sigma_1}\inf_{\sigma_2}\sup_{\sigma_3}\cdots\,\pay(\sigma_{\mathrm{terminal}}),\qquad \sigma_{k+1}\in\{T^{+}\sigma_k,\,T^{-}\sigma_k\}.
\]
When play can be infinite, an infinite play is a draw worth $0$. Let $W_{+}$ be the set of positions from which Max can force a terminal with $\pay=+1$, and $W_{-}$ the set from which Min can force a terminal with $\pay=-1$. Set $\val=+1$ on $W_{+}$, $\val=-1$ on $W_{-}$ and $\val=0$ elsewhere.

\begin{lemma}[well-defined value]\label{lem:well-defined}
$W_{+}$ and $W_{-}$ are disjoint, $\val$ agrees with backward induction for finite horizon, and both players have positional optimal strategies.
\end{lemma}

\begin{proof}
$W_{+}$ is the least fixed point of $X\mapsto X\cup\{\sigma\notin\tau:\ p(\sigma)=\mathrm{Max}$ and some option lies in $X\}\cup\{\sigma\notin\tau:\ p(\sigma)=\mathrm{Min}$ and both options lie in $X\}$, started from the terminals with $\pay=+1$; $W_{-}$ is defined symmetrically. Reachability games on finite graphs are determined with positional strategies (Zermelo). From outside $W_{+}$ Min has a positional strategy that never reaches a $+1$ terminal, and from outside $W_{-}$ Max has one that never reaches a $-1$ terminal, so both can guarantee $0$ on the rest. If $\sigma$ were in both sets, the two forcing strategies played against each other would end at a terminal with $\pay=+1$ and $\pay=-1$ at once. For finite horizon, induction on the height of $\sigma$ shows that $\val$ is the backward-induction value. The stage at which $\sigma$ enters $W_{+}$ is its attractor rank, the length of a shortest forced win.
\end{proof}

\begin{lemma}[normal form]\label{lem:normal-form}
Let a finite game with at most $b$ options per position, any turn order and payoffs in $\{-1,0,+1\}$ be given by circuits of size $s$ that compute the player to move, the number of options, the $i$-th option of a position and the payoffs. Then an arena with circuits of size $\poly(s,\log b)$ simulates it, and every play grows by a factor $O(1+\log b)$. The format of the problem statement therefore loses no generality.
\end{lemma}

\begin{proof}
Encode a choice among the options by $\max(1,\lceil\log_2 b\rceil)$ binary choices of the same player; codes beyond the last option of a position select that last option. Between two binary choices of one player insert a dummy opponent move with $T^{+}=T^{-}$, and do the same between consecutive moves of one player in the original game. The partial choice is stored in $O(\log b)$ extra bits, and dummy moves change no value.
\end{proof}

One could also read the operators as fixed per player: $T^{+}$ is Max's only move and $T^{-}$ is Min's only move. Then nobody chooses and $\val$ is read off one deterministic orbit. This zero-player reading changes the answers: $\val$ is $\Pclass$-complete for polynomial horizon and $\PSPACE$-complete for unbounded horizon (\cref{thm:alternation}), and in the landscape model one query at the end of the orbit suffices. We use the two-choice reading throughout.

\subsection{What an invariant is}
\begin{definition*}
An \emph{invariant} for a family of arenas is a map $\Phi$ from instances (an arena with a start position) to a set $\Omega$, together with a decoder $f\colon\Omega\to\{-1,0,+1\}$ such that $\val=f\circ\Phi$. It is \emph{polynomial-time} if $\Phi$ and $f$ are computable in time polynomial in $n$. The problem's $\Phi$ is the case $\Omega=\{-1,0,+1\}$ with $f$ the identity; \S\S\ref{sec:spectral}--\ref{sec:matching} use invariants valued in rational functions and divisor classes. An invariant is \emph{defined on} a structure $S(\SG)$ of the state graph if $\Phi(\arena,\sigma_0)$ depends only on $S(\SG)$ and $\sigma_0$.
\end{definition*}

\begin{proposition}[landscape invariants are blind]\label{prop:blind}
\begin{enumerate}[label=\textup{(\arabic*)}]
\item $\val$ is not a function of the labelled undirected state graph, so no invariant is defined on it, whatever its cost. Here the graph underlying $\SG$ carries the labels $p$, the terminal set $\tau$ and the payoffs $\pay$, with $\sigma_0$ marked; connected components, spectra and homology are functions of it.
\item At finite horizon the part of $\SG$ reachable from $\sigma_0$ is acyclic, so its adjacency matrix is nilpotent and its strong components are single positions. Neither the spectrum nor the strong components of $\SG$ determine $\val(\sigma_0)$.
\end{enumerate}
\end{proposition}

\begin{proof}
(1) Take $N=2$ and positions $u,v,s,t$ with $p(u)=p(s)=\mathrm{Max}$ and $p(v)=p(t)=\mathrm{Min}$, where $s,t$ are terminals with $\pay(s)=+1$ and $\pay(t)=-1$. In $\arena_1$, $u$ has the options $t$ and $v$, and $v$ has the option $s$ twice. In $\arena_2$, $u$ has the option $t$ twice, and $v$ has the options $s$ and $u$. Players alternate in both, and both have the undirected edges $ut,uv,vs$ with the same labels ($E$ is a set, so a repeated option gives one edge). In $\arena_1$, $\val(v)=+1$ and $\val(u)=\max(-1,+1)=+1$; in $\arena_2$, $\val(u)=-1$ and $\val(v)=\min(+1,-1)=-1$. Mark $\sigma_0=u$.
(2) A cycle reachable from $\sigma_0$ would give an infinite play, so the reachable part is acyclic and its adjacency matrix is nilpotent. A chain of $m$ positions ending at a terminal with payoff $c$ has spectrum $\{0\}$ and single-position components for each $c\in\{-1,0,1\}$, and positions that $\sigma_0$ cannot reach can be fixed identically in the three arenas.
\end{proof}

\Cref{prop:blind} is why \S\ref{sec:spectral} unfolds $\SG$ into its game tree. The unfolding turns orientation into the parent--child relation, and the root resolvents of two threshold trees then determine the value (\cref{cor:index}), although the spectrum of one tree alone does not. \Cref{prop:components} adds that for succinct $\SG$ even the connected components are $\PSPACE$-hard to compute.

\begin{remark}[what ``bypassing the search tree'' can mean]\label{rem:bypass}
Avoiding the $\sup/\inf$ iteration is not a formal constraint: Bouton's rule for Nim already avoids it. The formal requirement is the time bound alone. The question studied here is therefore: for which families, in which representation, is $\val$ computable in time polynomial in $n$?
\end{remark}

\begin{remark}[knowledge at one point does not localize]\label{rem:strategy-stealing}
Strategy stealing shows that the first player wins every minimum poset game, played on a poset with a least element where a move removes the chosen element and everything below it. It also shows that the first player never loses a symmetric Maker--Maker game. Yet finding such a first move is $\PSPACE$-complete in both classes \cite{BodwinGrossman}. A useful invariant must be known across the landscape, which \S\ref{sec:local} makes precise.
\end{remark}

\section{Local characterization: global invariants are locally consistent labelings}\label{sec:local}

A labeling $\Phi$ equals $\val$ exactly when it matches $\pay$ at terminals, obeys the one-step minimax rule everywhere and, when play can cycle, carries a ranking that certifies progress. Checking a candidate is local and cheap; finding one is the whole difficulty.

\begin{theorem}\label{thm:local}
Let $\Phi\colon\{0,1\}^N\to\{-1,0,+1\}$. Then $\Phi=\val$ if and only if:
\begin{enumerate}[label=\textup{(\arabic*)}]
\item \textup{(boundary)} $\Phi=\pay$ on $\tau$;
\item \textup{(Bellman consistency)} for non-terminal $\sigma$, $\Phi(\sigma)$ is the maximum of $\Phi(T^{+}\sigma)$ and $\Phi(T^{-}\sigma)$ if Max moves at $\sigma$, and their minimum if Min moves;
\item \textup{(ranking)} there are functions $\rk_{+}\colon\Phi^{-1}(+1)\to\NN$ and $\rk_{-}\colon\Phi^{-1}(-1)\to\NN$ such that at every non-terminal $\sigma$ with $\Phi(\sigma)=+1$: if Max moves, some option $\sigma'$ has $\Phi(\sigma')=+1$ and $\rk_{+}(\sigma')<\rk_{+}(\sigma)$; if Min moves, every option $\sigma'$ \textup{(}which has $\Phi(\sigma')=+1$ by \textup{(2)}\textup{)} satisfies $\rk_{+}(\sigma')<\rk_{+}(\sigma)$. The mirror condition holds for $\rk_{-}$ with Max and Min exchanged.
\end{enumerate}
For finite horizon condition \textup{(3)} holds automatically, with $\rk$ the remaining depth. Rankings of this kind are the progress measures of Jurdzi\'nski \cite{Jurdzinski}.
\end{theorem}

\begin{proof}
($\Leftarrow$) Let $\Phi(\sigma)=+1$. At Max positions in $\Phi^{-1}(+1)$ Max picks an option that keeps $\Phi=+1$ and lowers $\rk_{+}$; at Min positions every option does both, by (2) and (3). Along any play consistent with this strategy $\rk_{+}$ strictly decreases, so the play is finite and ends at a terminal with $\Phi=+1$, where $\pay=+1$ by (1). Hence $\sigma\in W_{+}$.
Now let $\Phi(\sigma)\le 0$. At Min positions with $\Phi\le0$, (2) gives an option with $\Phi\le 0$, and Min takes it; at Max positions with $\Phi\le0$, (2) gives $\Phi\le 0$ at both options. So play stays where $\Phi\le 0$, any terminal it reaches has $\pay\le 0$, and $\sigma\notin W_{+}$. Hence $\Phi^{-1}(+1)=W_{+}$; the mirror argument gives $\Phi^{-1}(-1)=W_{-}$, and therefore $\Phi^{-1}(0)=\val^{-1}(0)$.
($\Rightarrow$) $\val$ satisfies (1) and (2) by \cref{lem:well-defined}, and the stages at which positions enter the fixed points $W_{+}$ and $W_{-}$ serve as $\rk_{+}$ and $\rk_{-}$.
\end{proof}

Consistency alone is not enough when play can cycle: label $+1$ both positions of a $2$-cycle with no exit; the labeling is consistent, yet nobody can force termination and the true value is $0$. The ranking rules this out.

\begin{corollary}[certification versus discovery]\label{cor:certify}
In a succinct family, suppose a candidate $\Phi$ and rankings $\rk_{\pm}$ with values of polynomial bit length are computable in polynomial time. Then ``$\Phi=\val$ on the whole landscape'' is a $\coNP$ statement on each instance: a counterexample is one position that violates a local condition. The difficulty lies entirely in finding such a $\Phi$.
\end{corollary}

\begin{remark}[the template of known collapses]\label{rem:template}
Bouton's rule for Nim, the Fraenkel--Scheinerman--Ullman theorem and the Aspvall--Plass--Tarjan theorem can each be phrased this way: each exhibits a polynomial-time $\Phi$ and verifies conditions (1)--(2) by a local argument. The Sprague--Grundy theorem has the same local form, but its $\Phi$, the Grundy value, is polynomial-time only on special families; for undirected vertex geography it is $\PSPACE$-complete (\S\ref{ssec:fails}). The phrase of the problem statement, ``a global invariant dictates local terminal states'', is \cref{thm:local} read from right to left.
\end{remark}

\section{Background: the classical barriers}\label{sec:barriers}

Three classical barriers rule out a polynomial-time $\Phi$ for the class of all arenas, each in a stated model. The query barrier binds every algorithm, classical or quantum, that sees $\pay$ only through queries. The unbounded-horizon barrier is unconditional for deterministic time; its randomized and quantum versions are open. The polynomial-horizon barrier is equivalent to $\Pclass\neq\PSPACE$.

\subsection{The landscape barrier}

\begin{theorem}[query lower bounds]\label{thm:query}
Let the move structure be the complete binary tree of depth $h$, with Max and Min alternating, and let $\pay$ be available only through an oracle queried at the $2^h$ leaves. In the worst case, computing $\val$ at the root takes the following numbers of queries.
\begin{center}\small
\begin{tabularx}{\linewidth}{@{}l>{\raggedright\arraybackslash}X>{\raggedright\arraybackslash}X@{}}
\toprule
Model & Queries & Source\\ \midrule
Deterministic & $2^h$ (every leaf) & proof below\\
Randomized & $\Theta(2^{\alpha h})$, with $\alpha=\log_2\frac{1+\sqrt{33}}{4}\approx 0.7537$ & \cite{Snir,SaksWigderson} (zero error); \cite{Santha} (bounded error)\\
Quantum & $\Theta(2^{h/2})$ & lower bound \cite{BarnumSaks}; $O(2^{h/2})$ for balanced formulas \cite{ACRSZ,ReichardtSpalek}\\
\bottomrule
\end{tabularx}
\end{center}
\end{theorem}

\noindent\emph{Input model.} The move structure is given by circuits of size $\poly(N)$ and $\pay$ only by the oracle. Naming the $2^{h+1}-1$ nodes of the tree needs $N\ge h+1$ bits, so the bounds are super-polynomial in $N$ exactly when $h=\omega(\log N)$; at $h=N-1$ even the quantum bound is $2^{(N-1)/2}$. If the tree were given explicitly, $2^h$ queries would be linear in the input and no barrier would remain. For three-valued leaves the upper bounds double (two Boolean thresholds, $\val\ge 0$ and $\val\ge1$) and the lower bounds apply verbatim to leaves in $\{-1,+1\}$.

\begin{proof}[Proof of the deterministic bound]
By induction on $h$, an adversary can force every leaf to be queried and can choose the root's value in $\{-1,+1\}$ at the last query. For $h=0$ this is immediate. Let the root be a Max node with subtrees $L$ and $R$ (the Min case is symmetric) and run the inductive adversary separately inside $L$ and $R$. When one subtree has its last leaf queried, give it the value $-1$, which cannot decide a maximum. The root now equals the other subtree's value, which still needs all of its leaves and which the adversary chooses at the very last query.
\end{proof}

The Hamiltonian-oracle algorithm of Farhi, Goldstone and Gutmann \cite{FGG} probes exactly the zero-energy eigenvector identified in \S\ref{sec:spectral}; discrete-query algorithms reach $2^{(1/2+o(1))h}$ for arbitrary formulas \cite{CCJY,ACRSZ} and $O(2^{h/2})$ for balanced ones \cite{ACRSZ}, and the general adversary bound is tight \cite{Reichardt}. In the query model, spectral or quantum evaluation of minimax therefore gains a square root and no more. Succinct arenas pose a different question: a polynomial-time quantum algorithm for $\val$ at polynomial horizon would place $\PSPACE$ in $\BQP$, which is open.

\subsection{Succinct arenas: alternation and time}

\begin{theorem}[Chandra--Kozen--Stockmeyer; see Hearn--Demaine]\label{thm:alternation}
For circuit-described arenas, the complexity of computing $\val(\sigma_0)$ is:
\begin{center}\small
\begin{tabularx}{\linewidth}{@{}>{\raggedright\arraybackslash}Xll@{}}
\toprule
Players with real choices & Polynomial horizon & Unbounded horizon\\ \midrule
None ($T^{+}=T^{-}$ everywhere) & $\Pclass$-complete & $\PSPACE$-complete\\
One (the other makes only dummy moves) & $\NP$- or $\coNP$-complete & $\PSPACE$-complete\\
Two & $\PSPACE$-complete & $\EXPTIME$-complete\\
\bottomrule
\end{tabularx}
\end{center}
\end{theorem}

\begin{proof}[Proof sketch]
Upper bounds: depth-first evaluation of the game tree needs polynomial space when the horizon is polynomial, and computing the fixed points $W_{\pm}$ over all $2^N$ positions takes time $2^{O(N)}$. Lower bounds: configurations of a Turing machine with space bound $s(n)$ are $s(n)$-bit strings, and its one-step relation has circuits of size $\poly(s(n))$. Existential states become Max positions, universal states Min positions, and accepting or rejecting halts become terminals with $\pay=\pm1$; \cref{lem:normal-form} normalizes the rest. $\mathsf{APTIME}=\PSPACE$ and $\mathsf{APSPACE}=\EXPTIME$ \cite{CKS} give the two-player row; deterministic and nondeterministic machines give the other rows \cite{HearnDemaine}. Natural $\EXPTIME$-complete games include generalized chess \cite{FraenkelLichtenstein} and checkers \cite{Robson}, following \cite{StockmeyerChandra}.
\end{proof}

Two conventions matter in the lower bounds. An alternating machine rejects on an infinite branch, while an infinite play of an arena is a draw worth $0$; the reduction is exact for the threshold question $\val\ge1$, and a step counter that makes every branch halt gives $\val\in\{-1,+1\}$. The polynomial-horizon column uses the horizon promise of \S\ref{ssec:arenas}; clocked arenas give the same bounds.

\begin{corollary}[unbounded play, deterministic time]\label{cor:exp}
There is $\varepsilon>0$ such that no deterministic algorithm computes $\val$ in time $O(2^{n^{\varepsilon}})$ on succinct arenas of unbounded horizon. In particular no deterministic polynomial-time $\Phi$ exists. Randomized and quantum versions would follow from $\EXP\not\subseteq\BPP$ and $\EXP\not\subseteq\BQP$, which are open.
\end{corollary}

\begin{proof}
By the time hierarchy theorem \cite{HartmanisStearns} some language $L$ lies in $\mathsf{DTIME}(2^{n^2})$ but not in $\mathsf{DTIME}(2^{2n})$. \Cref{thm:alternation} reduces $L$ to $\val$ in polynomial time with outputs of size at most $n^c$. An algorithm for $\val$ running in time $2^{m^{1/c}}$ on inputs of size $m$ would decide $L$ in time $\poly(n)+2^{n}$, a contradiction. Take $\varepsilon=1/c$.
\end{proof}

\begin{corollary}[polynomial horizon]\label{cor:polyhorizon}
A polynomial-time $\Phi$ exists for succinct arenas of polynomial horizon if and only if $\Pclass=\PSPACE$.
\end{corollary}

\begin{proposition}[explicit graphs]\label{prop:explicit}
For explicit $\SG$, $\val$ is computable in time linear in the size of $\SG$ by retrograde analysis, and the problem is $\Pclass$-complete: it is alternating graph accessibility \cite{Immerman,GHR}. So no invariant computable in polylogarithmic parallel time exists unless $\NC=\Pclass$.
\end{proposition}

\subsection{Connected components of the state graph}

\begin{proposition}\label{prop:components}
For succinct $\SG$, deciding whether $\sigma_0$ and a given position $\sigma^{*}$ lie in the same connected component of the undirected graph underlying $\SG$ is $\PSPACE$-complete. This is an instance of the succinct-graph phenomenon of \cite{GalperinWigderson,PapadimitriouYannakakis,LozanoBalcazar}; we give a proof inside the arena format.
\end{proposition}

\begin{proof}
Membership: undirected reachability among $2^N$ positions is in $\mathsf{DSPACE}(N)$ \cite{Reingold} (already $\mathsf{NSPACE}(N)\subseteq\mathsf{DSPACE}(N^2)$ suffices \cite{Savitch}). Hardness: let a deterministic machine $\mathcal D$ decide $L\in\PSPACE$ in space $s(n)$, erasing its tape before it halts, so that it has a unique accepting configuration $c_{+}$ and a unique rejecting configuration $c_{-}$. Positions are pairs $(c,b)$ with a parity bit $b$; Max moves when $b=0$ and Min when $b=1$. A position whose $c$ is a non-halting configuration has $T^{+}=T^{-}=(\mathrm{next}(c),1-b)$. Add one move $(c_{+},1)\to(c_{+},0)$, and make $(c_{+},0)$, $(c_{-},0)$, $(c_{-},1)$ and every string that encodes no configuration terminal. The circuits have size $\poly(s(n))$. Every non-terminal position has exactly one out-neighbour and terminals have none, so a component with $k$ positions has at most $k$ edges and, being connected with at least $k-1$ edges, at most one terminal; a component containing a terminal is a tree directed towards it. Hence the component of $(c_{+},0)$ is the set of positions whose forward orbit reaches $(c_{+},0)$, and $(c_{\mathrm{start}},0)$ lies in it exactly when $\mathcal D$ accepts.
\end{proof}

So the components of $\SG$ are $\PSPACE$-hard to compute from a succinct description. Hardness is not the main obstruction, though: by \cref{prop:blind}, components and every other invariant of the undirected graph fail to determine $\val$ even when they are computed for free.

\subsection{Static sparsity is not enough}
A QBF is an arena whose positions are partial assignments. Atserias and Oliva \cite{AtseriasOliva} proved that QBF stays $\PSPACE$-complete on formulas of constant pathwidth, hence of constant treewidth. A sparse interaction graph does not collapse the minimax tree; \S\ref{ssec:width} shows what does.

\section{Algebraic and continuous extensions (background)}\label{sec:extensions}

Embedding $\{0,1\}^N$ into $\FF^N$ or $\RR^N$ yields canonical extensions of $\val$, and the classical results below show what each costs. Low-degree extensions over finite fields are as hard on average as $\val$ itself (\cref{thm:worst-average}) but make $\val$ verifiable (\cref{thm:ip}). Unit-cost arithmetic on exponentially long words evaluates $\val$ quickly (\cref{thm:unit-cost}), and a polynomial-time real-number algorithm with algebraic constants would give $\PSPACE=\CH$ (\cref{thm:bss}).

\subsection{Canonical extensions}
Fix a polynomial-horizon arena in normal form, with terminals made absorbing by dummy moves so that every play has exactly $h$ moves. A play from $\sigma_0$ is then a move string $x\in\{0,1\}^h$. For a threshold $t\in\{0,1\}$ let $W_t(x)=1$ if the play $x$ ends at a terminal with $\pay\ge t$ and $0$ otherwise. Composing the transition circuits $h$ times gives a circuit for $W_t$ of size $\poly(n)$. Arithmetize: replace a NOT gate $a$ by $1-a$ and an AND gate $(a,b)$ by $ab$; replace the move quantifiers (OR at Max moves, AND at Min moves) by
\[
\textstyle\bigvee_{x_i}f\ \mapsto\ 1-\bigl(1-f|_{x_i=0}\bigr)\bigl(1-f|_{x_i=1}\bigr),\qquad \bigwedge_{x_i}f\ \mapsto\ f|_{x_i=0}\cdot f|_{x_i=1}.
\]
The result is an integer expression $P_t$ equal to $\one[\val\ge t]$, and $\val(\sigma_0)=P_1+P_0-1$. This is a closed-form algebraic map that compresses the history of play into one static expression; its evaluation is $\PSPACE$-complete by \cref{thm:alternation}. A second canonical extension is the multilinear extension of the whole value function $f\colon\{0,1\}^n\to\{-1,0,+1\}$, with the arena description among its input bits:
\[
\widehat f(y)=\sum_{a\in\{0,1\}^n}f(a)\prod_{i=1}^{n}\bigl(a_iy_i+(1-a_i)(1-y_i)\bigr),\qquad y\in\FF^n.
\]

\subsection{Low-degree extensions over finite fields are as hard as the cube}

\begin{theorem}[Beaver--Feigenbaum; Lipton]\label{thm:worst-average}
Let $\FF$ be a finite field of characteristic greater than $2$ with at least $n+2$ elements. Suppose a randomized polynomial-time algorithm $\mathsf A$ satisfies $\Pr[\mathsf A(y)=\widehat f(y)]\ge 1-1/(3(n+1))$ for uniformly random $y\in\FF^n$. Then $f\in\BPP$. For polynomial horizon this gives $\PSPACE\subseteq\BPP$, and for unbounded horizon $\EXP=\BPP$.
\end{theorem}

\begin{proof}[Proof \textup{\cite{BeaverFeigenbaum,Lipton}}]
To compute $f(a)$ for $a\in\{0,1\}^n$, choose $b\in\FF^n$ uniformly and restrict $\widehat f$ to the line $s\mapsto a+sb$. The restriction $g$ is a univariate polynomial of degree at most $n$ with $g(0)=f(a)$, and for each fixed $s\neq0$ the point $a+sb$ is uniform in $\FF^n$. Query $\mathsf A$ at $s=1,\dots,n+1$; by the union bound all answers are correct with probability at least $2/3$. Interpolate $g$ and output $g(0)$. Characteristic $\neq2$ keeps $-1\neq+1$.
\end{proof}

List decoding lowers the required agreement much further for suitably encoded $\PSPACE$- and $\EXP$-complete functions \cite{STV,TrevisanVadhan}. So no low-degree extension can be easy on most of its domain while encoding $\val$ at the vertices: hardness is a property of the whole extension, not a boundary effect at the cube.

\subsection{What algebraic extensions buy: verification}
\begin{theorem}[Lund--Fortnow--Karloff--Nisan; Shamir]\label{thm:ip}
$\IP=\PSPACE$ \cite{LFKN,Shamir}. The expression $P_t$ admits an interactive proof with a polynomial-time verifier, via the sum-check protocol with degree reduction.
\end{theorem}
A claimed value of a polynomial-horizon game can therefore be checked in polynomial time against an all-powerful prover: the algebraic extension buys verification, not evaluation.

\begin{corollary}[non-uniform invariants]\label{cor:nonuniform}
If $\val$ had circuits $\Phi_n$ of size $\poly(n)$ for every input length $n$, then $\PSPACE=\MA$ for polynomial horizon and $\EXP=\MA$ for unbounded horizon \cite{BFNW}.
\end{corollary}
\begin{proof}[Proof sketch for $\PSPACE$]
The honest prover of the $\IP=\PSPACE$ protocol runs in polynomial space, so under the hypothesis it has polynomial-size circuits. Merlin sends such a circuit and Arthur runs the protocol against it; soundness holds against every prover, so false claims are still rejected.
\end{proof}

\subsection{The unit-cost mirage}
\begin{theorem}\label{thm:unit-cost}
On a random-access machine with unit-cost addition, subtraction, multiplication and bitwise Boolean operations on unbounded integers, $\val$ is computable in $\poly(n)$ steps for polynomial horizon. Hartmanis and Simon \cite{HartmanisSimon} showed that such machines decide exactly $\PSPACE$ in polynomial time.
\end{theorem}

\begin{proof}[Construction \textup{(word-parallel minimax)}]
Index move strings by $j\in\{0,\dots,2^h-1\}$, with $x_1$ the most significant bit. Represent a Boolean function of $x$ by the $2^h$-bit integer whose bit $j$ is its value at $j$. A left shift by $2^m$ positions is multiplication by $2^{2^m}$, obtained by $m$ squarings.
\begin{enumerate}[label=\arabic*.]
\item Build each projection word $X_i$, whose bit $j$ is the move bit $x_i$ of $j$. It has period $2^{h-i+1}$: take one block $(2^{2^{h-i}}-1)\cdot2^{2^{h-i}}$ and replicate it by $O(h)$ doublings $Y\mapsto Y\ \mathrm{OR}\ (Y\ll \text{a multiple of the period})$.
\item Evaluate the circuit for $W_t$ gate by gate with bitwise operations on these words (NOT is subtraction from the all-ones word). The result $L_t$ has bit $j$ equal to $W_t(j)$.
\item Fold from the last move upward: for $i=h,\dots,1$ replace $L$ by $(L\ \mathrm{op}\ (L\ll 2^{h-i}))\ \mathrm{AND}\ Z_i$, where op is OR when Max chooses $x_i$ and AND when Min does, and $Z_i$ keeps the positions whose move bits $x_i,\dots,x_h$ all equal $1$.
\item Bit $2^h-1$ of the final word is $\one[\val\ge t]$; it is read by an AND with $2^{2^h-1}=\prod_{k<h}2^{2^k}$.
\end{enumerate}
Each step costs $O(h)$ or $O(\text{circuit size})$ operations, so the total is polynomial.
\end{proof}

\noindent\emph{Accounting.} Under the logarithmic cost measure every word has $2^h$ bits and the same computation costs $2^{\Theta(h)}$: the apparent speed-up is exactly the power to operate on exponentially long words at unit cost.

\subsection{Real-number machines}
\begin{theorem}\label{thm:bss}
Suppose a Blum--Shub--Smale machine over $\RR$ with algebraic constants computes $\val$ in polynomial time on $0/1$ inputs. Then $\PSPACE=\CH$ for polynomial horizon, and $\EXP=\PSPACE=\CH$ for unbounded horizon, where $\CH$ is the counting hierarchy.
\end{theorem}

\begin{proof}
Algebraic constants do not add power: constant-free polynomial time over $\RR$ equals polynomial time with real algebraic constants \cite{ChapuisKoiran}. The Boolean part of constant-free polynomial time is $\Pclass^{\mathrm{PosSLP}}$, and $\mathrm{PosSLP}\in\CH$ \cite{ABKM}. So each threshold $\one[\val\ge t]$ would lie in $\Pclass^{\CH}=\CH$. For polynomial horizon $\PSPACE\subseteq\CH\subseteq\PSPACE$; for unbounded horizon $\EXP\subseteq\CH\subseteq\PSPACE\subseteq\EXP$.
\end{proof}

Neither collapse is known to be false, and both are widely considered implausible. Arbitrary real constants must be excluded, because one real number can encode unbounded advice.

\noindent\emph{Summary.} Each extension moves the cost to a stated resource: average-case hardness over finite fields, word length under unit-cost arithmetic, and the conjecture $\PSPACE\neq\CH$ over the reals, which is stronger than $\Pclass\neq\PSPACE$. None gives a polynomial-time evaluation in the bit model.

\section{Spectral mapping: an exact resolvent formula for the value}\label{sec:spectral}

Up to a sign fixed by who moves first, $\val$ is half the sum of the orders of vanishing at $z=0$ of the root resolvents of two threshold trees built from the unfolded game tree. This reformulates $\val$ exactly; it does not by itself make $\val$ cheaper to compute.

\subsection{From three values to normal play}
In the normal-play game on a finite rooted tree, the player to move takes the token from the current node to a child, and a player who cannot move loses. A node is an \emph{N-position} if the player to move wins and a \emph{P-position} otherwise; a node is N exactly when some child is P.

\begin{lemma}\label{lem:threshold}
Take a finite-horizon arena, a start $\sigma_0$ and $t\in\{0,1\}$. Let $\mathcal T_t$ be the game tree of $\sigma_0$ (nodes are histories of play), and at each leaf $\ell$ attach one pendant child exactly when $p(\ell)=\mathrm{Max}$ and $\pay(\ell)\ge t$, or $p(\ell)=\mathrm{Min}$ and $\pay(\ell)<t$. Then the root of $\mathcal T_t$ is an N-position if and only if the player to move at $\sigma_0$ wins the $t$-game: $\val(\sigma_0)\ge t$ when Max moves at $\sigma_0$, and $\val(\sigma_0)<t$ when Min moves there. $\mathcal T_t$ has at most twice as many nodes as the game tree.
\end{lemma}

\begin{proof}
Say the player to move at $v$ wins the $t$-game from $v$ if Max can force $\pay\ge t$ when Max moves at $v$, or Min can force $\pay<t$ when Min moves there. By induction on height, $v$ is N in $\mathcal T_t$ exactly when the player to move wins the $t$-game from $v$. A leaf with a pendant child can move to a node with no moves, so it is N; a leaf without one is P; this matches the attaching rule. At an internal node the player to move wins exactly when some child is a position from which the opponent, now to move (players alternate), loses; by induction that child is P, so $v$ is N.
\end{proof}

\subsection{The resolvent outcome theorem}
For a finite tree $T$ with adjacency matrix $A$ and a node $r$, the \emph{root resolvent} is
\[
R_r(z)=\bigl\langle e_r,(zI-A)^{-1}e_r\bigr\rangle=\frac{\varphi(T-r,z)}{\varphi(T,z)}=\sum_j\frac{|\langle e_r,v_j\rangle|^2}{z-\lambda_j},
\]
where $\varphi$ is the characteristic polynomial and $(\lambda_j,v_j)$ an orthonormal eigenbasis of $A$. Write $\ord_0$ for the order of vanishing at $z=0$, negative for a pole. Where a function has no pole at $0$, the spectral sum gives its derivative there as $-\sum_j|\langle e_r,v_j\rangle|^2/\lambda_j^2<0$; $-R_r$ is a Herglotz function.

\begin{theorem}[resolvent outcome theorem]\label{thm:resolvent-outcome}
Let $T$ be a finite tree rooted at $r$. Then $\ord_0R_r\in\{-1,+1\}$, and the following are equivalent:
\begin{enumerate}[label=\textup{(\arabic*)}]
\item $r$ is an N-position of normal play on $T$;
\item $\ord_0R_r=+1$, that is, $R_r$ has a simple zero at $0$;
\item $e_r\in\operatorname{Im}A$, that is, every vector in $\ker A$ vanishes at $r$;
\item $r$ is covered by every maximum matching of $T$.
\end{enumerate}
Correspondingly $r$ is a P-position exactly when $R_r$ has a simple pole at $0$, and the residue is the squared length of the projection of $e_r$ onto $\ker A$.
\end{theorem}

\begin{proof}
Let $r$ have children $c_1,\dots,c_k$ with subtrees $T_1,\dots,T_k$. Removing $r$ splits $T$ into the $T_i$, and the Schur complement of the $(r,r)$ entry gives
\begin{equation}\label{eq:schur}
R_r(z)=\Bigl(z-\sum_{i=1}^kR^{T_i}_{c_i}(z)\Bigr)^{-1},\qquad R_r(z)=\frac1z\ \text{for a leaf.}
\end{equation}
We prove by induction on height that either $R_r$ has a simple pole at $0$ with positive residue and $r$ is P, or $R_r(0)=0$ with negative derivative and $r$ is N. A leaf has $R_r=1/z$ and is P. If some child is P, its term has a simple pole with positive residue and every other term is such a pole or regular; positive residues cannot cancel, so the sum has a simple pole with residue $c>0$, and $R_r$ has a simple zero with derivative $-1/c$: $r$ is N. If every child is N, each term vanishes at $0$ with negative derivative, so the denominator is $z(1-\sum_iR_i'(0))+O(z^2)$ with $1-\sum_iR_i'(0)>1$, and $R_r$ has a simple pole with positive residue: $r$ is P. This proves (1)$\Leftrightarrow$(2) and $\ord_0R_r\in\{\pm1\}$.
For (2)$\Leftrightarrow$(3), the spectral sum has a pole at $0$ exactly when $e_r$ has a nonzero projection onto $\ker A$, that is, when $e_r\notin(\ker A)^{\perp}=\operatorname{Im}A$. For (3)$\Leftrightarrow$(4): in a forest $F$ a permutation contributes to $\det A(F[S])$ only if all its cycles are edges, so $\det A(F[S])$ is $\pm$ the number of perfect matchings of $F[S]$, which is $0$ or $1$. The rank of a symmetric matrix is the largest size of a nonsingular principal submatrix, so $\rk A(F)=2\nu(F)$, where $\nu$ is the matching number. Therefore $\ord_0R_r=\operatorname{nullity}A(T-r)-\operatorname{nullity}A(T)=2(\nu(T)-\nu(T-r))-1$, which is $+1$ exactly when every maximum matching covers $r$.
\end{proof}

On trees, (1)$\Leftrightarrow$(4) is the tree case of the Fraenkel--Scheinerman--Ullman theorem; the identity $\operatorname{nullity}A(T)=|T|-2\nu(T)$ is due to Cvetkovi\'c and Gutman \cite{CvetkovicGutman}. In Parter--Wiener terminology \cite{JLS}, $\ord_0R_r=+1$ says that $r$ is a Parter vertex for the eigenvalue $0$ and $\ord_0R_r=-1$ that it is a downer vertex; no vertex of a tree is neutral for $0$. The equivalence (1)$\Leftrightarrow$(3) is the exact classical counterpart of the zero-energy analysis of the NAND-tree algorithm \cite{FGG} and of span programs \cite{ReichardtSpalek}.

\subsection{The spectral index formula}
\begin{corollary}\label{cor:index}
For every finite-horizon arena and start $\sigma_0$, let $R^{(t)}$ be the root resolvent of $\mathcal T_t$ from \cref{lem:threshold}, and let $s(\sigma_0)=+1$ if Max moves at $\sigma_0$ and $s(\sigma_0)=-1$ if Min does. Then
\[
\val(\sigma_0)=\tfrac{s(\sigma_0)}2\Bigl(\ord_{z=0}R^{(1)}(z)+\ord_{z=0}R^{(0)}(z)\Bigr).
\]
\end{corollary}

\begin{proof}
By \cref{lem:threshold,thm:resolvent-outcome}, the root of $\mathcal T_t$ is N exactly when $\val\ge t$ if Max moves at $\sigma_0$, and exactly when $\val<t$ if Min does. In both cases $\one[\val\ge t]=(1+s(\sigma_0)\ord_0R^{(t)})/2$. Since $\val\in\{-1,0,1\}$, $\val=\one[\val\ge1]+\one[\val\ge0]-1$.
\end{proof}

Equivalently $\val(\sigma_0)=s(\sigma_0)\bigl(\one[e_r\in\operatorname{Im}A_1]+\one[e_r\in\operatorname{Im}A_0]-1\bigr)$ with $A_t$ the adjacency matrix of $\mathcal T_t$. For example, if Min moves at $\sigma_0$ and both options are terminals with $\pay=+1$, both roots are P-positions, each order is $-1$, and the formula gives $(-\frac12)(-1-1)=+1=\val(\sigma_0)$.

\begin{remark}\label{rem:index}
\Cref{cor:index} is an exact basis-free reformulation of $\val$: the order of a zero or pole of a resolvent entry at one point, or whether one vector lies in the image of a self-adjoint operator; equivalently, the multiplicity of the root $z=0$ of $\det(zI-A_t)$ before and after one vertex is deleted. It does not bypass the search tree: the proof of \cref{thm:resolvent-outcome} evaluates the order by the recursion \eqref{eq:schur}, which is backward induction on a tree with up to $2^h$ nodes, and \cref{thm:alternation} rules out a general polynomial-time evaluation. \S\ref{sec:matching} finds families where the tree collapses to a polynomial-size object.
\end{remark}

\section{Spectral collapse: the matching potential and where it breaks}\label{sec:matching}

For undirected vertex geography the exponential resolvent of \S\ref{sec:spectral} equals a ratio of two matching polynomials of the input graph, and its order at zero is computable by maximum matching. This gives a spectral proof of the Fraenkel--Scheinerman--Ullman theorem and isolates the structure responsible for the collapse; \S\ref{ssec:directed} extends it to directed geography and characterizes exactly where it survives.

\subsection{A transfer lemma}
\begin{lemma}[transfer]\label{lem:transfer}
Consider a family of finite impartial normal-play games. Suppose each position $q$ carries two nonzero rational functions $A(q),B(q)\in\Qz$ such that, for every position $q$ with options $q_1,\dots,q_k$,
\[
\frac{A(q)}{B(q)}=z-\sum_{i=1}^k\frac{B(q_i)}{A(q_i)}.
\]
Then $B(q)/A(q)$ is the root resolvent of the game tree at $q$. Consequently $q$ is an N-position if and only if $\ord_0B(q)-\ord_0A(q)=+1$.
\end{lemma}
\begin{proof}
Both families satisfy the recursion \eqref{eq:schur} with the same value $1/z$ at terminal positions ($k=0$). They agree by induction on height, and \cref{thm:resolvent-outcome} identifies the outcome.
\end{proof}

\begin{remark*}[the lemma is not a structure theorem]
Without restrictions on $A$ and $B$ every family satisfies \cref{lem:transfer}: take $A(q)$ and $B(q)$ to be the characteristic polynomials of the game tree at $q$ and of that tree minus its root. Then $\ord_0B-\ord_0A$ is computable in polynomial time exactly when outcomes are. Content comes only from a restriction on the form of $A$ and $B$: matching polynomials of the input graph in \S\ref{ssec:uvg}, and local weighted matching potentials in \S\ref{ssec:directed}.
\end{remark*}

\subsection{Undirected vertex geography}\label{ssec:uvg}
In undirected vertex geography (UVG) a token sits on a vertex $v$ of a graph $H$; the player to move slides it to a neighbour $w$ and $v$ is deleted; a player who cannot move loses. The game tree from $(G,u)$ is Godsil's path tree $T(G,u)$, whose nodes are the simple paths of $G$ that start at $u$. The matching polynomial is $\mu(H,z)=\sum_k(-1)^km_k(H)z^{|H|-2k}$, where $m_k(H)$ counts matchings with $k$ edges. Sorting matchings by whether and how they cover $v$ gives the vertex recurrence
\begin{equation}\label{eq:vertex-recurrence}
\mu(H,z)=z\,\mu(H-v,z)-\sum_{w\sim_Hv}\mu(H-v-w,z).
\end{equation}

\begin{theorem}[matching collapse]\label{thm:matching-collapse}
For a finite graph $G$ and vertex $u$, let $R$ be the root resolvent of the UVG game tree at $(G,u)$. Then
\[
R(z)=\frac{\mu(G-u,z)}{\mu(G,z)},\qquad \ord_{z=0}R=2\bigl(\nu(G)-\nu(G-u)\bigr)-1.
\]
So the player to move wins UVG at $(G,u)$ if and only if $\nu(G-u)=\nu(G)-1$, that is, $u$ is covered by every maximum matching of $G$ \textup{(}Fraenkel--Scheinerman--Ullman \cite{FSU}\textup{)}. Equivalently, $e_u$ lies in the image of the Tutte matrix of $G$ over $\QQ(x_e)$, which a random substitution tests with high probability \cite{Lovasz}.
\end{theorem}

\begin{proof}
For the position $(H,v)$ take $A=\mu(H)$ and $B=\mu(H-v)$. The options of $(H,v)$ are $(H-v,w)$ for the neighbours $w$ of $v$, so the right side of the transfer identity is $z-\sum_w\mu(H-v-w)/\mu(H-v)$, which equals $\mu(H)/\mu(H-v)$ by \eqref{eq:vertex-recurrence}; \cref{lem:transfer} gives $R=\mu(G-u)/\mu(G)$. The lowest-order term of $\mu(H)$ is $\pm(\text{number of maximum matchings})\,z^{|H|-2\nu(H)}$, so $\ord_0\mu(H)=|H|-2\nu(H)$, and the formula for $\ord_0R$ follows; \cref{thm:resolvent-outcome} turns it into the outcome.
For the Tutte statement, let $Q$ be the Tutte matrix: for each edge $e=vw$ with $v$ before $w$ it has $x_e$ in position $(v,w)$ and $-x_e$ in position $(w,v)$. $Q$ is skew-symmetric of rank $2\nu(G)$ over $\QQ(x_e)$ \cite{Tutte,Lovasz}. If some kernel vector is nonzero at $u$, then row $u$ and column $u$ are combinations of the others, and deleting both keeps the rank. If every kernel vector vanishes at $u$, deleting column $u$ lowers the rank; skew-symmetric ranks are even, so deleting row $u$ as well lowers it by exactly $2$. Hence $\nu(G-u)=\nu(G)-1$ exactly when $e_u\perp\ker Q$, that is, $e_u\in\operatorname{Im}Q$.
\end{proof}

\begin{remark*}
\begin{enumerate}[label=(\arabic*)]
\item The identity $R=\mu(G-u)/\mu(G)$ is Godsil's path-tree theorem \cite{Godsil81,GodsilBook}, whose proof is essentially the induction above. Reading the path tree as the UVG game tree and the order at zero as the outcome gives a spectral derivation of the FSU theorem; we have not found this derivation in the literature. That the multiplicity of the root $0$ of $\mu$ is the deficiency is Godsil's observation; the multiplicities of the other roots have a Gallai--Edmonds theory of their own \cite{Godsil95,KuChen}.
\item The collapse is a valuation phenomenon. The coefficients of $\mu(G,z)$ are $\sharpP$-hard to compute (for even $|G|$ the constant coefficient counts perfect matchings up to sign \cite{Valiant}); only the order of vanishing at $0$ is tractable, and the outcome needs nothing more.
\item By Heilmann--Lieb \cite{HeilmannLieb}, $\mu$ is real-rooted and the roots of $\mu(G-u)$ interlace those of $\mu(G)$; the outcome is the sign of the interlacing defect at the single point $z=0$. The multivariate form of the same theorem (real stability in vertex weights) is the engine of \S\ref{sec:realstab}.
\end{enumerate}
\end{remark*}

\subsection{Where spectral collapse fails}\label{ssec:fails}
The matching potential exists because deleting the token's vertex in UVG matches deleting a vertex in the matching recurrence. Most small changes break this match, and the resulting games are $\PSPACE$-complete, so none of them has a polynomial-time evaluation of any kind unless $\Pclass=\PSPACE$. Edge geography on bipartite graphs is the exception: it stays polynomial through a different algebraic criterion, over $\mathrm{GF}(2)$.

\begin{center}
\begin{tabularx}{\linewidth}{@{}L{0.24}L{0.2}>{\raggedright\arraybackslash}X L{0.19}@{}}
\toprule
Variant & Game tree & Complexity & Source\\ \midrule
UVG, outcome & path tree & polynomial (\cref{thm:matching-collapse}) & \cite{FSU}\\
Directed (generalized) geography & directed path tree & $\PSPACE$-complete, even planar bipartite of maximum degree $3$ & \cite{LichtensteinSipser}\\
Undirected edge geography & edge-simple walks & $\PSPACE$-complete & \cite{FSU}\\
\quad on bipartite graphs & edge-simple walks & polynomial, by linear algebra over $\mathrm{GF}(2)$ & \cite{FSU}\\
UVG, mis\`ere play & path tree, terminal rule reversed & $\PSPACE$-complete & \cite{RenaultSchmidt}\\
UVG, Grundy value & path tree & $\PSPACE$-hard at maximum degree $4$ (planar bipartite); polynomial at maximum degree $3$ & \cite{BFT}\\
Sum of two UVG positions & interleaved path trees & $\PSPACE$-complete & \cite{BFT}\\
\bottomrule
\end{tabularx}
\end{center}

The last two rows are the sharpest. The valuation $\ord_0$ records outcomes but not Grundy values, and it is not additive under disjunctive sums: a tractable spectral invariant can be strictly weaker than the compositional invariant of the same game. \S\ref{ssec:directed} refines the directed row: when every strong component is symmetric, a weighted version of $\mu$ survives and decides outcomes in polynomial time, and nowhere else.

\subsection{Quantum evaluation}
The zero-energy eigenvector of \cref{thm:resolvent-outcome} can be probed by a quantum walk on the tree. That is the Farhi--Goldstone--Gutmann algorithm, which evaluates a NAND tree with $L$ leaves in time about $\sqrt L$ in the Hamiltonian-oracle model; discrete-query algorithms reach $O(\sqrt L)$ for balanced trees \cite{ACRSZ}. By \cref{thm:query} this is optimal when the leaves are an oracle. The bound says nothing about succinctly described arenas, where a polynomial-time quantum algorithm would place $\PSPACE$ in $\BQP$.

\subsection{Directed resolvent collapse}\label{ssec:directed}
In directed vertex geography (DVG) on a digraph $D$ the token moves along an out-arc $v\to w$ to an unvisited vertex $w$, and $v$ is deleted. Positions are pairs $(W,v)$ with $W$ the unvisited set, $v\in W$, and the options of $(W,v)$ are $(W-v,w)$ for out-neighbours $w\in W$. A position is \emph{reachable} if some play from some start $(V(D),s)$ reaches it. The game tree is the directed path tree, and \cref{thm:resolvent-outcome} applies to it verbatim; write $R(W,v)$ for its root resolvent.

\begin{theorem}[unweighted potentials see no orientation]\label{thm:unweighted}
Let $\mathcal D$ be a class of digraphs closed under induced subdigraphs, and let $F$ assign to each $D\in\mathcal D$ an element $F(D)$ of a field $\mathbb K\supseteq\Qz$ with $F(\emptyset)\ne0$ and
\[
F(D)=z\,F(D-v)-\sum_{w\in\nplus_D(v)}F(D-v-w)\qquad\text{for every }D\in\mathcal D\text{ and every }v\in V(D).
\]
Then every digraph in $\mathcal D$ is symmetric, and $F(D)=F(\emptyset)\,\mu(D,z)$, with $\mu$ the matching polynomial of the underlying graph.
\end{theorem}
\begin{proof}
For a single vertex the recurrence gives $F(\{x\})=zF(\emptyset)$. Suppose some $D\in\mathcal D$ has an arc $u\to v$ without $v\to u$, so that $D[\{u,v\}]\in\mathcal D$. The recurrence at $u$ gives $F(D[\{u,v\}])=(z^2-1)F(\emptyset)$, while at $v$, which has no out-neighbour there, it gives $z^2F(\emptyset)$, a contradiction. So every arc is symmetric, out-neighbourhoods are neighbourhoods, and $F/F(\emptyset)$ obeys \eqref{eq:vertex-recurrence} with the same initial value.
\end{proof}

\Cref{thm:unweighted} applies to every hereditary class with a one-way arc, acyclic digraphs included, where DVG takes linear time, so it does not separate easy classes from hard ones. It says that a recurrence with the bare weight $z$ at every vertex cannot see orientation, and that one-way arcs must enter a potential through vertex weights. \Cref{thm:directed-collapse} shows how.

\begin{remark}[transposition-invariant data]\label{rem:transposition}
Reversing every arc transposes the adjacency matrix and negates the skew-symmetric adjacency matrix of an oriented graph. So characteristic polynomials, ranks of principal submatrices, Pfaffians up to sign and Kasteleyn-type counts agree on $D$ and its reversal. DVG outcomes do not: on the out-star $u\to a$, $u\to b$ the player to move at $u$ wins, and on its reversal that player loses. Data invariant under transposition cannot decide DVG, and a Kasteleyn orientation appears below only as a device for signing matchings.
\end{remark}

\begin{definition*}[SCC-symmetry]
A digraph is \emph{SCC-symmetric} if every arc lying on a directed cycle has its reverse arc. An arc $u\to v$ lies on a directed cycle exactly when $v$ reaches $u$, that is, when $u$ and $v$ share a strong component. So the definition says that every strong component induces a symmetric digraph and the asymmetric arcs are exactly the arcs between components. The class contains every DAG, every undirected graph, and every DAG of undirected graphs joined by one-way arcs.
\end{definition*}

\begin{lemma}[entry lemma]\label{lem:entry}
In DVG on any digraph, when the token enters a strong component at a vertex $w$, no vertex reachable from $w$ has been visited. So the subgame from that moment is DVG on the vertices reachable from $w$, started at $w$. Write $\Rdown(w)$ for its root resolvent, the \emph{fresh-entry resolvent} of $w$; it equals $R(V,w)$.
\end{lemma}
\begin{proof}
The visited vertices lie on a directed path ending at the vertex $x$ from which the token entered, with $x\to w$ and $x$ outside the component of $w$. If a visited vertex $y$ were reachable from $w$, the path from $y$ to $x$, the arc $x\to w$ and a path from $w$ to $y$ would form a directed cycle through $x\to w$, and $x,w$ would share a component. At the start nothing has been visited.
\end{proof}

\begin{definition*}[local weighted matching potential]
Let $D$ be a digraph and $\mathbb K\supseteq\Qz$ a field. A \emph{local weighted matching potential} for $D$ over $\mathbb K$ is a graph $U$ on $V(D)$ with vertex weights $a\colon V(D)\to\mathbb K$ such that for every reachable position $(W,v)$ of DVG on $D$
\[
\mu_a(U[W])\neq0\qquad\text{and}\qquad R(W,v)=\frac{\mu_a(U[W-v])}{\mu_a(U[W])},
\]
where $\mu_a(H)=\sum_{M}(-1)^{|M|}\prod_{x\notin V(M)}a_x$, the sum running over the matchings $M$ of $H$. The graph $U$ and the weights are fixed for $D$ and do not depend on the position; one vertex-deletion recurrence with weights fixed in advance produces every resolvent, which \cref{lem:transfer} alone does not provide. Unless said otherwise, $\mathbb K=\Qz$.
\end{definition*}

\begin{theorem}[directed resolvent collapse]\label{thm:directed-collapse}
Let $D$ be SCC-symmetric and let $U$ be its symmetric part, the graph of the pairs joined in both directions. For a vertex $x$ in the strong component $C(x)$ define the exit self-energy and the weight
\[
\eta_x(z)=\sum_{y\in\nplus(x),\,y\notin C(x)}\Rdown(y),\qquad a_x(z)=z-\eta_x(z).
\]
Then $(U,a)$ is a local weighted matching potential for $D$ over $\Qz$. For a reachable position $(W,v)$ with $v$ in the component $C$, and with $G(C)=U[C]$, the ratio depends only on $W\cap C$:
\[
R(W,v)=\frac{\mu_a\bigl(G(C)[W\cap C-v]\bigr)}{\mu_a\bigl(G(C)[W\cap C]\bigr)}.
\]
The self-energies are computed component by component, in reverse topological order of the condensation.
\end{theorem}

\begin{proof}
$U$ has no edges between components, so $\mu_a(U[S])$ is the product of the factors $\mu_a(G(C)[S\cap C])$ over the components, and every factor except that of $C$ cancels in the ratio (they are nonzero by \cref{prop:no-cancel}). By \cref{lem:entry} every exit target of a vertex of $C$ is unvisited while the token is in $C$, so the options of a reachable position $(W,v)$ are the positions $(W-v,w)$ for neighbours $w$ of $v$ in $G(C)[W\cap C]$, together with fresh entries into later components at the exit targets $y$, whose resolvents are $\Rdown(y)$. The recursion \eqref{eq:schur} therefore reads $1/R(W,v)=z-\eta(v)-\sum_wR(W-v,w)$. Sorting matchings by whether and how they cover $v$ gives the weighted vertex recurrence $\mu_a(H)=a_v\mu_a(H-v)-\sum_w\mu_a(H-v-w)$ over neighbours $w$ of $v$. Dividing by $\mu_a(H-v)$ shows that the claimed ratio satisfies the same recursion, with the same value $1/a_v$ when $v$ has no neighbour in $W\cap C$. Induction on $|W\cap C|$ gives equality.
\end{proof}

\begin{proposition}[valuation without cancellation]\label{prop:no-cancel}
Let $X_{\Pw}\subseteq V(C)$ be the set of vertices with an exit to a P-position, that is, with some exit target $y$ at which $\Rdown(y)$ has a pole. For $S\subseteq V(C)$ let $G(C)[S]^{+}$ be $G(C)[S]$ with a new pendant leaf attached to each vertex of $S\cap X_{\Pw}$. Then $\mu_a(G(C)[S])\neq0$ and
\[
\ord_{z=0}\mu_a\bigl(G(C)[S]\bigr)=|S|-2\nu\bigl(G(C)[S]^{+}\bigr).
\]
\end{proposition}
\begin{proof}
By \cref{thm:resolvent-outcome} each $\Rdown(y)$ has either a simple pole at $0$ with positive residue or a simple zero with negative derivative. So $a_x=-r_x/z+O(1)$ with $r_x>0$ for $x\in X_{\Pw}$, and $a_x=c_xz+O(z^2)$ with $c_x\ge1$ for $x\notin X_{\Pw}$. The term of a matching $M$ with unmatched set $Y$ then has order $|Y\setminus X_{\Pw}|-|Y\cap X_{\Pw}|=|S|-2\kappa(M)$, where $\kappa(M)=|M|+|Y\cap X_{\Pw}|$, and its leading coefficient has sign $(-1)^{\kappa(M)}$. All terms of minimal order share the sign $(-1)^{\kappa^{*}}$, with $\kappa^{*}$ the maximum of $\kappa$, so they cannot cancel. Finally $\kappa^{*}=\nu(G(C)[S]^{+})$, because a matching of the augmented graph is a matching of $G(C)[S]$ plus leaf edges at some unmatched vertices of $X_{\Pw}$.
\end{proof}

\begin{corollary}[outcomes in polynomial time]\label{cor:dvg-poly}
The player to move at a reachable position $(W,v)$ with $v$ in $C$ wins if and only if $v$ is covered by every maximum matching of $G(C)[W\cap C]^{+}$. So the explicit potential is $\mu_a$, with valuation $\Lambda(S)=|S|-2\nu(G(C)[S]^{+})$. Processing the components in reverse topological order decides DVG on SCC-symmetric digraphs with one Gallai--Edmonds decomposition per component: each entry vertex $w$ is classified by whether every maximum matching of $G(C)^{+}$ covers it, and that outcome decides whether $w$'s in-neighbours upstream belong to $X_{\Pw}$.
\end{corollary}
\begin{proof}
Put $S=W\cap C$. By \cref{thm:directed-collapse,prop:no-cancel}, $\ord_0R(W,v)=\Lambda(S-v)-\Lambda(S)=2(\nu(G(C)[S]^{+})-\nu(G(C)[S]^{+}-v))-1$, because the leaf of $v$ becomes isolated. \Cref{thm:resolvent-outcome} turns the sign into the outcome.
\end{proof}

\begin{proof}[Elementary proof of the outcome statement]
By \cref{lem:entry} an exit leads to a fresh position whose outcome is fixed in advance. An exit to an N-position is a losing move, and deleting losing moves changes no outcome; an exit to a P-position wins at once, exactly like a move to a pendant leaf. So from any reachable position inside $C$ the game is UVG on $G(C)[W\cap C]^{+}$, and \cref{thm:matching-collapse} applies. The spectral proof gives more: the exact resolvent of every position and the weights $r_x,c_x$ that enter the leading coefficient of $\mu_a$ (\S\ref{ssec:pfaffian}).
\end{proof}

\begin{proposition}[where the potential fails, small cases]\label{prop:fails}
\begin{enumerate}[label=\textup{(\arabic*)}]
\item No directed cycle $C_k$, $k\ge3$, has a local weighted matching potential over any field $\mathbb K\supseteq\Qz$.
\item The three digraphs on three vertices that are strongly connected but not symmetric have none, over any field.
\item Among the $218$ digraphs on four vertices, up to isomorphism, exactly the $104$ SCC-symmetric ones have a potential, over $\Qzbar$ \textup{(}\cref{thm:exact5}\textup{)}.
\end{enumerate}
\end{proposition}
\begin{proof}
(1) From a start $s$ the play on $C_k$ runs around the cycle and stops one vertex short of $s$, so the reachable positions are the pairs $(I,v)$ with $I$ a cyclic interval of length $\ell\le k$ that begins at $v$, and the game from $(I,v)$ is a directed path with $\ell$ vertices, with resolvent $\mu(P_{\ell-1})/\mu(P_\ell)$. Let $(U,a)$ be a potential. Since $I-v$ is the interval of length $\ell-1$ that begins at the successor of $v$, the defining identity telescopes to $\mu_a(U[I])=\mu(P_\ell)$ for every cyclic interval of length $\ell$. For $\ell=1$ this gives $a\equiv z$, so $\mu_a$ is the ordinary matching polynomial. For $\ell=2$ it gives $z^2-1$ for every pair of consecutive vertices, so $U$ contains the $k$ edges of the cycle. For $\ell=k$ it gives $\mu(U)=\mu(P_k)$, whose coefficient of $z^{k-2}$ is $-(k-1)$, while in $\mu(U)$ that coefficient is minus the number of edges, at least $k$.
(2) One of the three is $C_3$. The others, with arcs $0\rightleftarrows1$, $0\rightleftarrows2$, $1\to2$ and with arcs $0\rightleftarrows1$, $1\to2\to0$, have a vertex $s$ with arcs to both ends of a one-way arc $x\to y$ ($s=0,x=1,y=2$ and $s=1,x=2,y=0$). In both, every vertex ends a Hamiltonian path, so every singleton position is reachable and forces $a\equiv z$. The positions $(\{x,y\},x)$ and $(\{x,y\},y)$ are reachable from $s$, with resolvents $z/(z^2-1)$ and $1/z$, while a potential assigns both the value $z/(z^2-[xy\in U])$.
(3) is the exact computation of \cref{thm:exact5}.
\end{proof}

\begin{theorem}[formerly Conjecture 6.10]\label{thm:conj610}
A digraph has a local weighted matching potential over $\Qz$ if and only if it is SCC-symmetric. The same holds over $\Rz$ and over the field $\Puis$ of real Puiseux series convergent at $\infty$ \textup{(}\S\ref{ssec:real-setting}\textup{)}.
\end{theorem}
\begin{proof}
``If'' is \cref{thm:directed-collapse}. ``Only if'' is \cref{thm:main-real}: a local weighted matching potential with weights in $\Puis$ is a real-stable potential (\cref{ex:real-stable}), and a digraph with a real-stable potential is SCC-symmetric.
\end{proof}

With \cref{thm:unweighted} this gives the two-level picture: with the bare weight $z$ a vertex recurrence sees no orientation, and with real vertex weights it sees exactly the orientation of arcs between strong components. Over $\Qzbar$ the statement is false: the directed triangle together with four isolated vertices has a potential with algebraic weights (\cref{thm:complex-counterexample,cor:complex-false}), so the hypothesis that the weights are real cannot be dropped.

\begin{corollary}[edge geography with bipartite components]\label{cor:edge-geo}
In edge geography on a mixed graph, undirected edges and arcs are each usable once: the token moves along an unused edge in either direction or along an unused arc in its direction, and a player who cannot move loses. Suppose every arc joins two different connected components of the undirected part and the arcs make no cycle among these components. Then the strong components, with each edge counted in both directions, are exactly the components of the undirected part. From a position with the token in the component $C$, the outcome is that of undirected edge geography on the unused edges of $C$ plus one pendant edge at each vertex of $C$ with an arc to a P-position of a later component. If every component is bipartite, outcomes are computable in polynomial time.
\end{corollary}
\begin{proof}
When the token enters a component along an arc, no edge of that component or of any component reachable from it has been used (the argument of \cref{lem:entry} applied to components). So an arc into a later component leads to a position whose outcome depends only on its head. Arcs to N-positions are losing moves and can be deleted; an arc from $x$ to a P-position wins at once, like a pendant edge at $x$ (several such arcs at $x$ act like one, since using any of them ends the play in $C$). Arcs never return to $C$, so the game inside $C$ is undirected edge geography on the unused edges of $C$ plus these pendant edges. Pendant edges keep a graph bipartite, and FSU \cite{FSU} solve undirected edge geography on bipartite graphs in polynomial time by linear algebra over $\mathrm{GF}(2)$. Processing the components in reverse topological order classifies every head.
\end{proof}

Without bipartiteness each component is an instance of undirected edge geography, which is $\PSPACE$-complete \cite{FSU}, and the reduction only transfers hardness.

\begin{remark*}
\begin{enumerate}[label=(\arabic*)]
\item For a DAG every component is one vertex and \cref{cor:dvg-poly} is backward induction; for an undirected graph there are no exits, $a\equiv z$, and \cref{thm:matching-collapse} returns. \Cref{thm:directed-collapse} glues these two collapses along the condensation.
\item The self-energy $\eta_x$ is the cavity field of the Schur-complement method: each later component acts on $C$ only through one rational function at each exit target. The number of positions stays exponential inside large components, so the collapse is not a counting of positions.
\item Unless $\Pclass=\PSPACE$, hard instances of DVG need a one-way arc inside a strong component, which is exactly where \cref{thm:conj610} places the failure of the potential. Such arcs do not imply hardness: DVG on a directed cycle is trivial yet has no potential, and \S\ref{sec:oneway} decides all digraphs whose one-way arcs share a tail in each component. The absence of a potential is not a hardness criterion.
\item DVG also takes linear time on digraphs of bounded treewidth \cite{Bodlaender} (for edge geography Bodlaender needs bounded degree as well). That class contains the directed cycles, so the algorithm does not come from a local weighted matching potential; \S\ref{sec:treewidth} explains what a width-bounded object can and cannot carry.
\item We have not found \cref{thm:directed-collapse,cor:dvg-poly,cor:edge-geo} or \cref{thm:conj610} in the literature, including the 1993 geography papers \cite{Bodlaender,FraenkelSimonson,FSU} and Fraenkel's bibliography \cite{FraenkelBib}; since \cref{cor:dvg-poly} has a short elementary proof, an earlier occurrence of the outcome statement would not be surprising.
\end{enumerate}
\end{remark*}

\subsection{The Pfaffian layer}\label{ssec:pfaffian}
The valuation of the matching potential is polynomial for every graph, but its leading coefficient already counts maximum matchings. Kasteleyn's Pfaffian method computes that coefficient on planar graphs of bounded deficiency, and it is $\sharpP$-hard without that bound. Write $d=|G|-2\nu(G)$ for the deficiency, $\mm(G)$ for the number of maximum matchings and $\operatorname{pm}(G)$ for the number of perfect matchings.

\begin{proposition}[leading coefficient and residue]\label{prop:residue}
The lowest term of $\mu(G,z)$ is $(-1)^{\nu(G)}\mm(G)z^d$. If $(G,u)$ is a P-position of UVG, then the resolvent $R=\mu(G-u)/\mu(G)$ has the residue
\[
\Res_{z=0}R(z)=\frac{\mm(G-u)}{\mm(G)}=\Pr\bigl[u\notin V(M)\bigr],
\]
where $M$ is a uniformly random maximum matching of $G$.
Every maximum matching $M$ that misses $u$ is a winning strategy for the second player: answer each move to $w$ by moving to the $M$-partner of $w$. The residue is therefore the density of matching strategies among maximum matchings.
\end{proposition}
\begin{proof}
The lowest term of $\mu(G,z)$ comes from the maximum matchings. At a P-position $\nu(G-u)=\nu(G)$, so $\mu(G-u)$ has lowest term $(-1)^{\nu(G)}\mm(G-u)z^{d-1}$, and the ratio of lowest terms is the residue. The maximum matchings of $G-u$ are exactly the maximum matchings of $G$ that miss $u$. For the strategy, suppose the opponent moves from $u$ to $w$. Then $w$ is covered by $M$, since otherwise adding $uw$ would enlarge $M$. After the answer $w\to M(w)$, the matching $M$ minus the edge $wM(w)$ is a maximum matching of $G-u-w$ that misses the token vertex: a larger one plus $uw$ would beat $\nu(G)$. The opponent again moves from a vertex missed by a maximum matching, and induction on $|G|$ finishes.
\end{proof}

\begin{theorem}[the Kasteleyn frontier]\label{thm:kasteleyn}
Let $G$ have deficiency $d$.
\begin{enumerate}[label=\textup{(\arabic*)}]
\item The valuation $\ord_0\mu(G,z)=d$ is computable in polynomial time for every graph, and it decides every outcome.
\item Every maximum matching misses exactly $d$ vertices, so $\mm(G)=\sum\operatorname{pm}(G-S)$ over $d$-element sets $S$. For planar $G$ each $G-S$ is planar and $\operatorname{pm}(G-S)$ is the absolute value of the Pfaffian of a Kasteleyn orientation \cite{TemperleyFisher,Kasteleyn}. So on planar graphs $\mm(G)$, and every residue of \cref{prop:residue}, takes $\binom{|G|}{d}$ Pfaffians: polynomial time for bounded deficiency.
\item Computing $\mm(G)$ is $\sharpP$-complete in general \cite{Valiant}, and it stays $\sharpP$-complete on planar bipartite graphs of bounded degree: Vadhan \cite{Vadhan} proves $\sharpP$-completeness there for matchings ``and extremal variants'', which include maximum matchings; this is also the ground on which Mikl\'os and Kr\'esz withdrew an independent proof \cite{MiklosKresz}. Residues are $\sharpP$-hard on the same class, because an oracle for them yields $\mm(G)$ through the telescoping identity below.
\item Computing the whole polynomial $\mu(G,z)$ is $\sharpP$-hard even for planar bipartite graphs of bounded degree: the sum of the absolute values of its coefficients is the number of matchings \cite{Jerrum,Vadhan}.
\end{enumerate}
\[
\mm(G)=\frac{\operatorname{pm}(G_d)}{\prod_{i=1}^{d}\Res_{z=0}\dfrac{\mu(G_{i-1}-u_i,z)}{\mu(G_{i-1},z)}},\qquad G_0=G,\quad G_i=G_{i-1}-u_i.
\]
\end{theorem}
\begin{proof}[Proof of item \textup{(3)}]
Choose each $u_i$ missed by some maximum matching of $G_{i-1}$; one maximum-matching computation finds it. Then $\nu$ is unchanged, the deficiency drops by one, and $(G_{i-1},u_i)$ is a P-position whose residue is $\mm(G_i)/\mm(G_{i-1})$ by \cref{prop:residue}. The product telescopes; the last graph has a perfect matching, so $\mm$ equals $\operatorname{pm}$ there, and it is planar, so its $\operatorname{pm}$ is one Pfaffian. The other items were proved above or are the cited results.
\end{proof}

The same frontier governs the weighted potential $\mu_a$ of \cref{thm:directed-collapse} on planar components: its leading coefficient is a weighted count of maximum matchings of $G(C)[S]^{+}$, with weights $r_x$ on leaf edges and $c_x$ on unmatched vertices, which are themselves residues and slopes of downstream resolvents. The potential has three layers: its valuation decides the game on every graph; its leading coefficient counts matching strategies, computable by Pfaffians on planar graphs of bounded deficiency and $\sharpP$-hard on planar bipartite graphs; and the full polynomial is $\sharpP$-hard even on planar bipartite graphs of bounded degree.

\subsection{Compositional invariants and the torsion obstruction}\label{ssec:torsion}
An invariant that respects sums and decides outcomes must separate every pair of games that some sum tells apart, so it must carry the Grundy value for impartial games and the Conway value for partizan ones. If it takes values in an abelian group and detects P-positions as zero, it is killed by $2$, and the Grundy value sits in the $2$-torsion of the group. Let $\mathfrak C$ be a class of finite games closed under disjunctive sum and negation, with outcome $o(G)$. Games are equal, $G=H$, when $o(G+X)=o(H+X)$ for every $X\in\mathfrak C$; equivalently $G-H$ is a P-position \cite{Siegel}. For impartial games $G=H$ exactly when the Grundy values $g(G),g(H)$ agree.

\begin{theorem}[compositionality forces equality]\label{thm:compositional}
Let $\varphi\colon\mathfrak C\to\Omega$ satisfy $\varphi(G+H)=\Psi(\varphi(G),\varphi(H))$ and $o(G)=f(\varphi(G))$ for some maps $\Psi,f$. Then $\varphi(G)=\varphi(H)$ implies $G=H$. For impartial games $\varphi$ determines the Grundy value.
\end{theorem}
\begin{proof}
If $\varphi(G)=\varphi(H)$ then $o(G+X)=f(\Psi(\varphi(G),\varphi(X)))=o(H+X)$ for every $X$.
\end{proof}

\begin{theorem}[torsion obstruction]\label{thm:torsion}
Let $\mathfrak C$ consist of impartial games and let $\varphi\colon\mathfrak C\to A$ be additive into an abelian group with the zero test: $G$ is a P-position iff $\varphi(G)=0$. Then $2\varphi(G)=0$ for every $G$, and $\varphi(G)=\varphi(H)$ iff $g(G)=g(H)$. So $\mathfrak C$ realizes at most $|A[2]|$ Grundy values, where $A[2]$ is the subgroup of elements of order at most $2$. In particular:
\begin{enumerate}[label=\textup{(\arabic*)}]
\item no integer-valued additive invariant, such as the degree of a divisor, has the zero test once $*1\in\mathfrak C$;
\item even with an arbitrary outcome test, an integer-valued additive invariant separates at most two Grundy values, which is impossible once $*1,*2\in\mathfrak C$;
\item if $A$ is $\operatorname{Pic}^0$ of a finite graph or of a metric graph of genus $g$, then $|A[2]|\le 2^g$.
\end{enumerate}
\end{theorem}
\begin{proof}
$G+G$ is a P-position, so $2\varphi(G)=\varphi(G+G)=0$. Then $\varphi(G)-\varphi(H)=\varphi(G)+\varphi(H)=\varphi(G+H)$ vanishes exactly when $G+H$ is a P-position, that is, when $g(G)=g(H)$; this injects Grundy values into $A[2]$, and (1) follows because $\ZZ$ has no element of order $2$.
For (2) let $L\le\ZZ$ be generated by the values $2\varphi(G)$. Adding a game $G+G$ never changes an outcome, so if $\varphi(H_1)\equiv\varphi(H_2)\pmod L$, adding suitable nonnegative multiples of such games to each makes the values equal without changing outcomes; hence the outcome depends only on $\varphi\bmod L$. The reduced invariant is additive, outcome-determining and killed by $2$, so by \cref{thm:compositional} it injects Grundy values into $(\ZZ/L)[2]$, which has at most two elements (if $L=0$ then $\varphi\equiv0$ cannot separate $0$ from $*1$). The games $0,*1,*2,*3$ need four.
For (3), $\operatorname{Pic}^0$ of a finite graph is the discriminant group of its lattice of integral flows, a lattice of rank $g$, so it is generated by $g$ elements \cite{BHN}; $\operatorname{Pic}^0$ of a metric graph is a $g$-dimensional real torus \cite{MikhalkinZharkov}. In both cases $|A[2]|\le 2^g$.
\end{proof}

Wild outcome tests can do more: an additive map into $\RR/\ZZ$ with rationally independent values on $*1,*2,*4,\dots$ carries any number of Grundy bits if the test reads off parities of coefficients. Linear equivalence and rank are zero tests, so item (3) is the bound that applies to divisors.

\begin{example}[the torsion bound is attained]\label{ex:bouquet}
Let $B_k$ be the bouquet of $k$ digons: a vertex $o$ joined to each of $p_1,\dots,p_k$ by two parallel edges, of genus $k$. Firing $p_i$ moves two chips to $o$, so the principal divisors are spanned by $2(p_i-o)$ and $\operatorname{Pic}^0(B_k)\cong(\ZZ/2)^k$ is generated by the classes of $p_i-o$. For an impartial game $G$ with Grundy value below $2^k$ write its binary digits $b_1,\dots,b_k$ and put $\Dv(G)=\sum_ib_i(p_i-o)$, so $g(G)=\sum_ib_i2^{i-1}$. The bits of $g(G)\oplus g(H)$ are $b+b'-2bb'$, so $\Dv(G+H)\sim\Dv(G)+\Dv(H)$. A divisor of degree $0$ has rank $0$ if it is equivalent to $0$ and rank $-1$ otherwise, so $G$ is a P-position exactly when $\Dv(G)\sim0$, that is, $r(\Dv(G))=0$. Genus $k$ thus carries exactly $2^k$ Grundy values and item (3) is sharp. On the metric rose of $k$ circles of length $2$, with $p_i$ antipodal to $o$, the classes of $p_i-o$ generate a subgroup $(\ZZ/2)^k$ of the torus $\operatorname{Pic}^0$: on one circle a piecewise-linear function with integer slopes whose divisor is supported on $\{o,p\}$ has slope $s$ on both arcs oriented from $o$ to $p$, so its divisor is $2s(o-p)$.
\end{example}

\begin{corollary}[rank is not compositional]\label{cor:rank-not}
In \cref{ex:bouquet} with $k\ge2$, $r(\Dv(*1))=r(\Dv(*2))=-1$, yet $r(\Dv(*1)+\Dv(*1))=0$ while $r(\Dv(*1)+\Dv(*2))=-1$, so the rank of a sum is not a function of the ranks of its summands. Rank is only superadditive: if $r(D_1),r(D_2)\ge0$, any effective divisor of degree $r(D_1)+r(D_2)$ splits into parts of degrees $r(D_1)$ and $r(D_2)$ that $D_1$ and $D_2$ absorb, so $r(D_1+D_2)\ge r(D_1)+r(D_2)$.
\end{corollary}

Sums are linear in the Jacobian; the mex is not. The divisor of $G+H$ is the sum of the divisors, but $\Dv(G)$ encodes a mex taken over the whole game tree of $G$.

\subsection{Where the new invariants turn hard}\label{ssec:hard}

\begin{theorem}[hardness transfer]\label{thm:hardness}
\begin{enumerate}[label=\textup{(\arabic*)}]
\item Fix graphs $\Gamma_n$ of size polynomial in $n$, computable from $1^n$ in polynomial time. Let $\mathfrak C_n$ consist of the UVG positions with at most $n$ vertices, the nim heaps $*m$ with $m\le n$, and the sums $G+*m$. Suppose $\Dv$ maps $\mathfrak C_n$ to divisors on $\Gamma_n$ in polynomial time, with $\Dv(G+*m)\sim\Dv(G)+\Dv(*m)$, and each game in $\mathfrak C_n$ is a P-position exactly when its divisor is equivalent to $0$. Then $\Pclass=\PSPACE$.
\item Suppose a compositional, outcome-determining invariant has values of polynomial size, polynomial-time maps $\Psi$ and $f$, and polynomial-time evaluation on single UVG positions. Then it decides sums of two UVG positions in polynomial time, and $\Pclass=\PSPACE$.
\item The same holds on Morris's family of short partizan games. In particular Conway's canonical form cannot have values of polynomial size together with a polynomial-time composition map on that family unless $\Pclass=\PSPACE$.
\end{enumerate}
\end{theorem}
\begin{proof}
(1) $\Dv\sim0$ exactly when $\Dv$ lies in the integer column span of the Laplacian of $\Gamma_n$, which a Hermite normal form decides \cite{KannanBachem}. The Grundy value of a UVG position $G$ is at most the degree of its token vertex, hence at most $n-1$, and it is the unique $m\le n$ with $G+*m$ a P-position, that is, with $\Dv(G)+\Dv(*m)\sim0$. So Grundy values of UVG positions would be computable in polynomial time, and they are $\PSPACE$-hard to compute \cite{BFT}. (2) Evaluate on both summands, combine with $\Psi$ and apply $f$; sums of two UVG positions are $\PSPACE$-complete \cite{BFT}. (3) The same argument with Morris's theorem \cite{Morris}.
\end{proof}

\begin{proposition}[rank is a two-round game]\label{prop:rank-game}
By definition $r(\Dv)\ge k$ exactly when $\Dv-E$ is equivalent to an effective divisor for every effective $E$ of degree $k$: an adversary removes $k$ chips, then the player must win the dollar game. Winnability is decidable in polynomial time through $q$-reduced divisors \cite{BakerNorine,BakerShokrieh}, so deciding $r(\Dv)\ge k$ is in $\coNP$, and computing $r(\Dv)$ is $\NP$-hard \cite{KissToth}. A polynomial-time many-one reduction from a $\PSPACE$-complete game to the question $r(\Dv)\ge k$ would give $\PSPACE=\NP$, and a polynomial-time Turing reduction to the rank function would give $\PSPACE\subseteq\Pclass^{\NP}$.
\end{proposition}
\begin{proof}
A witness that $r(\Dv)<k$ is one effective $E$ of degree $k$ with $\Dv-E$ unwinnable, checked in polynomial time. A many-one reduction places $\PSPACE$ in $\coNP$; $\PSPACE$ is closed under complement, so it lies in $\NP$ as well. The value $r(\Dv)$ is found by binary search with polynomially many $\coNP$ queries.
\end{proof}

\noindent\emph{The directed frontier.} DVG is polynomial on SCC-symmetric digraphs (\cref{cor:dvg-poly}), on digraphs whose one-way arcs share a tail in each strong component (\cref{cor:star-dvg}) and on bounded treewidth \cite{Bodlaender}; it is $\PSPACE$-complete in general, even on planar bipartite digraphs of maximum degree $3$ \cite{LichtensteinSipser}. Mixed edge geography is polynomial when its arcs run acyclically between bipartite components of the undirected part (\cref{cor:edge-geo}). Partizan vertex geography is $\PSPACE$-complete even on bipartite digraphs of bounded degree, and impartial geography with two tokens on specified start vertices is $\PSPACE$-complete even on DAGs, where one token takes linear time \cite{FraenkelSimonson}.

\begin{center}
\small
\begin{longtable}{@{}L{0.24}L{0.19}L{0.26}L{0.21}@{}}
\toprule
Invariant & Tractable & Intractable & Source\\ \midrule\endhead
Valuation of $\mu$ (UVG outcomes) & every graph & none & \cref{thm:matching-collapse}\\
Valuation of $\mu_a$ (DVG outcomes) & SCC-symmetric digraphs & general digraphs: $\PSPACE$-complete, even planar bipartite of maximum degree $3$ & \cref{cor:dvg-poly}; \cite{LichtensteinSipser}\\
Local weighted matching potential & exists exactly on SCC-symmetric digraphs (real weights) & over $\Qzbar$: on strictly more digraphs; characterization open & \cref{thm:directed-collapse,thm:conj610,thm:complex-counterexample}, \cref{prob:complex}\\
Real-stable potentials (matching, Hermitian determinantal) & exactly SCC-symmetric digraphs & never with a one-way arc on a cycle & \cref{thm:main-real}\\
DVG by dynamic programming & bounded treewidth, linear time & general hardness as above & \cite{Bodlaender}\\
Mixed edge geography & arcs acyclic between bipartite undirected components & general components: $\PSPACE$-complete & \cref{cor:edge-geo}; \cite{FSU}\\
Transposition-invariant data & none & cannot decide DVG & \cref{rem:transposition}\\
Leading coefficient of $\mu$; residues & planar, bounded deficiency (Pfaffians) & $\sharpP$-hard on planar bipartite graphs & \cref{thm:kasteleyn}; \cite{Vadhan}\\
Full matching polynomial & bounded treewidth & $\sharpP$-hard on planar bipartite graphs of bounded degree & \cite{Jerrum,Vadhan}\\
Linear equivalence, dollar game ($r\ge0$) & every graph & none & \cite{BakerShokrieh}\\
Divisor rank, $r\ge k$ & fixed $k$ & $\NP$-hard to compute; decision in $\coNP$ & \cref{prop:rank-game}; \cite{KissToth}\\
Additive divisor invariant with a zero test & computable from the Grundy value (\cref{ex:bouquet}) & $\PSPACE$-hard for UVG positions & \cref{thm:torsion,thm:hardness}\\
Compositional invariant with polynomial-size values & sums of nim heaps & sums of two UVG positions; Morris sums & \cref{thm:hardness}\\
\bottomrule
\end{longtable}
\end{center}

None of these invariants helps for the class of all arenas, where \cref{thm:alternation} applies. On structured classes algebra does remove the exponential cost: matchings and the acyclic order absorb the alternation of directed geography, Pfaffians count strategies up to bounded deficiency, and tropical divisors linearize sums but not the mex.

\section{Algebraic collapse: where exact boundaries are known}\label{sec:algebraic}

Exact collapse theorems exist, and each is stated in terms of a small algebraic or combinatorial object attached to the rules: a group, a clone of polymorphisms, a matroid, a hypergraph of bounded rank, or a decomposition that follows the order of play.

\subsection{Compositional invariants: Sprague--Grundy and its limits}
For impartial normal-play games the Grundy value $g(P)=\operatorname{mex}\{g(Q):Q\text{ an option of }P\}$ turns disjunctive sums into XOR, $g(G+H)=g(G)\oplus g(H)$, and $P$ is an N-position exactly when $g(P)\ne0$ \cite{Sprague,Grundy}; Bouton's rule for Nim \cite{Bouton} is the case where $g$ is the heap size. A homomorphism into a small group does not by itself give tractability:
\begin{itemize}
\item Grundy values of UVG are $\PSPACE$-hard although outcomes are polynomial \cite{BFT};
\item Red--Blue Hackenbush values are dyadic rationals that add under disjunctive sum, yet computing the value is $\NP$-hard \cite{WinningWays,DemaineHearn};
\item sums of individually trivial partizan games are $\PSPACE$-complete \cite{Morris}, and sums of two polynomial-time impartial games can be $\PSPACE$-complete \cite{BFT}.
\end{itemize}
An invariant helps compositionally only if it is itself polynomial-time computable and has polynomial-size representations.

\subsection{Quantified constraints: the polymorphism clone}\label{ssec:qcsp}
A quantified formula $Q_1x_1\dots Q_nx_n$ over constraints from a finite language $\Gamma$ is a game: players set the variables in order, and Max ($\exists$) wins when every constraint holds. A polymorphism of $\Gamma$ is an operation $f\colon D^k\to D$ that preserves every relation of $\Gamma$ coordinatewise. The complexity of $\mathrm{QCSP}(\Gamma)$ depends only on the clone of surjective polymorphisms \cite{BBCJK}.

\begin{theorem}[Boolean case; Schaefer, see Creignou--Khanna--Sudan]\label{thm:schaefer}
For a finite Boolean language with constants, $\mathrm{QCSP}(\Gamma)$ is in $\Pclass$ when $\mathrm{Pol}(\Gamma)$ contains $\min$ (Horn), $\max$ (dual Horn), majority (2-CNF) or minority $x\oplus y\oplus z$ (affine), and $\PSPACE$-complete otherwise \cite{SchaeferSAT,CKS01}.
\end{theorem}
The Horn case is due to Kleine B\"uning, Karpinski and Fl\"ogel \cite{KBKF}. Two of the four collapses make the global invariant explicit.

\begin{proposition}[affine collapse]\label{prop:affine}
If every constraint is a linear equation over $\mathrm{GF}(2)$, $\val$ is computable by Gaussian elimination in polynomial time.
\end{proposition}
\begin{proof}
Let $S_k\subseteq\mathrm{GF}(2)^k$ be the set of assignments to $x_1,\dots,x_k$ from which $\exists$ wins the rest of the game; $S_n$ is the affine solution set, and by induction every $S_k$ is affine or empty. Keep the system for $S_k$ in reduced echelon form with $x_k$ as the preferred pivot. If $x_k$ does not occur, the slices $x_k=0$ and $x_k=1$ both equal the projection, so $S_{k-1}$ is the projection for either quantifier. If an equation reads $x_k=\ell(x_1,\dots,x_{k-1})+c$, then for $Q_k=\exists$ the set $S_{k-1}$ is cut out by the remaining equations, and for $Q_k=\forall$ the two slices are disjoint, so $S_{k-1}=\emptyset$. Each step is one elimination on at most $n$ equations; an inconsistent system shows up as an equation $0=1$, after which the set stays empty. The formula is true exactly when $S_0\ne\emptyset$.
\end{proof}

\begin{theorem}[Aspvall--Plass--Tarjan \cite{APT}]\label{thm:apt}
Build the implication digraph of a quantified 2-CNF formula on the $2n$ literals, each clause $a\lor b$ giving the arcs $\neg a\to b$ and $\neg b\to a$. The formula is true if and only if none of the following holds:
\begin{enumerate}[label=\textup{(\arabic*)}]
\item an existential literal and its complement lie in one strong component;
\item a universal literal lies in one strong component with an existential literal quantified before it;
\item there is a path from a universal literal to a different universal literal, possibly its own complement.
\end{enumerate}
The test runs in linear time.
\end{theorem}
Here the ``global connected components'' of the problem statement appear literally: as strong components of a graph on $2n$ literals, not of the $2^n$-vertex state space.

\noindent\emph{Beyond the Boolean case.} Chen developed collapsibility and the polynomially generated powers (PGP) property: when the relevant powers of the polymorphism algebra are generated by polynomially many tuples, the universal player's options reduce to polynomially many scenarios \cite{Chen08,Chen11}. He conjectured that PGP places $\mathrm{QCSP}(\Gamma)$ in $\NP$ and that its absence (EGP) makes it $\PSPACE$-complete \cite{Chen12}. Zhuk and Martin \cite{ZhukMartin} refuted this: on three elements with constants $\mathrm{QCSP}(\Gamma)$ is in $\Pclass$, $\NP$-complete, $\coNP$-complete or $\PSPACE$-complete, and the classification exhibits EGP languages that are not $\PSPACE$-complete; they also found a $\mathsf{DP}$-complete language on four elements and a $\Theta_2^{\Pclass}$-complete one on ten, and proved the conjecture for conservative languages. Zhuk \cite{Zhuk24} then proved that for every finite constraint language $\mathrm{QCSP}(\Gamma)$ is either in $\Pi_2^{\Pclass}$ or $\PSPACE$-complete.

\begin{remark*}
For quantified constraints the exact alternative to $\PSPACE$-completeness is a collapse of alternation depth: every such game is either $\PSPACE$-complete or decidable within $\Pi_2^{\Pclass}$, with two rounds of alternation. Below that line sit $\Pclass$, $\NP$, $\coNP$, $\mathsf{DP}$, $\Theta_2^{\Pclass}$ and $\Pi_2^{\Pclass}$, so the true boundary is richer than a dichotomy.
\end{remark*}

\subsection{Matroids: the exchange axiom supplies the strategy}\label{ssec:matroid}
\begin{theorem}[Lehman \cite{Lehman}]\label{thm:lehman}
In the Shannon switching game on a graph $G$ with terminals $a,b$, Short moving second wins if and only if $G$ contains a subgraph that contains $a$ and $b$ and has two edge-disjoint spanning trees. Lehman proved the analogue for all matroids.
\end{theorem}
The condition is tested in polynomial time as a matroid partition problem \cite{EdmondsPartition}; Short moving first wins exactly when some first move leaves a second-player win, which takes at most $|E|$ further tests. The exchange axiom supplies the strategy: when Cut deletes an edge of one tree, Short claims an edge of the other tree that reconnects it. The vertex version has no such exchange structure and is $\PSPACE$-complete \cite{EvenTarjan}, and so is Hex \cite{Reisch}.

\subsection{Positional games: the rank-3/rank-4 threshold}\label{ssec:positional}
If $\sigma=(S_0,B_0)$ records the sets already claimed by two players, the arena is a positional game on a hypergraph. For Maker--Breaker games the boundary is exact, provided two recent preprints hold up.
\begin{center}\small
\begin{tabularx}{\linewidth}{@{}l>{\raggedright\arraybackslash}Xl@{}}
\toprule
Edge size & Maker--Breaker complexity & Source\\ \midrule
rank at most $3$ & polynomial, with a structural characterization & \cite{GGS} (preprint)\\
$4$-uniform, bounded degree & $\PSPACE$-complete & \cite{Galliot25} (preprint)\\
$5$-uniform & $\PSPACE$-complete & \cite{Koepke}\\
$6$-uniform & $\PSPACE$-complete & \cite{RahmanWatson}\\
rank $11$ & $\PSPACE$-complete & \cite{SchaeferGames}\\
\bottomrule
\end{tabularx}
\end{center}
On rank-$3$ hypergraphs a winning Maker wins within $O(\log n)$ rounds \cite{GGS}; that depth alone gives only an $n^{O(\log n)}$ search, so the polynomial algorithm rests on the structural characterization, not on the short horizon. For Maker--Maker games, rank $4$ and $4$-uniform hypergraphs are $\PSPACE$-complete \cite{GalliotSenizergues,Galliot25}, and rank $3$ is open.

\subsection{Width must follow the order of play}\label{ssec:width}
\begin{center}\small
\begin{tabularx}{\linewidth}{@{}>{\raggedright\arraybackslash}X>{\raggedright\arraybackslash}X>{\raggedright\arraybackslash}p{0.2\linewidth}@{}}
\toprule
Structure of the formula & Complexity & Source\\ \midrule
constant pathwidth, unbounded alternation & $\PSPACE$-complete & \cite{AtseriasOliva}\\
treewidth $k$, quantifier depth $\ell$ & $\ell$-fold exponential in $k$, times $\poly(n)$ & \cite{Chen04}\\
same parameters, lower bound & nothing significantly better than $\ell$-fold exponential in $k$, under ETH & \cite{FHP}\\
width measured along the quantifier prefix & fixed-parameter tractable & \cite{EGO,AtseriasOliva}\\
\bottomrule
\end{tabularx}
\end{center}
Static sparsity of the interaction graph, a global structural invariant of the landscape, does not collapse the minimax tree; it collapses the tree only at a cost that is a tower in the number of alternations, and under ETH the tower is necessary. What collapses the tree is width measured in space-time, along the order in which the players move.

\section{The exact boundary}\label{sec:boundary}

In full generality no decidable criterion on the rules captures the exact boundary (\cref{thm:sigma2}). Inside syntactic classes the boundary is the one \S\ref{sec:algebraic} draws.

\subsection{Collapse is equivalent to locally computable optimal play}
\begin{theorem}\label{thm:strategies}
For a family of succinct arenas with polynomial horizon, closed under moving the start to any position reachable from it, the following are equivalent:
\begin{enumerate}[label=\textup{(\arabic*)}]
\item $\val$ is computable in polynomial time;
\item both players have optimal strategies computable in polynomial time from the position.
\end{enumerate}
\end{theorem}
\begin{proof}
(1$\Rightarrow$2) At a position evaluate $\val$ at both options, which are instances of the family by closure, and move to the better one; with finite horizon this is optimal by backward induction. (2$\Rightarrow$1) Play the two strategies against each other from $\sigma_0$. Max's strategy guarantees at least $\val(\sigma_0)$ and Min's at most $\val(\sigma_0)$, so the play ends with payoff exactly $\val(\sigma_0)$; it has polynomial length and each move costs polynomial time.
\end{proof}

\noindent\emph{Unbounded horizon.} Direction (2$\Rightarrow$1) weakens to (2$\Rightarrow\val\in\PSPACE$): the simulated play is deterministic, so a counter up to $2^N$ detects a cycle in polynomial space. Zero-player arenas have trivial strategies and are $\PSPACE$-complete, so (2$\Rightarrow$1) fails unless $\Pclass=\PSPACE$. Direction (1$\Rightarrow$2) needs more than $\val$, because a player on a $+1$ position must also make progress; that is the ranking of \cref{thm:local}.

\begin{corollary}\label{cor:chess}
For an $\EXPTIME$-complete family such as generalized chess, at least one player has no polynomial-time optimal strategy unless $\EXP=\PSPACE$.
\end{corollary}

\subsection{No decidable exact criterion}
A \emph{uniform family} is given by a clocked polynomial-time generator $\mathfrak M$ that maps $(1^n,y)$, $|y|=n$, to a succinct arena and a start position. Its value function $\val[\mathfrak M]$ sends $y$ to the value of that arena at that start. Call $\mathfrak M$ \emph{collapsible} if $\val[\mathfrak M]$ is computable in polynomial time.

\begin{theorem}\label{thm:sigma2}
The set of collapsible generators is $\Sigma^0_2$-complete, in particular undecidable. For generators of clocked arenas with polynomial horizon it is $\Pi^0_1$-hard if $\Pclass\ne\PSPACE$, and it contains every such generator if $\Pclass=\PSPACE$.
\end{theorem}
\begin{proof}
Upper bound: $\mathfrak M$ is collapsible if and only if some clocked polynomial-time machine $\mathfrak M'$ agrees with $\val[\mathfrak M]$ at every input $y$; the inner condition is decidable for each $y$, because each arena is finite, so the statement is $\Sigma^0_2$.
Hardness: we reduce $\mathrm{FIN}=\{e:W_e\text{ finite}\}$, which is $\Sigma^0_2$-complete. Fix an effective enumeration $(\mathfrak N_i)$ of clocked polynomial-time machines (every polynomial-time language is decided by some $\mathfrak N_i$), and a language $L\in\EXPTIME$ on which every $\mathfrak N_i$ errs at all but finitely many input lengths: on the $i$-th string of length $n$, for $i\le\log n$, simulate $\mathfrak N_i$ for $2^n$ steps and answer the opposite (strings with $i>\log n$ are rejected). Let $y\mapsto\arena(y)$ be the reduction of $L$ to $\val$ from \cref{thm:alternation}, with the alternating machine clocked so that $\val(\arena(y))=\pm1$. Given $e$, the generator $\mathfrak M(e)$ on input $(1^n,y)$ runs the enumeration of $W_e$ for $n$ steps; if a new element appears at step $n$ it outputs $\arena(y)$, and otherwise a trivial arena of value $+1$. If $W_e$ is finite, $\val[\mathfrak M(e)]$ is constant at all large lengths, so $\mathfrak M(e)$ is collapsible. If $W_e$ is infinite, $\val[\mathfrak M(e)]$ encodes $L$ at infinitely many lengths, and a polynomial-time algorithm for it would decide $L$ correctly at those lengths, contrary to the choice of $L$.
For polynomial horizon take $L$ $\PSPACE$-complete and reduce the complement of the halting problem: $\mathfrak M(H)$ on $(1^n,y)$ outputs $\arena(y)$ if the machine $H$ halts within $n$ steps on the empty input and a trivial arena otherwise. If $H$ never halts, $\val[\mathfrak M(H)]$ is constant; if it halts, $\val[\mathfrak M(H)]$ agrees with $L$ at all large lengths and is not polynomial-time unless $\Pclass=\PSPACE$. If $\Pclass=\PSPACE$, every generator with polynomial horizon is collapsible.
\end{proof}

No decidable property of generators characterizes collapsibility, and \cref{thm:sigma2} places the question at the $\Sigma^0_2$ level of the arithmetical hierarchy, like the index sets of complexity classes studied by H\'ajek \cite{Hajek} and Kozen \cite{Kozen}. Exact boundaries are therefore stated inside syntactic classes: polymorphism clones for constraint languages, rank for hypergraphs, exchange structure for matroidal winning families, width along the prefix, and, for directed geography, SCC-symmetry (\cref{thm:conj610}) and common tails (\S\ref{sec:oneway}).

\subsection{The phase table}\label{ssec:phase}
Alternation, horizon and width along the order of play govern the complexity of $\val$. A block of real choice is a maximal stretch of play in which only one player has real choices, the other making dummy moves.

\begin{center}
\small
\begin{longtable}{@{}L{0.36}L{0.32}L{0.25}@{}}
\toprule
Regime & Complexity of $\val$ & Polynomial-time $\Phi$?\\ \midrule\endhead
No real choices, polynomial horizon & $\Pclass$-complete & yes\\
No real choices, unbounded horizon & $\PSPACE$-complete & iff $\Pclass=\PSPACE$\\
$k$ blocks of real choice, polynomial horizon & $\Sigma_k^{\Pclass}$-complete if Max chooses first, $\Pi_k^{\Pclass}$-complete if Min does & iff $\Pclass=\NP$\\
Unbounded alternation, polynomial horizon & $\PSPACE$-complete & iff $\Pclass=\PSPACE$\\
Unbounded alternation, unbounded horizon & $\EXPTIME$-complete & no (deterministic, unconditional); randomized and quantum open\\
Landscape (oracle) model, $h=\omega(\log N)$ & $2^h$ deterministic, $\Theta(2^{0.7537h})$ randomized, $\Theta(2^{h/2})$ quantum queries & no, unconditionally, in the query model\\
Explicit state graph & $\Pclass$-complete, linear time & yes; in $\NC$ only if $\NC=\Pclass$\\
Constant pathwidth, unbounded alternation & $\PSPACE$-complete & iff $\Pclass=\PSPACE$\\
Treewidth $k$, $\ell$ alternations & $\ell$-fold exponential in $k$, times $\poly(n)$ & yes for fixed $k,\ell$; tower tight under ETH\\
Quantified constraints, finite language & in $\Pi_2^{\Pclass}$ or $\PSPACE$-complete & case by case (\S\ref{ssec:qcsp})\\
Maker--Breaker, rank at most $3$ & polynomial & yes\\
Maker--Breaker, $4$-uniform & $\PSPACE$-complete & iff $\Pclass=\PSPACE$\\
Directed geography, SCC-symmetric & polynomial (\cref{cor:dvg-poly}) & yes\\
Directed geography, one-way arcs with a common tail per component & polynomial (\cref{cor:star-dvg}) & yes\\
Directed geography, planar bipartite, maximum degree $3$ & $\PSPACE$-complete & iff $\Pclass=\PSPACE$\\
Partizan vertex geography, bounded degree & $\PSPACE$-complete & iff $\Pclass=\PSPACE$\\
Undirected vertex geography & polynomial (\cref{thm:matching-collapse}) & yes\\
Nim and Wythoff with binary heap sizes & polynomial \cite{Bouton,Wythoff} & yes\\
Shannon switching game on graphs & polynomial \cite{Lehman,EdmondsPartition} & yes\\
Undirected edge geography, bipartite & polynomial \cite{FSU} & yes\\
Mixed edge geography, arcs acyclic between bipartite components & polynomial (\cref{cor:edge-geo}) & yes\\
\bottomrule
\end{longtable}
\end{center}

\section{Real stability forces reciprocity: the proof of Theorem 6.10}\label{sec:realstab}

This section proves that a potential whose generating polynomial is real stable, for instance a local weighted matching potential with real weights, exists only on SCC-symmetric digraphs. The argument uses three facts: the Heilmann--Lieb theorem, which makes the generating polynomial of a real matching potential real stable; Br\"and\'en's criterion, which makes its Rayleigh differences nonnegative on all of $\RR^n$; and the expansion of resolvents at $z=\infty$, in which a one-way arc on a directed cycle leaves a first-order trace that no real-stable polynomial can carry.

\subsection{Setting}\label{ssec:real-setting}
Let $\Puis$ be the field of real Puiseux series in $z^{-1}$ that converge for all sufficiently large real $z$: series $f=\sum_{q\ge q_0,\,q\in\frac1N\ZZ}c_qz^{-q}$ with real coefficients and some $N\ge1$. It contains $\Rz$ (Laurent expansion at $\infty$) and the real algebraic functions that are real for large real $z$. A nonzero $f\in\Puis$ has a leading term $c\,z^{-q}$ with $c\neq0$; we write $f=O(z^{-q})$ if $f=0$ or its leading exponent is at least $q$, and $f=o(z^{-q})$ if the leading exponent is larger than $q$. Every $f\in\Puis$ has a real value $f(z)$ for all large real $z$, and $f(z)\neq0$ for all large $z$ when $f\ne0$.

A polynomial $P\in\RR[h_1,\dots,h_n]$ is \emph{real stable} if $P\equiv0$ or $P(h)\neq0$ whenever $\Im h_i>0$ for all $i$. Specializing a variable to a real number, differentiating, scaling the variables by a positive number, translating them by a real vector and dividing by a nonzero real number preserve real stability (the first by Hurwitz's theorem). For a multi-affine $P$ and $i\ne j$ let
\[
\Delta_{ij}P=\partial_iP\,\partial_jP-P\,\partial_i\partial_jP,
\]
a polynomial in the variables other than $h_i,h_j$ (the \emph{Rayleigh difference}).

\begin{lemma}[Br\"and\'en's criterion, necessity \cite{Branden}]\label{lem:branden}
If $P$ is multi-affine and real stable, then $\Delta_{ij}P(x)\ge0$ for all $i\ne j$ and all $x\in\RR^n$.
\end{lemma}
\begin{proof}
Specialize the variables other than $h_i,h_j$ to the real coordinates of $x$. The result $a+bh_i+ch_j+dh_ih_j$ is real stable or zero, and $\Delta_{ij}P(x)=bc-ad$. Suppose $bc-ad<0$. If $d\neq0$, the zero set is $h_j=(-bh_i-a)/(dh_i+c)$, a real M\"obius transformation of determinant $ad-bc>0$, which maps the upper half-plane into itself; a point $h_i$ of the upper half-plane gives a zero with both coordinates there. If $d=0$ then $b,c\ne0$ with $b/c<0$, and $h_j=-(a+bh_i)/c$ has $\Im h_j=-(b/c)\Im h_i>0$. Both contradict real stability. (Br\"and\'en proves the converse as well.)
\end{proof}

\begin{theorem}[Heilmann--Lieb \cite{HeilmannLieb}]\label{thm:hl}
For a finite graph $U$ with nonnegative edge weights $w$, the multivariate matching polynomial
\[
\mu_{y,w}(U)=\sum_M(-1)^{|M|}\prod_{e\in M}w_e\prod_{x\notin V(M)}y_x
\]
is real stable in the vertex variables $y$.
\end{theorem}
\noindent(Heilmann and Lieb show that $\sum_M\prod_{e\in M}w_e\prod_{x\notin V(M)}x_x$ has no zero with all $\operatorname{Re}x_x>0$; the substitution $x=iy$ and complex conjugation give the stated form.)

\subsection{Real-stable potentials}
Throughout, $D$ is a digraph with vertex set $V$ and $\eta\colon V\to\Puis$ is a family of \emph{self-energies} with $\eta_v=O(z^{-1})$; plain DVG has $\eta\equiv0$. The game $(D,\eta)$ has the positions and moves of DVG and the root resolvents
\[
\frac1{R(W,v)}=z-\eta_v-\sum_{w\in\nplus(v)\cap(W-v)}R(W-v,w).
\]
For a directed path $P=(t_0,\dots,t_k)$ of $D$ write $\Gamma(P)=\prod_{i}R(V\setminus\{t_0,\dots,t_{i-1}\},t_i)$ for the product of resolvents along the play $P$ started at $(V,t_0)$, and $\back(P)$ for the number of arcs $t_j\to t_i$ with $i<j$. A set $X\subseteq V$ is a \emph{path set} if it is the vertex set of a directed path.

\begin{definition}[real-stable potential]\label{def:real-stable}
A \emph{real-stable potential} for $(D,\eta)$ is a map $\Mpot\colon2^V\to\Puis$ such that
\begin{enumerate}[label=\textup{(P\arabic*)}]
\item at every reachable position $(W,v)$: $\Mpot(W)\neq0$ and $\Mpot(W-v)=R(W,v)\,\Mpot(W)$;
\item for all sufficiently large real $z$ the following polynomial is real stable:
\[
P_z(h)=\sum_{X\subseteq V}\Mpot(V\setminus X)(z)\,h^X.
\]
\end{enumerate}
\end{definition}

\begin{example}[real-stable potentials]\label{ex:real-stable}
\begin{enumerate}[label=(\alph*)]
\item \emph{Matching potentials with real weights.} Let $U$ be a graph on $V$ with nonnegative real edge weights $w$ and vertex weights $a\in\Puis^V$, and $\Mpot(S)=\mu_{a,w}(U[S])$. Because $\mu$ is multi-affine in the vertex weights and $\partial_{y_x}\mu_{y,w}(U)=\mu_{y,w}(U-x)$, $P_z(h)=\mu_{a(z)+h,w}(U)$, which is real stable by \cref{thm:hl}. With unit edge weights this is the local weighted matching potential of \S\ref{ssec:directed}.
\item \emph{Frozen environments.} With $U$ on $V\sqcup F$ and $\Mpot(S)=\mu_{a,w}(U[S\cup F])$, $P_z(h)=\mu_{(a_V+h,\,a_F),w}(U)$ is a specialization of a real stable polynomial.
\item \emph{Hermitian determinantal potentials.} Let $A(z)$ be a matrix indexed by $V\sqcup F$ whose entries have real and imaginary parts in $\Puis$ and which is Hermitian for all large real $z$, and put $\Mpot(S)=\det A(z)[S\cup F]$. Expanding the determinant along the diagonal gives $P_z(h)=\det\bigl(A(z)+\sum_{v\in V}h_vE_{vv}\bigr)$, where $E_{vv}$ is a diagonal matrix unit. A polynomial $\det(B+\sum_ih_iA_i)$ with $B$ Hermitian and every $A_i$ positive semidefinite is real stable or identically zero \cite[Prop.~2.4]{BorceaBranden}: if $\sum_iA_i$ is positive definite, then for $\Im h>0$ and $x\neq0$ the number $x^{*}(B+\sum_ih_iA_i)x$ has positive imaginary part, and the general case follows by replacing $A_i$ with $A_i+\varepsilon I$ and letting $\varepsilon\to0$ (Hurwitz). Since $P_z(0)=\Mpot(V)(z)\neq0$, $P_z$ is real stable. By Cramer's rule, \textup{(P1)} says that every resolvent is a diagonal entry of the inverse of a principal submatrix, $R(W,v)=\bigl(A[W\cup F]^{-1}\bigr)_{vv}$: the potential is the Green's function of a Hermitian (lossless) system.
\end{enumerate}
\end{example}

\begin{theorem}[real-stable potentials force SCC-symmetry]\label{thm:main-real}
If $(D,\eta)$ has a real-stable potential, then every arc of $D$ that lies on a directed cycle has its reverse arc; that is, $D$ is SCC-symmetric.
\end{theorem}

\begin{corollary}\label{cor:characterization}
For a digraph $D$ the following are equivalent:
\begin{enumerate}[label=\textup{(\roman*)}]
\item $D$ is SCC-symmetric;
\item $D$ has a local weighted matching potential over $\Qz$;
\item $D$ has a local weighted matching potential over $\Puis$ \textup{(}in particular over $\Rz$\textup{)}, even with nonnegative real edge weights;
\item $D$ has a real-stable potential.
\end{enumerate}
No digraph that is not SCC-symmetric has a Hermitian determinantal potential \textup{(}\cref{ex:real-stable}\textup{(c))}.
\end{corollary}
\begin{proof}
(i)$\Rightarrow$(ii) is \cref{thm:directed-collapse}; (ii)$\Rightarrow$(iii) is trivial; (iii)$\Rightarrow$(iv) is \cref{ex:real-stable}(a); (iv)$\Rightarrow$(i) is \cref{thm:main-real}. The last sentence is \cref{ex:real-stable}(c) with \cref{thm:main-real}.
\end{proof}

\subsection{Four lemmas}

\begin{lemma}[resolvents at infinity]\label{lem:asymptotics}
Let $e_v$ be the coefficient of $z^{-1}$ in $\eta_v$. For every position, $R(W,v)=z^{-1}+\bigl(d^{+}_W(v)+e_v\bigr)z^{-3}+o(z^{-3})$, where $d^{+}_W(v)=|\nplus(v)\cap(W-v)|$. Consequently, for a directed path $P$,
\[
z^{|V(P)|}\,\Gamma(P)=1+\ell(P)\,z^{-2}+o(z^{-2}),\qquad \ell(P)=\sum_{t\in V(P)}\bigl(d^{+}(t)+e_t\bigr)-\back(P).
\]
\end{lemma}
\begin{proof}
By induction on $|W|$ every child resolvent is $z^{-1}+O(z^{-3})$, and $\eta_v=e_vz^{-1}+o(z^{-1})$, so $1/R(W,v)=z-(d^{+}_W(v)+e_v)z^{-1}+o(z^{-1})$, which inverts to the claim. (Self-energies in $\Puis$ may contain exponents between $1$ and $2$, so the error terms cannot be improved to $O(z^{-4})$ and $O(z^{-3})$ in general; for $\eta\equiv0$ they can, see \S\ref{ssec:fields-setting}.) In the product, the $i$-th factor has $d^{+}_{W_i}(t_i)$ equal to $d^{+}(t_i)$ minus the number of out-neighbours of $t_i$ among $t_0,\dots,t_{i-1}$; summing gives $\ell(P)$.
\end{proof}

Let $Z$ be a directed cycle $q_0\to q_1\to\dots\to q_{m-1}\to q_0$ with vertex set $T$, $m\ge3$. A \emph{cyclic interval} of $Z$ is a set of consecutive vertices of $Z$, including $\emptyset$ and $T$.

\begin{lemma}[tameness]\label{lem:tameness}
Let $\gamma_X\in\Puis$ for $X\subseteq T$, and suppose that for every sufficiently large real $z$ the polynomial $G_z(v)=\sum_{X}\gamma_X(z)v^X$ is real stable. If $\gamma_I=1+O(z^{-2})$ for every cyclic interval $I$ of $Z$, then $\gamma_X=1+O(z^{-2})$ for every $X\subseteq T$.
\end{lemma}

\begin{proof}
Fix $r\in T$ and let $J=T\setminus\{r\}$, linearly ordered along $Z$ starting after $r$. A subset of $J$ that is contiguous in this order is a cyclic interval of $Z$. For $Y\subseteq J$ let $\operatorname{hull}(Y)$ be the smallest contiguous subset containing $Y$ and $\operatorname{gap}(Y)=|\operatorname{hull}(Y)\setminus Y|$. We prove $\gamma_Y=1+O(z^{-2})$ for all $Y\subseteq J$ by induction on $(\operatorname{gap}(Y),|Y|)$ in lexicographic order; every $X\ne T$ lies in some such $J$, and $T$ is an interval.
If $\operatorname{gap}(Y)=0$, $Y$ is a cyclic interval. Otherwise let $p<q$ be the first and last elements of $Y$, choose $g\in\operatorname{hull}(Y)\setminus Y$ (so $p<g<q$), and put $S=Y\setminus\{p,q\}$. The seven sets $S$, $S+p$, $S+q$, $S+g$, $S+p+g$, $S+q+g$ and $Y+g$ are subsets of $J$ that precede $Y$ in the order: $Y+g$ has one gap fewer; $S$, $S+p=Y-q$ and $S+q=Y-p$ are smaller and have no more gaps than $Y$; and $S+g$, $Y-q+g$, $Y-p+g$ have strictly fewer gaps, because moving an endpoint of $Y$ onto $g$ shortens the hull by more than it adds. By induction their coefficients are $1+O(z^{-2})$.
Now $\partial^S G_z$ is real stable, and so is its restriction $f(v_p,v_q,v_g)$ to the variables $v_p,v_q,v_g$ with all others set to $0$; its coefficients are $\gamma_{S\cup Y'}$ for $Y'\subseteq\{p,q,g\}$. Write $\gamma_Y=1-\delta$. By \cref{lem:branden}, for every real $t$,
\begin{multline*}
\Delta_{pq}f(t)=\bigl(\gamma_{S+p}+\gamma_{S+p+g}t\bigr)\bigl(\gamma_{S+q}+\gamma_{S+q+g}t\bigr)-\bigl(\gamma_S+\gamma_{S+g}t\bigr)\bigl(\gamma_Y+\gamma_{Y+g}t\bigr)\\
=\mathsf a+\mathsf bt+\mathsf ct^2\ge0,
\end{multline*}
with $\mathsf a=\delta(1+O(z^{-2}))+O(z^{-2})$, $\mathsf b=\delta(1+O(z^{-2}))+O(z^{-2})$ and $\mathsf c=O(z^{-2})$. A real quadratic that is nonnegative on $\RR$ has $\mathsf c\ge0$, $\mathsf a\ge0$ and $\mathsf b^2\le4\mathsf a\mathsf c$. If $\delta$ were not $O(z^{-2})$, its leading exponent would be some $q<2$, and then $\mathsf b^2$ would have leading exponent $2q$ while $4\mathsf a\mathsf c=O(z^{-q-2})$ with $q+2>2q$; so $\mathsf b^2>4\mathsf a\mathsf c$ for large $z$, a contradiction. Hence $\delta=O(z^{-2})$.
\end{proof}

\begin{lemma}[first-order factorization]\label{lem:first-order}
Let $T$ be finite, $x\ne y\in T$, $R=T\setminus\{x,y\}$, and let $G_z(v)=\sum_{X\subseteq T}\gamma_Xv^X$ with $\gamma_X=1+\ell_Xz^{-2}+o(z^{-2})$ for real numbers $\ell_X$. Then for every $v\in\RR^R$,
\[
\Delta_{xy}G_z(v)=z^{-2}\,Q(v)F(v)+o(z^{-2}),\qquad Q(v)=\prod_{r\in R}(1+v_r),\quad F(v)=\sum_{S\subseteq R}\kappa(S)v^S,
\]
where $\kappa(S)=\ell_{S+x}+\ell_{S+y}-\ell_{S+x+y}-\ell_S$.
\end{lemma}
\begin{proof}
Write $G_z=\prod_{t\in T}(1+v_t)+z^{-2}E_z$ with $E_z=\sum_X(\ell_X+o(1))v^X$, and decompose any multi-affine $G$ as $G_0+v_xG_x+v_yG_y+v_xv_yG_{xy}$ with $G_0,G_x,G_y,G_{xy}$ free of $v_x,v_y$; then $\Delta_{xy}G=G_xG_y-G_0G_{xy}$. For the product all four parts equal $Q$, so
\begin{multline*}
\Delta_{xy}\bigl(\textstyle\prod(1+v_t)+\varepsilon E\bigr)=(Q+\varepsilon E_x)(Q+\varepsilon E_y)-(Q+\varepsilon E_0)(Q+\varepsilon E_{xy})\\
=\varepsilon\,Q\,(E_x+E_y-E_{xy}-E_0)+\varepsilon^2(E_xE_y-E_0E_{xy}).
\end{multline*}
With $\varepsilon=z^{-2}$ the coefficient of $v^S$ in $E_x+E_y-E_{xy}-E_0$ is $\kappa(S)+o(1)$.
\end{proof}

\begin{lemma}[vanishing]\label{lem:vanishing}
Let $Q=\prod_{r\in R}(1+v_r)$ and let $F$ be multi-affine in $v_R$. If $Q(v)F(v)\ge0$ for all $v\in\RR^R$, then $F=cQ$ for a constant $c\ge0$.
\end{lemma}
\begin{proof}
Induction on $|R|$; for $R=\emptyset$, $F$ is a constant $c=QF\ge0$. Take $r\in R$ and write $F=F^0+v_rF^1$ and $Q=(1+v_r)Q'$ with $F^0,F^1,Q'$ free of $v_r$. At a point $u$ of the other coordinates with $Q'(u)\ne0$, the quadratic $v_r\mapsto(1+v_r)(F^0(u)+v_rF^1(u))Q'(u)$ is nonnegative and vanishes at $v_r=-1$, so it cannot change sign there, and $F^0(u)=F^1(u)$. Such points are dense, so $F^0=F^1$ and $F=(1+v_r)F^1$. Then $Q'F^1\ge0$ off the hyperplane $v_r=-1$, hence everywhere, and induction gives $F^1=cQ'$.
\end{proof}

\subsection{Proof of \texorpdfstring{\cref{thm:main-real}}{Theorem 9.5}}
Let $\Mpot$ be a real-stable potential and suppose that an arc $x\to y$ lies on a directed cycle while $y\to x$ is not an arc. Choose a directed path $y=q_0\to q_1\to\dots\to q_k=x$; then $k\ge2$, and together with $x\to y$ it forms a directed cycle $Z$ with vertex set $T=\{q_0,\dots,q_k\}$. Every cyclic interval of $Z$ is the vertex set of a sub-path of $Z$, hence a path set.

For $X\subseteq V$ put $c(X)=\Mpot(V\setminus X)/\Mpot(V)\in\Puis$. If $X=V(P)$ for a directed path $P$ started at $(V,t_0)$, every position of the play $P$ is reachable, and telescoping \textup{(P1)} along it gives $c(X)=\Gamma(P)$; in particular $\Gamma(P)$ depends only on $V(P)$. For large $z$ the polynomial $P_z(h)/\Mpot(V)(z)=\sum_Xc(X)(z)h^X$ is real stable by \textup{(P2)}; setting $h_u=0$ for $u\notin T$ and rescaling $h=zv$ shows that
\[
G_z(v)=\sum_{X\subseteq T}\gamma_Xv^X,\qquad \gamma_X=z^{|X|}c(X),
\]
is real stable. By \cref{lem:asymptotics}, $\gamma_X=1+\ell(P)z^{-2}+o(z^{-2})$ whenever $X=V(P)$; in particular $\gamma_I=1+O(z^{-2})$ on the cyclic intervals of $Z$. By \cref{lem:tameness}, $\gamma_X=1+O(z^{-2})$ for every $X\subseteq T$, so $\gamma_X=1+\ell_Xz^{-2}+o(z^{-2})$ with $\ell_X\in\RR$ (and $\ell_X=\ell(P)$ for path sets).

Apply \cref{lem:first-order} to the pair $(x,y)$ with $R=\{q_1,\dots,q_{k-1}\}$. The eight sets entering $\kappa(\emptyset)$ and $\kappa(R)$ are path sets: $\emptyset$, $\{x\}$, $\{y\}$, $\{x,y\}$ (the arc $x\to y$), $R$, $R+y$, $R+x$ and $T$ (the sub-paths $(q_1,\dots,q_{k-1})$, $(q_0,\dots,q_{k-1})$, $(q_1,\dots,q_k)$ and $(q_0,\dots,q_k)$). In each combination the sums $\sum(d^{+}(t)+e_t)$ cancel, so only back-arc counts remain:
\begin{align*}
\kappa(\emptyset)&=\back(x,y)=[\,y\to x\,]=0,\\
\kappa(R)&=\back(q_0,\dots,q_k)-\back(q_1,\dots,q_k)\\
&\qquad-\back(q_0,\dots,q_{k-1})+\back(q_1,\dots,q_{k-1})\\
&=[\,q_k\to q_0\,]=[\,x\to y\,]=1,
\end{align*}
because an arc between $q_0$ and $R$, or between $R$ and $q_k$, is counted equally often with each sign, and the only remaining pair is $\{q_0,q_k\}$, whose arc $q_k\to q_0$ points backward along $(q_0,\dots,q_k)$. Hence $F(0)=\kappa(\emptyset)=0$ while the coefficient of $v^R$ in $F$ is $1$, so $F$ is not a constant multiple of $Q$. By \cref{lem:vanishing} there is $v^{*}\in\RR^R$ with $Q(v^{*})F(v^{*})<0$, and by \cref{lem:first-order} $\Delta_{xy}G_z(v^{*})<0$ for all large $z$. This contradicts \cref{lem:branden}. \qed

\begin{remark}[what the proof uses]\label{rem:locality}
Only the play on one directed cycle through the one-way arc enters, together with the starts on that cycle. Nothing is assumed about $U$, about the rest of $D$, about the sign or degree of the weights, or about a set potential; the ``free'' coefficients $\gamma_X$ on sets that are not path sets, which \S11.7 of the earlier draft identified as the obstacle, are controlled by \cref{lem:tameness}. The quantities $\kappa(R)=1$ and $\kappa(\emptyset)=0$ are the two claims in the proof of \cref{cor:sign-coherent}; what was missing there is \cref{lem:first-order,lem:vanishing}, which turn them into a contradiction without any sign assumption, by evaluating the Rayleigh difference at real points where $Q(v)F(v)<0$.
\end{remark}

\subsection{Frozen environments}
Let $C$ be a digraph with self-energies $\eta$ as above; in the reduction of \cref{prop:restriction}, $\eta_x$ is the sum of the fresh-entry resolvents of the exits of $x$, which is $O(z^{-1})$. A \emph{frozen potential} for $(C,\eta)$ with environment $F$ is a graph $U$ on $V(C)\sqcup F$ with vertex weights $a$ (and edge weights $w\ge0$) such that $\Mpot_F(W)=\mu_{a,w}(U[W\cup F])$ satisfies \textup{(P1)} for the game $(C,\eta)$.

\begin{theorem}[formerly Conjecture 10.9$'$]\label{thm:frozen-real}
If $(C,\eta)$ has a frozen potential with weights in $\Puis$, for any finite environment, then $C$ is SCC-symmetric. In particular no strongly connected, non-symmetric digraph, with any real self-energies vanishing at infinity, has a real frozen potential.
\end{theorem}
\begin{proof}
$\Mpot_F$ is a real-stable potential for $(C,\eta)$ by \cref{ex:real-stable}(b); apply \cref{thm:main-real}.
\end{proof}

\begin{proposition}[restriction to a component]\label{prop:restriction}
Let $(U,a)$ be a local weighted matching potential for $D$ over a field $\mathbb K\supseteq\Qz$, and $C$ a strong component. Then $(U,a)$ is a frozen potential for $(D[C],\eta)$ with environment $V\setminus C$, where $\eta_x=\sum\Rdown(y)$ over the exits $y$ of $x$.
\end{proposition}
\begin{proof}
In a play that starts in $C$ no vertex outside $C$ is visited while the token stays in $C$, because it cannot leave $C$ and return. So the unvisited set has the form $S\cup(V\setminus C)$ with $S\subseteq C$, and by \cref{lem:entry} every exit leads to a fresh subgame; hence $R(S\cup(V\setminus C),v)$ is the resolvent of the game $(D[C],\eta)$ at $(S,v)$.
\end{proof}

\Cref{prop:restriction} is the reduction by which the earlier draft derived Theorem~6.10 from the frozen statement; \cref{thm:main-real} makes it unnecessary for real weights. Over $\Qzbar$ it runs the other way: \S\ref{sec:complex} finds a counterexample by realizing a complex frozen environment with genuine vertices.

\subsection{Sharpness and scope}

\begin{proposition}[the band is sharp for the triangle, for every environment]\label{prop:triangle-band}
Let $C_3$ be the directed triangle $0\to1\to2\to0$ with $\eta=0$, and let $z$ be real with $|z|>2$. Then no real stable multi-affine polynomial $G(v)=\sum_{X\subseteq\{0,1,2\}}\gamma_Xv^X$ has $\gamma_X=z^{|X|}\Gamma(P)(z)$ for the path data of $C_3$. Consequently, at every real $z$ with $|z|>2$ there is no frozen potential of $C_3$ with real weights, for any environment and any graph $U$. For real $z$ with $0<|z|<2$ and $z\notin\{\pm1,\pm\sqrt2\}$ a real frozen potential exists \textup{(}\cref{thm:triangle-solved}\textup{)}.
\end{proposition}
\begin{proof}
Every subset of $\{0,1,2\}$ is a path set, with $\gamma_\emptyset=1$, $\gamma_{\{i\}}=g_1=(s+1)/s$, $\gamma_{\{i,j\}}=\gamma_{\{0,1,2\}}=g_2=(s+2)/s$, where $s=z^2-2$ (\cref{thm:triangle-solved}). For the pair $(x,y)=(2,0)$ and $v_1=t$,
\[
\Delta_{20}G(t)=(g_1+g_2t)^2-(1+g_1t)(g_2+g_2t)=\frac{1+(s+2)(t+t^2)}{s^2},
\]
and at $t=-\frac12$ this equals $(4-z^2)/(4(z^2-2)^2)<0$. By \cref{lem:branden} $G$ is not real stable. A real frozen potential at $z$ would make $G$ real stable (\cref{ex:real-stable}(b), with the path data evaluated at $z$).
\end{proof}

\begin{remark}\label{rem:scope}
\begin{enumerate}[label=(\arabic*)]
\item \emph{Reality is essential.} Over $\Qzbar$ frozen environments rescue every directed cycle (\cref{thm:every-cycle}), and genuine vertices realize them: the directed $m$-cycle with $m+1$ isolated vertices has a potential over $\Qzbar$ (\cref{thm:complex-counterexample}), so \cref{thm:conj610} fails for complex weights. In both cases some environment weights are non-real at every large real $z$ (\cref{thm:every-cycle}(4), \cref{prop:Ym-arithmetic}(2)), in agreement with \cref{thm:frozen-real,thm:main-real}.
\item \emph{Lossless means reciprocal.} By \cref{ex:real-stable}(c) and \cref{thm:main-real}, the diagonal Green's functions $\bigl(A[W\cup F]^{-1}\bigr)_{vv}$ of a Hermitian matrix, however large its hidden part $F$, never reproduce the resolvents of a digraph with a one-way arc on a directed cycle. In physical language, one-way transport cannot be engineered from reciprocal couplings without loss or gain; the complex weights of \S\ref{sec:complex} play the role of the non-Hermitian elements of reservoir engineering \cite{MetelmannClerk}. The earlier draft stated this only as an analogy.
\item \emph{Earlier partial results.} For weights in $\Puis$, \cref{thm:main-real} contains the normal and sign-coherent cases (\cref{thm:normal,cor:sign-coherent}), the unilateral case (\cref{cor:unilateral}), directed triangles (\cref{cor:triangles}), isolated directed $4$-cycles (\cref{cor:c4}), and directed cycles with symmetric environments (\cref{thm:meanfield}). \Cref{thm:normal,thm:no-phantom,thm:cocycle,thm:locality} hold over arbitrary fields and remain relevant over $\Qzbar$; \cref{prop:path-lemma}, \cref{lem:frozen-path,thm:meanfield} give pointwise or structural information that \cref{thm:main-real} does not.
\item \emph{No effective threshold.} The proof excludes real-stable completions with coefficients in $\Puis$, which is an asymptotic statement. An effective version, with uniform constants in the induction of \cref{lem:tameness} and a compactness argument in \cref{lem:vanishing}, would give for each non-SCC-symmetric $D$ an explicit $z_D$ such that the path data of $D$ admit no real-stable completion at any real $z>z_D$; we have not written it out. For the directed triangle $z_D=2$ (\cref{prop:triangle-band}). An explicit bound for general cycles is \cref{prob:band}.
\end{enumerate}
\end{remark}

\section{Potentials over arbitrary fields}\label{sec:fields}

\Cref{thm:main-real} settles real weights. This section collects what holds over an arbitrary field $\mathbb K\supseteq\Qz$, in particular over $\Qzbar$, where \S\ref{sec:complex} shows that the conclusion of \cref{thm:conj610} fails. Four mechanisms are available over every field: a balance law at $z=\infty$, a cocycle carried by the game alone, locality in $U$, and the cavity identity. Three further obstructions use real weights through sums of squares or Rayleigh differences at a few points. For weights in $\Puis$ they are special cases of \cref{thm:main-real}; we keep their short proofs because each isolates one mechanism and shows exactly where reality enters, which is what \S\ref{sec:complex} exploits.

\subsection{Setting}\label{ssec:fields-setting}
Throughout, $D$ is a digraph on $V$ with $\eta\equiv0$ unless said otherwise, $(U,a)$ is a local weighted matching potential over a field $\mathbb K\supseteq\Qz$, and
\[
\Mpot(S)=\mu_a(U[S])\quad(S\subseteq V),\qquad c(X)=\frac{\Mpot(V\setminus X)}{\Mpot(V)},
\]
and $\rho(S,s)=\Mpot(S-s)/\Mpot(S)$ when $\Mpot(S)\neq0$.
We call $\rho$ the \emph{Green's function} of the potential; at reachable positions $\rho=R$. For a vertex $v$ its \emph{exits} are $\Delta^{+}(v)=\nplus(v)\setminus N_U(v)$ and its \emph{phantom neighbours} are $\Delta^{-}(v)=N_U(v)\setminus\nplus(v)$. Two vertices are \emph{comparable} if one reaches the other, and $D$ is \emph{unilateral} if any two of its vertices are comparable. As in \S\ref{sec:realstab}, $\Gamma(P)$ is the product of resolvents along a play $P$ started at $(V,t_0)$, and $\back(P)$ counts the arcs $t_j\to t_i$ with $i<j$.

\begin{lemma}[field independence]\label{lem:field-independence}
For a digraph $D$ and a graph $U$ on $V(D)$ the following are equivalent: \textup{(i)} $D$ has a potential $(U,a)$ over some field $\mathbb K\supseteq\Qz$; \textup{(ii)} it has one over $\Qzbar$; \textup{(iii)} it has one with weights in the field $\CC\{\!\{z^{-1}\}\!\}$ of complex Puiseux series in $z^{-1}$ that converge for all large $|z|$.
\end{lemma}
\begin{proof}
Introduce unknowns $a_x$ $(x\in V)$ and $t$. The polynomials $\Mpot(W-v)-R(W,v)\Mpot(W)$ over the reachable positions $(W,v)$, together with $t\prod_W\Mpot(W)-1$, lie in $\Qz[a,t]$, and their common zeros in $\mathbb K^{V}\times\mathbb K$ are exactly the potentials over $\mathbb K$ with graph $U$. If they have a common zero in some $\mathbb K\supseteq\Qz$, the ideal they generate in $\Qz[a,t]$ does not contain $1$, and by Hilbert's Nullstellensatz it has a zero over $\Qzbar$. So (i) implies (ii). By the Newton--Puiseux theorem $\Qzbar$ embeds into $\CC\{\!\{z^{-1}\}\!\}$ (every algebraic function has convergent Puiseux expansions at $z=\infty$), so (ii) implies (iii), and (iii) implies (i) trivially.
\end{proof}

So ``a potential over $\Qzbar$'', ``a potential over $\CC(z)$'s algebraic closure'' and ``a potential over some field'' mean the same thing, and the existence question for fixed $U$ is decided by one Gr\"obner basis over $\Qz$ (\cref{thm:exact5}). For a nonzero Puiseux series $f$ we write $\deg f$ for the exponent of its leading term at $z=\infty$ (so $\deg z=1$), and $f=O(z^{-q})$ as in \S\ref{ssec:real-setting}; when the weights lie in $\Qzbar$ these notions refer to a fixed embedding into $\CC\{\!\{z^{-1}\}\!\}$, and every statement below holds for each embedding. \Cref{lem:asymptotics} holds verbatim for self-energies that are complex Puiseux series, with the same $o(\cdot)$ error terms; for $\eta\equiv0$ each $R(W,v)$ is an odd function of $z$, so its error terms improve to $R(W,v)=z^{-1}+d^{+}_W(v)z^{-3}+O(z^{-5})$ and $z^{|V(P)|}\Gamma(P)=1+\ell(P)z^{-2}+O(z^{-4})$ with $\ell(P)=\sum_{t\in V(P)}d^{+}(t)-\back(P)$.

\begin{lemma}[telescoping]\label{lem:telescoping}
If $P$ is a directed path of $D$, then $c(V(P))=\Gamma(P)$. In particular $\Gamma(P)$ depends only on the vertex set of $P$.
\end{lemma}
\begin{proof}
Every position of the play $P$ from $(V,t_0)$ is reachable, and multiplying the identities $\Mpot(W-v)=R(W,v)\Mpot(W)$ along the play gives $\Mpot(V\setminus V(P))=\Gamma(P)\Mpot(V)$.
\end{proof}

\subsection{Three identities}

\begin{lemma}[cavity identity]\label{lem:cavity}
At every reachable position $(W,v)$,
\[
a_v-z=\sum_{w\in\Delta^{-}(v)\cap W}\rho(W-v,w)\;-\sum_{t\in\Delta^{+}(v)\cap W}R(W-v,t).
\]
\end{lemma}
\begin{proof}
$\Mpot(W-v)=R(W,v)\Mpot(W)\neq0$, because resolvents of positions are nonzero rational functions. Sorting matchings by how they cover $v$ gives $\Mpot(W)=a_v\Mpot(W-v)-\sum_{w}\Mpot(W-v-w)$ over the $U$-neighbours $w\in W$ of $v$; dividing by $\Mpot(W-v)$ gives $1/R(W,v)=a_v-\sum_w\rho(W-v,w)$. The game recursion gives $1/R(W,v)=z-\sum_tR(W-v,t)$ over the out-neighbours $t\in W$ of $v$. If $w$ is both a $U$-neighbour and an out-neighbour, $(W-v,w)$ is reachable, so $\rho(W-v,w)=R(W-v,w)$ and the two terms cancel.
\end{proof}

\begin{lemma}[weighted two-point identity]\label{lem:two-point}
For every graph $U$, all vertex weights in a commutative ring, every $S\subseteq V(U)$ and $x\neq y$ in $S$,
\[
\Mpot(S-x)\,\Mpot(S-y)-\Mpot(S)\,\Mpot(S-x-y)=\sum_{P}\Mpot\bigl(S\setminus V(P)\bigr)^{2},
\]
the sum running over the paths $P$ of $U[S]$ from $x$ to $y$.
\end{lemma}
\begin{proof}
Induction on $|S|$. Expand $\Mpot(S)$ and $\Mpot(S-y)$ by the vertex recurrence at $x$; the terms with $a_x$ cancel. What remains is $\Mpot(S-x-y)^2$ when $xy\in U$, plus, for each $U$-neighbour $w\neq y$ of $x$, the left side of the identity for $S-x$ and the pair $(w,y)$. By induction that is the sum over paths from $w$ to $y$ avoiding $x$, and prefixing $x$ gives every other path from $x$ to $y$.
\end{proof}

This is the identity of Heilmann--Lieb and Godsil \cite{HeilmannLieb,GodsilBook}; the proof uses only the vertex recurrence. Write $\tau(S;x,y)$ for the left side divided by $\Mpot(S)^2$. When $\Mpot(S)$ and $\Mpot(S-x)$ are nonzero,
\begin{equation}\label{eq:green-difference}
\rho(S,y)-\rho(S-x,y)=\frac{\tau(S;x,y)}{\rho(S,x)}.
\end{equation}
For weights in $\Puis$, $\tau(S;x,y)$ is a sum of squares of elements of $\Puis$, hence nonnegative at all large real $z$.

\subsection{Obstructions over every field}

\begin{theorem}[normal potentials]\label{thm:normal}
Let $(U,a)$ be a potential for $D$ with weights in $\CC\{\!\{z^{-1}\}\!\}$ and $\deg a_x\ge1$ for every $x$. Then $D$ is SCC-symmetric, $U$ contains every symmetric pair of $D$, and every other edge of $U$ joins two incomparable vertices.
\end{theorem}
\begin{proof}
\emph{Step 1.} A matching $N$ of $U[S]$ contributes $\prod_{x\in S}a_x$ times $O(z^{-2|N|})$, because each covered vertex removes a factor of degree at least $1$. So $\Mpot(S)=\bigl(\prod_{x\in S}a_x\bigr)(1+O(z^{-2}))$ and $\rho(S,s)=a_s^{-1}(1+O(z^{-2}))$. At the reachable position $(V,x)$ this equals $R(V,x)=z^{-1}(1+O(z^{-2}))$, so $a_x=z+O(z^{-1})$ and $\rho(S,s)=z^{-1}+O(z^{-3})$ for all $S$ and $s$.

\emph{Step 2 (balance).} At a reachable position $(W,v)$ the vertex recurrence gives $\Mpot(W)/\Mpot(W-v)=a_v-|N_U(v)\cap(W-v)|\,z^{-1}+O(z^{-3})$, and the game recursion gives $1/R(W,v)=z-d^{+}_W(v)z^{-1}+O(z^{-3})$. So the coefficient of $z^{-1}$ in $a_v-z$ equals $|N_U(v)\cap(W-v)|-d^{+}_W(v)$ for every reachable $(W,v)$. Comparing with $W=V$ gives the \emph{balance law}
\[
|\nplus(v)\cap X|=|N_U(v)\cap X|\qquad\text{for every reachable position }(V\setminus X,v).
\]

\emph{Step 3.} Let $x\to y$ be an arc and let $y$ reach $x$. Take a shortest path $y=p_0\to p_1\to\dots\to p_k=x$. For $0\le j\le k$ the position $(V\setminus\{p_j,\dots,p_{k-1}\},x)$ is reachable from $p_j$, so subtracting the balance laws for $j$ and $j+1$ gives $[p_j\in\nplus(x)]=[p_j\in N_U(x)]$; at $j=0$ this gives $xy\in U$. If $y\not\to x$, the position $(V-x,y)$ is reachable from $x$, and the balance law at $y$ with $X=\{x\}$ reads $0=[xy\in U]=1$. So every arc inside a strong component lies in $U$ and has its reverse. The same suffix argument applied to a $U$-edge $xw$ with $w$ reaching $x$ gives $x\to w$, so $x,w$ lie in one strong component and form a symmetric pair; by symmetry of $U$ the same holds when $x$ reaches $w$.
\end{proof}

\begin{theorem}[no phantom neighbour]\label{thm:no-phantom}
Let the weights lie in any field. If a vertex $v$ has no phantom neighbour, then $a_v=z-\sum_{t\in\Delta^{+}(v)\cap W}R(W-v,t)$ at every reachable position $(W,v)$, and no exit of $v$ reaches $v$.
\end{theorem}
\begin{proof}
The formula is \cref{lem:cavity} with $\Delta^{-}(v)=\emptyset$. By \cref{lem:asymptotics} its right side is $z-|\Delta^{+}(v)\cap W|\,z^{-1}+O(z^{-3})$, and $a_v$ does not depend on the position, so $|\Delta^{+}(v)\cap W|$ is the same at every reachable position $(W,v)$; at the start $(V,v)$ it equals $|\Delta^{+}(v)|$. If an exit $t$ of $v$ reached $v$, a shortest path from $t$ to $v$ would give a reachable position $(W',v)$ with $t\notin W'$, and then $|\Delta^{+}(v)\cap W'|<|\Delta^{+}(v)|$.
\end{proof}

\begin{theorem}[back-arc cocycle]\label{thm:cocycle}
Let the weights lie in any field, or more generally let $\Mpot$ be any \emph{set potential}: a function on the unvisited sets of reachable positions, with values in a field containing $\Qz$, such that $\Mpot(W)\neq0$ and $\Mpot(W-v)=R(W,v)\Mpot(W)$ at every reachable $(W,v)$. Then any two directed paths with the same vertex set have the same number of back arcs. In particular, if $D[X]$ has a directed Hamiltonian cycle, then $D[X]$ is Eulerian: every vertex of $X$ has as many out-neighbours as in-neighbours in $X$.
\end{theorem}
\begin{proof}
By \cref{lem:telescoping} both paths give $\Gamma=c(X)$, so their $z^{-2}$ terms in \cref{lem:asymptotics} agree, and the out-degree sums are equal. Moving the first vertex $x_i$ of a rotation of the Hamiltonian cycle to the end changes the back count by $d^{+}_X(x_i)-d^{-}_X(x_i)$.
\end{proof}

This settles \cref{prop:fails}(2) with no reference to $U$: both non-cycle digraphs there have the Hamiltonian cycle $0\to1\to2\to0$ and a vertex with two out-neighbours and one in-neighbour. Directed cycles are Eulerian and need the argument of \cref{prop:fails}(1).

\begin{theorem}[locality]\label{thm:locality}
Let the weights lie in any field. For every arc $u\to v$ inside a strong component, $u$ and $v$ lie in one component of $U$. Hence every strong component lies inside one component of $U$.
\end{theorem}
\begin{proof}
Since $\mu_a$ factorizes over the components of $U[W]$, and every factor is nonzero when $\Mpot(W)\neq0$, $\rho(W,u)$ depends only on the component of $u$ in $U[W]$. Suppose that $u$ and $v$ lie in different components of $U$. Take a shortest path $v=p_0\to\dots\to p_k=u$, put $W_1=V\setminus\{p_1,\dots,p_{k-1}\}$ and $W_0=W_1-v$. The positions $(W_1,u)$ and $(W_0,u)$ are reachable, from $p_1$ and from $v$ respectively, and the components of $u$ in $U[W_1]$ and $U[W_0]$ coincide, so their Green's functions agree. Their resolvents differ: $1/R(W_1,u)-1/R(W_0,u)=-R(W_1-u,v)=-z^{-1}+O(z^{-3})$.
\end{proof}

\subsection{Obstructions that use real weights}
For weights in $\Puis$ each result of this subsection follows from \cref{thm:main-real}; the direct proofs use only the two-point identity or Rayleigh differences at three points. The notation $O(\cdot)$ refers to $\Puis$.

\begin{proposition}[path lemma]\label{prop:path-lemma}
Let the weights lie in $\Puis$. If $t_0\to t_1\to\dots\to t_m$ is a directed path with $m\ge1$ and $t_m\not\to t_0$, then $t_0t_m\notin U$.
\end{proposition}
\begin{proof}
Put $T=\{t_1,\dots,t_{m-1}\}$ and $S=V\setminus T$. The sets $T$, $T+t_0$, $T+t_m$ and $T+t_0+t_m$ are the vertex sets of the directed paths $(t_1..t_{m-1})$, $(t_0..t_{m-1})$, $(t_1..t_m)$ and $(t_0..t_m)$, so their values $c(\cdot)$ are products $\Gamma$ (\cref{lem:telescoping}). \Cref{lem:two-point} at $S$, divided by $\Mpot(V)^2$, gives
\begin{multline*}
c(T{+}t_0)\,c(T{+}t_m)-c(T)\,c(T{+}t_0{+}t_m)=\sum_{P}\Bigl(\frac{\Mpot(S\setminus V(P))}{\Mpot(V)}\Bigr)^{2}\\
\ge\;c(T{+}t_0{+}t_m)^{2}\quad\text{if }t_0t_m\in U,
\end{multline*}
at every large real $z$. In \cref{lem:asymptotics} the out-degree sums cancel, and the back counts of the four paths differ only by the arcs from $T$ to $t_0$, from $t_m$ to $T$ and $t_m\to t_0$; the combination leaves $[t_m\to t_0]=0$. So the left side is $O(z^{-2|T|-6})$, while the right side is $z^{-2|T|-4}(1+O(z^{-2}))$, which is impossible at large $z$.
\end{proof}

\begin{corollary}\label{cor:u-edges}
Let the weights lie in $\Puis$. Every edge of $U$ is a symmetric pair inside one strong component or joins two incomparable vertices.
\end{corollary}
\begin{proof}
If $x$ reaches $w$ and $xw\in U$, \cref{prop:path-lemma} applied to a path from $x$ to $w$ forces $w\to x$; then $w$ reaches $x$, and the arc $w\to x$ with the edge $wx$ forces $x\to w$.
\end{proof}

\begin{corollary}[unilateral digraphs]\label{cor:unilateral}
Let the weights lie in $\Puis$. In every potential each tail $x$ of a one-way arc inside a strong component has a $U$-neighbour incomparable to $x$. Hence a unilateral digraph, in particular a strongly connected digraph or one with a directed Hamiltonian path, has a potential only if it is SCC-symmetric.
\end{corollary}
\begin{proof}
Let $x\to y$ be one-way inside a strong component $C$, and suppose every $U$-neighbour of $x$ is comparable to $x$. By \cref{cor:u-edges} they are symmetric partners, hence out-neighbours, so $x$ has no phantom neighbour. By \cref{thm:no-phantom} no exit of $x$ reaches $x$, so $y$ is not an exit and $xy\in U$, contradicting \cref{cor:u-edges}.
\end{proof}

\begin{corollary}[sign-coherent potentials]\label{cor:sign-coherent}
Let the weights lie in $\Puis$ and suppose that for every $S\subseteq V$ the value $\Mpot(S)$ is nonzero with the sign of $\Mpot(V)$ at all large real $z$. Then $D$ is SCC-symmetric. Normal real potentials are sign-coherent, since $\Mpot(S)=z^{|S|}(1+o(1))$, and so is the potential of \cref{thm:directed-collapse}.
\end{corollary}
\begin{proof}
Let $x\to y$ be one-way inside a strong component, take a shortest path $y=p_0\to\dots\to p_k=x$ (so $k\ge2$) and put $T=\{p_1,\dots,p_{k-1}\}$.

\emph{Claim 1: $\tau(V\setminus T;x,y)=z^{-4}(1+O(z^{-2}))$.} The sets $T$, $T+x$, $T+y$, $T+x+y$ are the vertex sets of the paths $(p_1..p_{k-1})$, $(p_1..x)$, $(y,p_1..p_{k-1})$ and $(y,p_1..x)$, with back counts $b$, $b+B_x$, $b+B_y$ and $b+B_x+B_y+1$, where $B_x$ counts arcs from $x$ into $T$ and $B_y$ arcs from $T$ into $y$, and the last $1$ is the arc $x\to y$. By \cref{lem:asymptotics}, $c(T{+}x)c(T{+}y)-c(T)c(T{+}x{+}y)=z^{-2|T|-4}(1+O(z^{-2}))$, while $c(T)^2=z^{-2|T|}(1+O(z^{-2}))$.

\emph{Claim 2: $\tau(V;x,y)=c(x)c(y)-c(x,y)=O(z^{-6})$}, because the only possible back arc of the path $(x,y)$ would be $y\to x$.

\emph{Claim 3: $\tau(V\setminus T;x,y)\le\tau(V;x,y)$ at large real $z$.} Every Green's function is positive at large real $z$ by sign-coherence, so \eqref{eq:green-difference} gives $\rho(S-u,b)\le\rho(S,b)$ whenever the sets involved are nonempty. For a path $P$ of $U$ from $x$ to $y$ that avoids $T$, remove its vertices $q_1,\dots,q_r$ one at a time:
\begin{align*}
0<\frac{\Mpot(V\setminus T\setminus V(P))}{\Mpot(V\setminus T)}&=\prod_i\rho\bigl(V\setminus T\setminus\{q_1,\dots,q_{i-1}\},q_i\bigr)\\
&\le\prod_i\rho\bigl(V\setminus\{q_1,\dots,q_{i-1}\},q_i\bigr)=\frac{\Mpot(V\setminus V(P))}{\Mpot(V)}.
\end{align*}
So each term of $\tau(V\setminus T;x,y)$ is at most the term of the same path in $\tau(V;x,y)$, whose other terms are nonnegative. Claims 1 and 2 contradict Claim 3 at large $z$.
\end{proof}

\begin{corollary}[directed triangles]\label{cor:triangles}
Let the weights lie in $\Puis$. Let $i\to j$ be a one-way arc and $k$ a vertex adjacent to both $i$ and $j$ such that $D[\{i,j,k\}]$ has a directed Hamiltonian path. Then $\back\{i,j,k\}=\back\{i,k\}+\back\{j,k\}$, where $\back X$ is the back count of any directed Hamiltonian path of $D[X]$ \textup{(}well defined by \cref{thm:cocycle}\textup{)}. Consequently every strongly connected set of three vertices induces a symmetric subdigraph, and no one-way arc inside a strong component lies on a directed triangle; no assumption on $U$ is needed.
\end{corollary}
\begin{proof}
By \cref{thm:hl} and \cref{ex:real-stable}(a), at every large real $z$ the polynomial $f(h)=\mu_{a+h}(U)/\mu_a(U)$ is real stable, and so is its restriction to $h_i,h_j,h_k$ with the other variables set to $0$. Its coefficients are the $c(T)$, $T\subseteq\{i,j,k\}$, and each such $T$ is the vertex set of a directed path. Write $f=A+Bh_i+Ch_j+Dh_ih_j$ with $A,B,C,D$ affine in $h_k$. By \cref{lem:branden}, for every real $h_k$,
\[
BC-AD=\gamma+\beta h_k+\alpha h_k^{2}\ge0,\qquad\text{so}\qquad \beta^{2}\le4\alpha\gamma,
\]
with $\gamma=c(i)c(j)-c(i,j)$, $\beta=c(i)c(j,k)+c(j)c(i,k)-c(i,j,k)-c(k)c(i,j)$ and $\alpha=c(i,k)c(j,k)-c(k)c(i,j,k)$. By \cref{lem:asymptotics} the out-degree sums cancel, $\gamma=\back\{i,j\}z^{-4}+O(z^{-6})=O(z^{-6})$, $\alpha=O(z^{-6})$, and $\beta=\bigl(\back\{i,j,k\}-\back\{i,k\}-\back\{j,k\}\bigr)z^{-5}+O(z^{-7})$. If the bracket were nonzero, $\beta^2$ would have order $z^{-10}$ while $4\alpha\gamma=O(z^{-12})$. A strongly connected non-symmetric triple is a directed triangle or one of the two digraphs of \cref{prop:fails}(2); the bracket equals $1$ for the triangle, and the other two are not Eulerian (\cref{thm:cocycle}).
\end{proof}

\begin{corollary}[an isolated directed $4$-cycle]\label{cor:c4}
Let the weights lie in $\Puis$. If a chordless directed $4$-cycle $x_0\to x_1\to x_2\to x_3\to x_0$ is a connected component of $D$, then $D$ has no potential, whatever the rest of $D$ and $U$.
\end{corollary}
\begin{proof}
Plays inside the cycle ignore the rest of $D$, so $c(T)=\mu(P_{4-|T|})/\mu(P_4)$ for every cyclic interval $T$. As in \cref{cor:triangles}, $g(h)=\sum_{T\subseteq\{x_0,\dots,x_3\}}c(T)h^T$ is real stable at large real $z$, hence so is $P(y)=\mu(P_4)\,g(y-z\one)$. Let $Q$ be the multivariate matching polynomial of the undirected $4$-cycle in the variables $y$. The Taylor coefficients of $Q$ at $y=z\one$ are $\mu(C_4-T)$; they agree with those of $P$ on every nonempty interval, and on $\emptyset$ they give $\mu(C_4)=\mu(P_4)-(z^2-1)$. Only $\{x_0,x_2\}$ and $\{x_1,x_3\}$ are not intervals, so for real $p,q$ depending on $z$,
\[
P(y)=Q(y)+(z^{2}-1)+(p-1)(y_0-z)(y_2-z)+(q-1)(y_1-z)(y_3-z).
\]
Direct expansion gives $\Delta_{02}P=1-p\,\mu(P_4)$ at $y_1=y_3=z$, so $p\le1/\mu(P_4)$ by \cref{lem:branden}, and likewise $q\le1/\mu(P_4)$. Next $\Delta_{13}P=(1-q)(z^{2}p+1)$ at $y_0=y_2=0$, so $p\ge-z^{-2}$, and likewise $q\ge-z^{-2}$. Finally, at $y_0=-1$ and $y_3=1$,
\[
\Delta_{12}P=4-z^{2}+z^{2}(p+q)+z(p-q)+pq\,(z^{2}-1),
\]
and for $z\ge3$ and $p,q\in[-z^{-2},1/\mu(P_4)]$ the last three terms are each at most $1$, so $\Delta_{12}P\le7-z^{2}<0$, contradicting \cref{lem:branden}. (The three expansions are checked symbolically by \code{scripts/check_c4.py}.)
\end{proof}

\subsection{Exact verification on at most five vertices}

\begin{theorem}[exact verification]\label{thm:exact5}
A digraph with at most five vertices has a local weighted matching potential over $\Qzbar$ if and only if it is SCC-symmetric.
\end{theorem}
\begin{proof}[Method \textup{(computer-assisted)}]
For fixed $D$ and $U$ the weights are the unknowns. The polynomials $\Mpot(W-v)-R(W,v)\Mpot(W)$ over reachable positions generate an ideal of $\Qz[a_1,\dots,a_n]$; saturating it by every $\Mpot(W)$ discards solutions with a vanishing denominator. By \cref{lem:field-independence} the saturated ideal is the unit ideal exactly when no potential with graph $U$ exists over any field, so no test point is involved.
\begin{itemize}
\item Three and four vertices: all $16$ and $218$ digraphs against all $8$ and $64$ graphs $U$. The $13$ and $104$ SCC-symmetric digraphs have potentials; for the $3$ and $114$ others every saturated ideal is the unit ideal.
\item Five vertices: $8{,}055$ of the $9{,}608$ digraphs are not SCC-symmetric. A potential over any field gives a set potential in the sense of \cref{thm:cocycle}, and whether a set potential exists is decided exactly over $\Qz$ by propagating $\Mpot$ from $V$ along the reachable positions; $7{,}976$ of the $8{,}055$ have none. For the other $79$ every graph $U$ contained in the symmetric part of $D$ is excluded by \cref{thm:no-phantom} (a tail $x$ of a one-way arc $x\to y$ inside a strong component would have no phantom neighbour and the exit $y$ reaching it), and all $80{,}798$ remaining saturated ideals are the unit ideal.
\end{itemize}
The computations are those of the original verification package, reproduced in \code{verification/legacy} (\cref{app:verification}).
\end{proof}

\subsection{The shape of a counterexample over an arbitrary field}

\begin{proposition}\label{prop:shape}
Let $D$ be a digraph that is not SCC-symmetric and has a potential $(U,a)$ over a field $\mathbb K\supseteq\Qz$. Then:
\begin{enumerate}[label=\textup{(\arabic*)}]
\item $D$ has at least six vertices \textup{(\cref{thm:exact5})};
\item any two directed paths with the same vertex set have the same back count, and every vertex set spanning a directed Hamiltonian cycle induces an Eulerian subdigraph \textup{(\cref{thm:cocycle})};
\item every strong component lies inside one component of $U$ \textup{(\cref{thm:locality})};
\item for no embedding of the weights into $\CC\{\!\{z^{-1}\}\!\}$ do all weights have degree at least $1$ \textup{(\cref{thm:normal})};
\item if $x\to y$ is a one-way arc inside a strong component and $xy\notin U$, then $x$ has a phantom neighbour \textup{(\cref{thm:no-phantom})};
\item for every strong component $C$, $(U,a)$ is a frozen potential for $(D[C],\eta)$ with environment $V\setminus C$ and the exit self-energies $\eta$ of \cref{thm:directed-collapse} \textup{(\cref{prop:restriction}, whose proof uses no property of the field)};
\item $D$ has no potential with all weights in $\Puis$ \textup{(\cref{thm:conj610})}.
\end{enumerate}
\end{proposition}

Items (1)--(6) are all the constraints we know over $\Qzbar$, and the counterexamples of \S\ref{sec:complex} meet them with seven vertices. There the real-weight obstructions visibly fail: the directed triangle $D_3$ of \cref{thm:complex-counterexample} has one-way arcs on a directed triangle, which \cref{cor:triangles} forbids for real weights, and two of its seven weights are non-real at every large real $z$ (\cref{prop:Ym-arithmetic}).

\section{The complex frontier: Theorem 6.10 fails over \texorpdfstring{$\Qzbar$}{the algebraic closure}}\label{sec:complex}

By \cref{lem:field-independence}, ``over $\Qzbar$'', ``over the algebraic closure of $\CC(z)$'' and ``over some field containing $\Qz$'' are the same question. This section answers it: \emph{the complex form of \cref{thm:conj610} is false} (\cref{thm:complex-counterexample}). The route runs through frozen environments. Over $\Qzbar$ a frozen environment rescues every directed cycle (\cref{thm:every-cycle}), and the earlier draft asked whether genuine vertices of a digraph, whose own positions must also satisfy the potential identities, can realize such an environment. They can: a directed $m$-cycle together with $m+1$ isolated vertices has a potential whose cycle weights vanish and whose environment weights are the roots of an explicit polynomial $Y_m$. The mechanism is Vieta's formulas: with vanishing cycle weights every identity prescribes an elementary symmetric function of the environment weights, and over an algebraically closed field every monic polynomial splits. Over $\RR$ the same prescription asks for a real-rooted polynomial, and $Y_m$ is not real-rooted at large real $z$, in accordance with \cref{thm:main-real}.

\subsection{Frozen potentials over an arbitrary field}

\begin{definition}\label{def:frozen}
Let $C$ be a digraph with self-energies $\eta\colon V(C)\to\mathbb K$, where $\mathbb K\supseteq\Qz$ is a field; the game $(C,\eta)$ has the positions of DVG on $C$ and the resolvents $1/R(W,v)=z-\eta_v-\sum_{w}R(W-v,w)$. A \emph{frozen potential} for $(C,\eta)$ with environment $F$ (a finite set disjoint from $V(C)$) is a graph $U$ on $V(C)\sqcup F$ with weights $a\colon V(C)\sqcup F\to K$ such that $\Mpot_F(W)=\mu_a(U[W\cup F])$ satisfies $\Mpot_F(W)\neq0$ and $\Mpot_F(W-v)=R(W,v)\Mpot_F(W)$ at every reachable position $(W,v)$ of $(C,\eta)$.
\end{definition}

The environment is never visited and acts only through $U$. By \cref{prop:restriction}, whose proof uses no property of the field, a potential for $D$ restricts to a frozen potential for each strong component with environment the rest of $D$. For a directed cycle $C$ with $\eta=0$ the reachable unvisited sets are the cyclic intervals, and $R(W,v)=\mu(P_{\ell-1})/\mu(P_\ell)$ on an interval of length $\ell$ that begins at $v$ (proof of \cref{prop:fails}(1)); so a frozen potential is exactly a function $\Mpot_F$ with $\Mpot_F(I)=\lambda\,\mu(P_{|I|})(z)$ on every cyclic interval $I$, for a constant $\lambda=\Mpot_F(\emptyset)\neq0$.

\begin{proposition}[the environment barrier for the triangle]\label{prop:env-barrier}
Let $C$ be the directed triangle with $\eta=0$. With one frozen vertex there is no frozen potential over $\Qzbar$. With two frozen vertices exactly four graphs $U$ admit one: $U$ contains all six edges between $C$ and $F$, with or without the edge inside $F$ and with or without the three edges inside $C$. Each solution set is finite; every weight on $C$ has degree $3$ over $\Qz$ and every weight on $F$ has degree $6$. For the smallest $U$ an environment weight $a$ satisfies
\[
(z^{2}-4)\,a^{6}+12a^{4}+4z\,a^{3}-12a^{2}+8=(za^{3}+2)^{2}+4(1-a^{2})^{3}=0,
\]
an irreducible sextic over $\Qz$ that has no real root at any real $z$ with $|z|\ge4$.
\end{proposition}
\begin{proof}
The enumeration of graphs $U$ and the degrees are exact computations (saturated Gr\"obner bases over $\Qz$ and factorization, \cref{app:verification}). The sextic follows from \cref{thm:triangle-solved} by eliminating $t$, and its two forms agree by expansion. For real $z$ with $|z|\ge4$ it is positive at every real $a$: when $|a|\le1$ both terms are nonnegative and they do not vanish together, and when $|a|>1$, $|za^{3}+2|\ge4|a|^3-2>2|a|^{3}$ gives $(za^{3}+2)^{2}>4a^{6}>4(a^{2}-1)^{3}$.
\end{proof}

\begin{theorem}[the rescued triangle, solved]\label{thm:triangle-solved}
Let $C$ be the directed triangle with $\eta=0$, $F=\{f_1,f_2\}$, and $U$ the six edges between $C$ and $F$. The frozen potentials are exactly the following: the three cycle weights equal $\alpha=z+t$, where $t$ is a root of $t^{3}-3t=z$, and the two environment weights are the roots of $(t^{2}-1)x^{2}-2tx+2=0$. With $z=v^{3}+v^{-3}$ this reads $t=v+v^{-1}$, $\alpha=\mu(P_3)(t)$ and $x=(t\pm\sqrt{2-t^{2}})/(t^{2}-1)$. For real $z\notin\{0,\pm1,\pm\sqrt2\}$ a real frozen potential exists if and only if $|z|<2$, and one branch has weights that are real power series in $z$ at $z=0$.
\end{theorem}
\begin{proof}
Write $\Mpot(S)=\mu_a(U[S\cup F])$ for $S\subseteq C$, $p=a_{f_1}a_{f_2}$ and $s=a_{f_1}+a_{f_2}$. The frozen identities say $\Mpot(S)=\lambda\mu(P_{|S|})(z)$ on every $S\subseteq C$ (every subset of the triangle is a cyclic interval), with $\lambda=\Mpot(\emptyset)=p$. Since $\Mpot(\{j\})=a_jp-s$, the singletons give $a_jp-s=zp$, so every cycle weight equals $\alpha=z+t$ with $t=s/p$. With equal cycle weights, $\Mpot(C\setminus\{i\})=\alpha^{2}p-2\alpha s+2$ and $\Mpot(C)=\alpha^{3}p-3\alpha^{2}s+6\alpha$. Dividing by $p$ and substituting $\alpha=z+t$, the identity for the pairs becomes $2/p=t^{2}-1$, and the identity for $C$ becomes $(z+t)^{2}(z-2t)+3(z+t)(t^{2}-1)=z^{3}-2z$, that is, $t^{3}-3t=z$. So $p=2/(t^{2}-1)$ and $s=tp$, and the environment weights are the roots of $(t^{2}-1)x^{2}-2tx+2$. Conversely these values satisfy every identity, and $\Mpot(W)=p\,\mu(P_{|W|})(z)\neq0$.

For real $|z|>2$ the cubic has one real root, with $|t|>2$, so $2-t^{2}<0$ and $x$ is not real; at the other two roots $\alpha=z+t$ is not real. At $z=\pm2$ the roots are $t=\pm2$, again with $2-t^2<0$, and the double root $t=\mp1$, where the quadratic degenerates. For $|z|<2$ write $z=2\cos3\theta$ with $0<\theta<\pi/3$; the roots are $t=2\cos(\theta+2\pi k/3)$, one of which has $|\cos|<\frac12$, so $t^{2}<1$ and all weights are real, except where $t^2=1$ or $p$ or a denominator vanishes, which happens only at $z\in\{0,\pm1,\pm\sqrt2\}$. Near $z=0$ the root $t=-z/3+O(z^{3})$ is a real power series, and so are $\alpha$ and $x$, with $x(0)=\mp\sqrt2$.
\end{proof}

\begin{theorem}[every directed cycle is rescued by a frozen environment]\label{thm:every-cycle}
Let $m\ge3$, let $C$ be the directed $m$-cycle with $\eta=0$, let $|F|=m-1$, and let $U$ consist of all edges between $C$ and $F$. Let $t$ be a root of $2T_m(t/2)=z$, where $T_m$ is the Chebyshev polynomial of the first kind, let $\alpha=U_m(t/2)=\mu(P_m)(t)$, and let
\[
Q(x)=\sum_{r=0}^{m-1}F_r(\alpha)\,\frac{x^{r}}{r!},\qquad F_r(a)=\sum_{k=0}^{r}\binom{r}{k}(-a)^{r-k}\,\mu(P_k)(z).
\]
Then the weight $\alpha$ on every cycle vertex and the $m-1$ roots of $Q$ on the environment form a frozen potential. Moreover:
\begin{enumerate}[label=\textup{(\arabic*)}]
\item With $z=v^{m}+v^{-m}$ and $t=v+v^{-1}$, $F_r(\alpha)=\mathsf D^{r}(v^{2m}-v^{2r})/(v^{2m}-1)$ with $\mathsf D=v^{m}-\alpha$, so the environment weights are $x=y/\mathsf D$ where $v^{2m}\exp_{m-1}(y)=\exp_{m-1}(v^{2}y)$ and $\exp_n(y)=\sum_{j\le n}y^{j}/j!$ is the truncated exponential.
\item Near $z=\infty$, $\alpha-z=U_{m-2}(t/2)$ grows like $z^{(m-2)/m}$, and each environment weight is $-y_0v^{-(m-2)}(1+o(1))$, of order $z^{-(m-2)/m}$, where $y_0$ runs over the zeros of $\exp_{m-1}$.
\item The $m$ branches are the sheets of the Chebyshev cover $t\mapsto2T_m(t/2)$, branched over $\pm2$ and $\infty$; its monodromy group is dihedral of order $2m$, the loop around $z=\infty$ permuting the branches cyclically and the loops around $\pm2$ acting as reflections.
\item At every sufficiently large real $z$, no branch has all weights real.
\end{enumerate}
\end{theorem}
\begin{proof}
$U$ is invariant under all permutations of $C$ and the cycle weights are equal, so $\Mpot(W)$ depends only on $k=|W|$; write $\Mpot_k$. Sorting matchings by the cycle vertices matched into $F$ gives the exponential generating function $\sum_k\Mpot_kx^{k}/k!=e^{\alpha x}\prod_j(\beta_j-x)$, with $\beta_j$ the environment weights. The frozen identities say $\Mpot_k=\lambda\mu(P_k)(z)$ for $0\le k\le m$ with $\lambda\neq0$. Multiplying by $e^{-\alpha x}$ and comparing coefficients up to $x^m$, this holds exactly when $\prod_j(\beta_j-x)=\lambda Q(x)$ and $F_m(\alpha)=0$. With $\mu(P_k)(z)=(w^{k+1}-w^{-k-1})/(w-w^{-1})$ and $z=w+w^{-1}$,
\[
F_r(a)=\frac{w\,(w-a)^{r}-w^{-1}(w^{-1}-a)^{r}}{w-w^{-1}}.
\]
Put $w=v^{m}$ and $\alpha=\sum_{j=0}^{m}v^{m-2j}=U_m(t/2)$. Then $w^{-1}-\alpha=v^{2}(w-\alpha)$, which gives the formula in (1), and it vanishes at $r=m$. Since $F_m$ has degree $m$ in $a$, its roots are exactly the $m$ values $U_m(t/2)$. Finally $\lambda=\prod_j\beta_j\neq0$, because $Q(0)=1$ and the leading coefficient $F_{m-1}(\alpha)/(m-1)!$ is nonzero by (1).

For (2), $U_m-U_{m-2}=2T_m$ gives $\alpha-z=U_{m-2}(t/2)$, and $\mathsf D=-v^{m-2}(1+O(v^{-2}))$. Dividing the equation of (1) by $v^{2m}$ gives $\exp_{m-1}(y)=\sum_{r<m}v^{2r-2m}y^{r}/r!$, whose right side tends to $0$ uniformly on compact sets, since $|v|\to\infty$; by Hurwitz's theorem the $m-1$ roots tend to the zeros of $\exp_{m-1}$. Item (3) is the classical monodromy of Chebyshev polynomials. For (4), suppose that along real $z\to\infty$ one branch had all environment weights real. Then all $y_j=\mathsf D\beta_j$ lie on the line $\RR\mathsf D$, so in the limit all zeros of $\exp_{m-1}$ lie on one line through $0$, closed under conjugation. It is not $\RR$, because $\exp_n$ has at most one real zero ($\exp_n>0$ for even $n$, and $\exp_n$ is increasing for odd $n$), so it is $i\RR$; but the zeros of $\exp_{m-1}$ sum to $-(m-1)\neq0$.
\end{proof}

\begin{remark*}[inside the band]
A scan of the family of \cref{thm:every-cycle} over real $z\in(-2,2)$ finds branches with all weights real for $m=3$ on the whole band (\cref{thm:triangle-solved}), for $m=4$ on about $(0.178,2)$ and for $m=5$ on about $(-1.024,1.025)$; it finds none for $m=6,7$ on a grid of $799$ points. These are numerical observations, not theorems (\cref{prob:band}).
\end{remark*}

\subsection{Genuine digraphs realize complex environments}\label{ssec:counterexample}

\begin{lemma}[disjoint unions]\label{lem:disjoint-union}
If $D_1$ and $D_2$ have potentials $(U_1,a_1)$ and $(U_2,a_2)$ over a field $\mathbb K$, then $(U_1\sqcup U_2,a_1\sqcup a_2)$ is a potential for $D_1\sqcup D_2$ over $\mathbb K$.
\end{lemma}
\begin{proof}
A reachable position of $D_1\sqcup D_2$ with the token in $D_1$ has the form $(W_1\cup V(D_2),v)$ with $(W_1,v)$ reachable in $D_1$, and its resolvent is $R_{D_1}(W_1,v)$. The weighted matching polynomial factorizes: $\mu_a(U[W_1\cup V(D_2)])=\mu_{a_1}(U_1[W_1])\,\mu_{a_2}(U_2[V(D_2)])$, where the second factor is nonzero because $(V(D_2),s)$ is reachable in $D_2$ for every $s$. So the identities of $D_1$ carry over, and symmetrically for $D_2$.
\end{proof}

For $m\ge3$ let $\exp_m(w)=\sum_{k=0}^{m}w^{k}/k!$ be the truncated exponential, and let
\begin{equation}\label{eq:Ym}
Y_m(y)=y^{m+1}-\frac{m!\,z}{\mu(P_m)(z)}\sum_{\ell=0}^{m}\mu(P_\ell)(z)\,\frac{y^{\ell}}{\ell!}\ \in\Qz[y].
\end{equation}
For instance $\mu(P_3)\,Y_3/z=(z^{2}-2)y^{4}-z(z^{2}-2)y^{3}-3(z^{2}-1)y^{2}-6zy-6$.

\begin{theorem}[a counterexample over $\Qzbar$]\label{thm:complex-counterexample}
Let $m\ge3$ and let $D_m$ be the disjoint union of the directed cycle $c_0\to c_1\to\dots\to c_{m-1}\to c_0$ and a set $F$ of $m+1$ isolated vertices. Let $U$ be the complete bipartite graph between the cycle and $F$, put $a_{c}=0$ on the cycle, and let $(a_f)_{f\in F}$ be the $m+1$ roots of $Y_m$ in $\Qzbar$, in any order.
\begin{enumerate}[label=\textup{(\arabic*)}]
\item $(U,a)$ is a local weighted matching potential for $D_m$ over the splitting field of $Y_m$. In particular $D_m$, which is not SCC-symmetric, has a potential over $\Qzbar$.
\item Conversely, if $(U,a)$ is a potential for $D_m$ with this $U$ and $a\equiv0$ on the cycle, then $(a_f)_{f\in F}$ are the roots of $Y_m$.
\end{enumerate}
\end{theorem}

\begin{proof}
\emph{Positions.} Every vertex is a start. From a cycle vertex the token runs around the cycle and never meets $F$, so the reachable positions are $(I\cup F,v)$, with $I$ a cyclic interval of length $\ell\in\{1,\dots,m\}$ that begins at $v$, and the starts $(V,f)$, $f\in F$, which have no option. The game from $(I\cup F,v)$ is DVG on a directed path with $\ell$ vertices, so $R(I\cup F,v)=\mu(P_{\ell-1})/\mu(P_\ell)$, and $R(V,f)=1/z$.

\emph{Matching polynomials.} Let $J$ be a set of $\ell$ cycle vertices. A matching of $U[J\cup F]$ that leaves a vertex of $J$ uncovered contributes the factor $a_c=0$, so only matchings that cover $J$ count. Such a matching is a choice of an $\ell$-subset $S\subseteq F$ and a bijection $J\to S$, and it contributes $(-1)^{\ell}\prod_{f\in F\setminus S}a_f$. Hence
\[
\mu_a(U[J\cup F])=(-1)^{\ell}\,\ell!\;e_{m+1-\ell}(a_F),\qquad \mu_a(U[V-f])=(-1)^{m}\,m!\quad(f\in F),
\]
where $e_j(a_F)$ is the $j$-th elementary symmetric function of the environment weights; in the second formula every vertex of $F-f$ is matched, since $|F-f|=m$.

\emph{Vieta.} Put $c=(-1)^{m}m!\,z/\mu(P_m)$. Comparing \eqref{eq:Ym} with
\[
Y_m(y)=\sum_{j=0}^{m+1}(-1)^{j}e_j(a_F)\,y^{m+1-j}
\]
at the coefficient of $y^{\ell}$ gives
\[
e_{m+1-\ell}(a_F)=(-1)^{\ell}\,c\,\frac{\mu(P_\ell)}{\ell!}\qquad(0\le\ell\le m),
\]
so that $\mu_a(U[J\cup F])=c\,\mu(P_{|J|})$ for every set $J$ of cycle vertices.

\emph{Identities.} At $(I\cup F,v)$ with $|I|=\ell$: $\mu_a(U[I\cup F])=c\,\mu(P_\ell)\neq0$, and $\mu_a(U[(I-v)\cup F])=c\,\mu(P_{\ell-1})=R(I\cup F,v)\,\mu_a(U[I\cup F])$. At $(V,f)$: $\mu_a(U[V])=c\,\mu(P_m)=(-1)^{m}m!\,z\neq0$, and $\mu_a(U[V-f])=(-1)^{m}m!=\frac1z\,\mu_a(U[V])$. This proves (1): the arcs of the cycle are one-way arcs inside a strong component.

For (2), the identities at the positions $(I\cup F,v)$ with $|I|=\ell$ say $\mu_a(U[I\cup F])/\mu(P_\ell)$ is independent of $\ell$, and nonzero; call it $c'$. The identity at $(V,f)$ then says $(-1)^{m}m!=c'\mu(P_m)/z$, so $c'=c$, and the matching formula determines $e_1(a_F),\dots,e_{m+1}(a_F)$ as above. A monic polynomial is determined by its elementary symmetric functions.
\end{proof}

\begin{corollary}[the complex form of Theorem 6.10 is false]\label{cor:complex-false}
Let $\mathbb K$ be an algebraically closed field containing $\Qz$, for instance $\Qzbar$ or the algebraic closure of $\CC(z)$. The digraphs with a local weighted matching potential over $\mathbb K$ strictly contain the SCC-symmetric digraphs. The smallest counterexample produced here is $D_3$, with seven vertices; for every $m\ge3$ and every digraph $D'$ with a potential over $\mathbb K$ (for instance any SCC-symmetric digraph), $D_m\sqcup D'$ is a counterexample.
\end{corollary}
\begin{proof}
\Cref{thm:complex-counterexample} and \cref{lem:disjoint-union}.
\end{proof}

\Cref{thm:complex-counterexample} was also verified by computer, independently of the proof: \code{complex_counterexample.verify_symbolic} enumerates every reachable position with the package's game recursion, computes each $\mu_a(U[S])$ symbolically in indeterminate environment weights, rewrites it in elementary symmetric functions and checks every identity exactly, for $3\le m\le5$; \code{verify_numeric} checks the identities with the roots of $Y_m$ at real and complex $z$ to $50$ digits (residuals below $10^{-50}$), for $m\le6$.

The counterexample is not confined to a bare cycle: exits do not destroy it.

\begin{proposition}[cycles with exits]\label{prop:sinks}
Let $m\ge3$ and $r\ge0$, and let $D_{m,r}$ be $D_m$ with $r$ private sinks attached to every cycle vertex \textup{(}arcs $c_i\to s_{i,1},\dots,s_{i,r}$\textup{)}. Put $x=z-r/z$ and let $Y_{m,r}$ be \eqref{eq:Ym} with $\mu(P_\ell)(z)$ replaced by $\mu(P_\ell)(x)$ throughout \textup{(}the factor $m!\,z$ is unchanged\textup{)}. Then $U$ as in \cref{thm:complex-counterexample}, weight $0$ on the cycle, weight $z$ on every sink \textup{(}isolated in $U$\textup{)} and the roots of $Y_{m,r}$ on $F$ form a potential for $D_{m,r}$ over $\Qzbar$.
\end{proposition}
\begin{proof}
Let $S$ be the set of sinks. While the token is on the cycle every sink is unvisited, and a move to a sink ends the play there; so the reachable positions are $(I\cup F\cup S,v)$ with the cycle intervals $I$ as before, the positions $((I-v)\cup F\cup S,s)$ with $s$ a sink of $v$, and the starts $(V,f)$ and $(V,s)$. Every sink has fresh-entry resolvent $1/z$, so each cycle vertex carries the self-energy $r/z$ and $R(I\cup F\cup S,v)=\mu(P_{\ell-1})(x)/\mu(P_\ell)(x)$; at a sink $R=1/z$. Since the sinks are isolated in $U$ with weight $z$, $\mu_a(U[T\cup S'])=z^{|S'|}\mu_a(U[T])$ for $T\cap S=\emptyset$ and $S'\subseteq S$, so every identity at a sink reads $1/z=1/z$, and the other identities are those of the proof of \cref{thm:complex-counterexample} with $\mu(P_\ell)(z)$ replaced by $\mu(P_\ell)(x)$ and $c=(-1)^mm!\,z/\mu(P_m)(x)$. Nonvanishing holds because $\mu(P_\ell)(x)$ is a nonzero rational function.
\end{proof}

\begin{proposition}[arithmetic of the environment weights]\label{prop:Ym-arithmetic}
Let $m\ge3$.
\begin{enumerate}[label=\textup{(\arabic*)}]
\item $Y_m$ is unramified over $z=\infty$: its roots are $m+1$ distinct Laurent series in $z^{-1}$. Exactly $m$ of them are $w_i/z+O(z^{-3})$, where $w_1,\dots,w_m$ are the \textup{(}distinct\textup{)} zeros of $\exp_m$, and the remaining one is $z+m/z+O(z^{-3})$.
\item At every sufficiently large real $z$ the weights of \cref{thm:complex-counterexample} include exactly $m-1$ non-real ones if $m$ is odd and exactly $m$ if $m$ is even; the other weights, including the $m$ zeros on the cycle, are real.
\item For $3\le m\le6$ the Galois group of $Y_m$ over $\CC(z)$, which is the monodromy group of the $(m+1)$-sheeted cover $\{Y_m=0\}\to\mathbb P^{1}_z$, is the full symmetric group $S_{m+1}$. For $m=3$ the discriminant of $\mu(P_3)Y_3/z$ in $y$ is $-108\,(z^{2}-2)^{2}\,(2z^{8}+26z^{6}+99z^{4}-10z^{2}-676)$.
\end{enumerate}
\end{proposition}
\begin{proof}
(1) Put $u=1/z$ and $y=wu$. Since $\mu(P_\ell)(z)=z^{\ell}g_\ell(u^{2})$ with $g_\ell\in\ZZ[u^{2}]$ and $g_\ell(0)=1$,
\[
-\frac{z^{m-1}}{m!}\,Y_m(wu)=\frac{1}{g_m(u^{2})}\sum_{\ell=0}^{m}g_\ell(u^{2})\frac{w^{\ell}}{\ell!}-\frac{u^{2}w^{m+1}}{m!},
\]
a polynomial in $w$ with coefficients in $\QQ[[u^{2}]]$ that reduces to $\exp_m(w)$ at $u=0$. The zeros of $\exp_m$ are simple, since $\exp_m-\exp_m'=w^{m}/m!$ does not vanish at a zero of $\exp_m$ other than $0$, and $\exp_m(0)=1$. By Hensel's lemma each zero $w_i$ lifts to a root $w_i(u)=w_i+O(u^{2})$ in $\QQ(w_i)[[u^{2}]]$, convergent for small $|u|$, which gives $m$ roots $y=w_i(u)u$. The last root is $e_1(a_F)$ minus their sum; $e_1(a_F)=z$ by the Vieta formula at $\ell=m$, and $\sum_iw_i=-m$, which gives $z+m/z+O(z^{-3})$.

(2) For real $z$ the polynomial $Y_m$ has real coefficients, so its non-real roots come in conjugate pairs. The polynomial $\exp_m$ has no real zero for even $m$ ($\exp_m>0$ on $\RR$) and exactly one for odd $m$ ($\exp_m'=\exp_{m-1}>0$). A root $w_i(u)u$ with $\Im w_i\neq0$ has imaginary part $u\,\Im w_i+O(u^{3})\neq0$ for small real $u$. A root near a real $w_i$, or the large root, is the only root in a small disc symmetric about the real axis, so it equals its conjugate.

(3) Computer-assisted (\code{scripts/ym_arithmetic.py}). For $3\le m\le6$ the primitive part of $\mu(P_m)Y_m$ (which is $\mu(P_m)Y_m/z$ for odd $m$) is irreducible in $\QQ[z,y]$, so $Y_m$ is irreducible over $\Qz$, and its specialization at $z=3$ has Galois group $S_{m+1}$ over $\QQ$ (PARI \code{polgalois}); since specialization at a place where the polynomial stays separable of the same degree embeds the Galois group of the specialization into the Galois group over $\Qz$, the latter is $S_{m+1}$. The discriminant has an irreducible factor of multiplicity one (of degree $8,20,26,42$ for $m=3,4,5,6$) coprime to the leading coefficient; over each of its roots $z_0$ the discriminant, which up to a unit is $\prod_{i<j}(y_i-y_j)^2$, has a simple zero. Three coinciding roots, or two coinciding pairs, would make it vanish to order at least $2$; so exactly two roots meet at $z_0$, and they cannot both be single-valued near $z_0$, since then $(y_1-y_2)^2$ would vanish to even order. A small loop around $z_0$ therefore exchanges them and fixes the other roots: the local monodromy is a transposition. The Galois group over $\CC(z)$ is a normal subgroup of the Galois group over $\Qz$ (the constant field of the splitting field is Galois over $\QQ$), and the only normal subgroup of $S_{m+1}$ that contains a transposition is $S_{m+1}$.
\end{proof}

\begin{proposition}[a triangle with isolated vertices]\label{prop:c3-minimal}
The disjoint union of the directed triangle and $k$ isolated vertices has a potential over $\Qzbar$ if and only if $k\ge4$.
\end{proposition}
\begin{proof}
For $k\ge4$ combine \cref{thm:complex-counterexample} ($m=3$) with \cref{lem:disjoint-union}, an isolated vertex having the potential with $U=\emptyset$ and weight $z$. For $k\le2$ the digraph has at most five vertices (\cref{thm:exact5}). For $k=3$: by \cref{thm:locality} the triangle lies in one component of $U$; up to rotations of the triangle and permutations of the isolated vertices there are $2{,}208$ graphs $U$ on six vertices, $1{,}928$ of them with this property, and for each of these the saturated ideal of \cref{thm:exact5} over $\Qz$ is the unit ideal (\code{scripts/search6_c3k3.py}, $878$\,s with Singular).
\end{proof}

\begin{remark}[what the counterexample says about proof strategies]\label{rem:monodromy}
\begin{enumerate}[label=(\arabic*)]
\item \emph{Monodromy is not the obstruction.} The frozen rescues of \cref{thm:every-cycle} are ramified over $z=\infty$, where the Chebyshev cover permutes the $m$ branches cyclically, and their monodromy group is dihedral of order $2m$. The genuine potentials of \cref{thm:complex-counterexample} are unramified over $\infty$ (\cref{prop:Ym-arithmetic}(1)) and, for $m\le6$, have the full symmetric group $S_{m+1}$ as monodromy group. So no argument that obstructs a genuine realization through dihedral monodromy, or through ramification at $\infty$, can be correct: it would also exclude $D_m$.
\item \emph{Reality is the obstruction.} For a one-way arc of the cycle of $D_m$, every subset $X$ of the cycle that enters the proof of \cref{thm:main-real} is a path set when $m=3$, so the polynomial $G_z(v)=\sum_X\gamma_Xv^{X}$ of \S\ref{sec:realstab} is determined by the path data alone; it is a real polynomial, the same for every potential of $D_3$, and the proof shows that it has a negative Rayleigh difference at large real $z$. What fails for $D_3$ is therefore only \textup{(P2)}: the Heilmann--Lieb theorem needs real vertex weights, and the environment weights of order $z^{-1}$ that leave the real axis (\cref{prop:Ym-arithmetic}(2)) violate that hypothesis.
\item \emph{Vieta realizability.} With vanishing cycle weights every identity fixes one elementary symmetric function of the environment weights, and the $m+1$ identities at the starts $(V,f)$ collapse to one because $\mu_a(U[V-f])$ does not depend on $f$. Over an algebraically closed field any prescribed values of $e_1,\dots,e_{m+1}$ are realized; over $\RR$ only those of real-rooted polynomials. Vanishing cycle weights force $|F|\ge m+1$, since $U[V-f]$ must still have a matching that covers the cycle when $U$ has no edge inside the cycle; this is why $D_m$ has two more environment vertices than the frozen rescue of \cref{thm:every-cycle}.
\end{enumerate}
\end{remark}

\subsection{Smallest frozen environments}

\begin{proposition}[smallest environments]\label{prop:smallest-env}
For $(m,|F|)$ equal to $(3,1)$, $(4,1)$, $(4,2)$, $(5,1)$ or $(5,2)$, no graph $U$ on $C_m\sqcup F$ admits a frozen potential for $(C_m,0)$, even over $\Qzbar$. So the environment of \cref{thm:every-cycle} is as small as possible for $m=3,4$, and a $5$-cycle needs at least three frozen vertices.
\end{proposition}
\begin{proof}
Saturated Gr\"obner bases over $\Qz$ for every $U$ up to the symmetries of the cycle and of $F$: $276$, $4{,}496$, $6{,}560$ and $216{,}288$ graphs for $(4,1)$, $(4,2)$, $(5,1)$ and $(5,2)$; the case $(3,1)$ is \cref{prop:env-barrier}.
\end{proof}

\begin{proposition}[symmetric environments]\label{prop:symmetric-env}
Let the weights lie in $\Qzbar$, $\eta=0$, and let $U$ be invariant under every permutation of $C=C_m$ that fixes $F$ pointwise; equivalently all cycle vertices have the same set $N$ of $U$-neighbours in $F$, and $U[C]$ is empty or complete. If a frozen potential exists, then $|N|\ge m-1$ when $U[C]$ is empty, for every $m\ge3$, and when $U[C]$ is complete, for $3\le m\le10$. \Cref{thm:every-cycle} attains $|F|=|N|=m-1$ with $U[C]$ empty, and the frozen potential that $D_m$ induces on its cycle \textup{(}\cref{prop:restriction}\textup{)} has $|N|=m+1$.
\end{proposition}
\begin{proof}
On a singleton interval $\{c\}$ no other cycle vertex is present, and the vertex recurrence at $c$ with the frozen identity $\Mpot(\{c\})=z\lambda$ gives every cycle vertex the weight $\alpha=z+\sum_{f\in N}\rho(F,f)$, where $\rho(F,f)=\mu_a(U[F-f])/\mu_a(U[F])$. Let $\kappa=1$ if $U[C]$ is complete and $\kappa=0$ otherwise. Sorting matchings as in the proof of \cref{thm:every-cycle}, with the factor $e^{-x^2/2}$ counting matchings inside a complete graph, gives $\sum_k\Mpot_kx^{k}/k!=e^{\alpha x-\kappa x^{2}/2}Q(x)$, where $Q$ is a polynomial of degree at most $|N|$. Write $z=w+w^{-1}$; then $w$ is transcendental over $\QQ$. The frozen identities $\Mpot_k=\lambda\mu(P_k)(z)$ for $k\le m$ give
\[
Q(x)\equiv\lambda\,e^{\kappa x^{2}/2}\,\frac{w\,e^{Ax}-w^{-1}e^{Bx}}{w-w^{-1}}\pmod{x^{m+1}},\qquad A=w-\alpha,\quad B=w^{-1}-\alpha .
\]
If $|N|\le m-2$, the coefficients of $x^{m-1}$ and $x^{m}$ on the right vanish. For $\kappa=0$ this says $wA^{m-1}=w^{-1}B^{m-1}$ and $wA^{m}=w^{-1}B^{m}$. If $A=0$ then $B=0$ and $w-w^{-1}=A-B=0$; otherwise dividing gives $A=B$, hence again $w^{2}=1$, impossible for transcendental $w$. For $\kappa=1$ the two coefficients, multiplied by $(m-1)!/\lambda$ and $m!/\lambda$, are polynomials in $\alpha$ with leading coefficients $\pm1$, and for each $m\le10$ their resultant in $\alpha$ is a nonzero Laurent polynomial in $w$ (exact computation), which does not vanish at a transcendental $w$.
\end{proof}

\begin{proposition}[the $5$-cycle and the $6$-cycle]\label{prop:c5c6}
Over $\Qzbar$, with $\eta=0$ and $U[C]=\emptyset$, the directed $5$-cycle has no frozen potential with $|F|=3$, and the directed $6$-cycle has none with $|F|\le3$. Up to rotation of $C$ and permutation of $F$ there are $9{,}168$ graphs $U$ for the $5$-cycle, and $14$, $724$ and $59{,}976$ for the $6$-cycle with $|F|=1,2,3$; all give the unit ideal after saturation.
\end{proposition}
\begin{proof}
Saturated Gr\"obner bases over $\Qz$, as for \cref{prop:smallest-env}.
\end{proof}

\begin{conjecture}[minimal frozen environments]\label{conj:minimal-env}
Over $\Qzbar$, every frozen potential of the directed $m$-cycle with $\eta=0$ has at least $m-1$ frozen vertices.
\end{conjecture}

The conjecture holds for $m=3,4$ (\cref{prop:smallest-env}, with \cref{thm:exact5} for $F=\emptyset$), for every $m$ when $U$ is symmetric with $U[C]$ empty, for $m\le10$ when $U$ is symmetric with $U[C]$ complete (\cref{prop:symmetric-env}), for $m=5$ when $U[C]=\emptyset$ and in general up to two frozen vertices, and for $m=6$ with $U[C]=\emptyset$ up to three frozen vertices. (With $F=\emptyset$ and $U[C]=\emptyset$ the frozen identities on intervals of lengths $1$ and $2$ contradict each other.) For genuine digraphs the relevant count is larger: $D_3$ needs four isolated vertices (\cref{prop:c3-minimal}).

\subsection{Real environments, pointwise in \texorpdfstring{$z$}{z}}
For weights in $\Puis$, \cref{thm:frozen-real} excludes every frozen potential of a non-symmetric strongly connected digraph. The two results below say more for directed cycles: they hold at each real $z$ outside the band $|z|\le2$ separately, not only for large $z$. Let $C$ be a chordless directed cycle $x_0\to\dots\to x_{m-1}\to x_0$, $m\ge3$, with self-energies $\eta$ vanishing at infinity, let $(U,a)$ be a frozen potential with environment $F$, and write $\Mpot(W)=\mu_a(U[W\cup F])$, $\lambda=\Mpot(\emptyset)$. The frozen identities make $\Mpot$ a set potential for $(C,\eta)$, so $\eta$ has period $2$ and all rotation continuants of $C$ have one value $\cont(C)$ (\cref{thm:set-cycles}); telescoping gives $\Mpot(I)=\lambda\cont(I)$ on every cyclic interval $I$, where $\cont(I)$ is the continuant of the weights $z-\eta$ along $I$ in cycle order (\S\ref{ssec:setpot-cycles}).

\begin{lemma}[frozen path lemma]\label{lem:frozen-path}
If the weights lie in $\Puis$, then $U$ has no edge inside $C$. When $\eta=0$ this holds pointwise: at every real $z$ with $|z|\ge2$, real weights satisfying the frozen identities at $z$ have $U[C]=\emptyset$.
\end{lemma}
\begin{proof}
Let $x\neq y$ in $C$. Since $m\ge3$, one of the two cyclic intervals with end vertices $x$ and $y$ has at least three vertices; let $I$ be that interval, running from $x$ to $y$. Then $I-x$, $I-y$ and $I-x-y$ are cyclic intervals, and \cref{lem:two-point} at $S=I\cup F$ gives
\[
\Mpot(I-x)\,\Mpot(I-y)-\Mpot(I)\,\Mpot(I-x-y)=\sum_{P}\mu_a\bigl(U[(I\cup F)\setminus V(P)]\bigr)^{2}
\]
over the $U$-paths $P$ from $x$ to $y$ in $U[I\cup F]$. The left side is $\lambda^{2}[\cont(I-x)\cont(I-y)-\cont(I)\cont(I-x-y)]=\lambda^{2}$ by the determinant identity for continuants (each transfer matrix has determinant $1$); when $I=C$ this uses the equality of the rotation continuants. If $xy\in U$, the path $P=xy$ contributes $\lambda^{2}\cont(I-x-y)^{2}$. At a real $z$ where the weights are real every term on the right is nonnegative, so $\cont(I-x-y)(z)^{2}\le1$. But $I-x-y$ has $|I|-2\ge1$ vertices, so $\cont(I-x-y)=z^{|I|-2}(1+o(1))$ exceeds $1$ in absolute value at large real $z$; when $\eta=0$ it is $\mu(P_{|I|-2})(z)$, whose absolute value is at least $|I|-1\ge2$ at every real $z$ with $|z|\ge2$.
\end{proof}

\begin{lemma}[the environment seen from the cycle]\label{lem:env-seen}
Suppose $U[C]=\emptyset$; the weights may lie in any field. For every cycle vertex $c$, with $N(c)$ its set of $U$-neighbours in $F$,
\[
a_c=z-\eta_c+\sum_{f\in N(c)}\rho(F,f),\qquad \rho(F,f)=\frac{\mu_a(U[F-f])}{\mu_a(U[F])}.
\]
For consecutive cycle vertices $c,c^{+}$, the sum of $\bigl(\mu_a(U[F\setminus V(P)])/\lambda\bigr)^{2}$ over the $U$-paths $P$ from $c$ to $c^{+}$ with inner vertices in $F$ equals $1$.
\end{lemma}
\begin{proof}
The frozen identity on $\{c\}$ gives $\Mpot(\{c\})=\lambda(z-\eta_c)$, and the vertex recurrence at $c$ gives $\Mpot(\{c\})=a_c\lambda-\sum_{f\in N(c)}\mu_a(U[F-f])$. For the second claim apply \cref{lem:two-point} to $S=\{c,c^{+}\}\cup F$: its left side is $\lambda^{2}((z-\eta_c)(z-\eta_{c^+})-\cont(c,c^{+}))=\lambda^{2}$.
\end{proof}

\begin{theorem}[no real mean-field rescue outside the band]\label{thm:meanfield}
Let $\eta=0$ and let every vertex of $C$ have the same set $N\subseteq F$ of $U$-neighbours in $F$. At no real $z$ with $|z|>2$ do real weights satisfy the frozen identities. The band edge is sharp: for $m=3$ the environments of \cref{thm:triangle-solved} are of this kind and are real at every $z$ in the band except $0,\pm1,\pm\sqrt2$. A constant self-energy $e$ shifts $z$ to $z-e$, so the same holds for $\eta\equiv e$ at every real $z$ with $|z-e|>2$.
\end{theorem}
\begin{proof}
Fix a real $z$ with $|z|>2$ and real weights satisfying the frozen identities at $z$, and write $z=w+w^{-1}$ with $w$ real, $|w|>1$. By \cref{lem:frozen-path}, $U[C]=\emptyset$, and by \cref{lem:env-seen} every cycle weight equals the real number $\alpha=z+\sum_{f\in N}\rho(F,f)$. Sorting matchings of $U[W\cup F]$ by the set of cycle vertices matched into $F$ gives $\Mpot(W)=\Mpot_k$ with $k=|W|$ and $\sum_k\Mpot_kx^{k}/k!=e^{\alpha x}Q(x)$, where
\[
Q(x)=\sum_{J\subseteq N}(-x)^{|J|}\,\mu_a\bigl(U[F\setminus J]\bigr)=\mu_{a-x\one_N}\bigl(U[F]\bigr)
\]
is the matching polynomial of $U[F]$ with $x$ subtracted from the weight of every vertex of $N$. By \cref{thm:hl} restricted to this real line, $Q$ has only real zeros, and $Q(0)=\lambda\neq0$. Writing $Q(x)=\sum_jq_jx^{j}/j!$, the identities $\Mpot_k=\lambda\mu(P_k)(z)$ for $k\le m$ give $q_k=\lambda F_k(\alpha)$ for $k\le m$, with $F_k$ as in \cref{thm:every-cycle}.

The Jensen polynomial $g_3(t)=\sum_{k=0}^{3}\binom3kq_kt^{k}$ of the real-rooted polynomial $Q$ is real-rooted \cite{PolyaSchur}: $g_3(t)=t^{3}h(1/t)$ with $h(y)=Q(d/dy)\,y^{3}$, and up to a constant $Q(d/dy)$ is a product of factors $d/dy-\varrho$ with $\varrho$ real, each of which maps a real-rooted polynomial $p$ to $p'-\varrho p=e^{\varrho y}(e^{-\varrho y}p)'$, real-rooted by Rolle's theorem; $h$ has degree $3$ because $q_0\neq0$. Since $m\ge3$, with $A=w-\alpha$ and $B=w^{-1}-\alpha$,
\[
g_3(t)=\lambda\,\frac{w(1+At)^{3}-w^{-1}(1+Bt)^{3}}{w-w^{-1}}.
\]
For a cube root $\zeta$ of $w^{-2}$ with $A-\zeta B\neq0$ put $t=(\zeta-1)/(A-\zeta B)$. Then $1+At=\zeta(A-B)/(A-\zeta B)$ and $1+Bt=(A-B)/(A-\zeta B)$, so $g_3(t)$ is a multiple of $w\zeta^{3}-w^{-1}=0$. The two non-real cube roots satisfy $A-\zeta B\neq0$, because $A,B$ are real and $A-B=w-w^{-1}\neq0$; the real M\"obius map $\zeta\mapsto(\zeta-1)/(A-\zeta B)$, of determinant $A-B\neq0$, sends them to two non-real roots of $g_3$, a contradiction.
\end{proof}

\subsection{What remains open over \texorpdfstring{$\Qzbar$}{the algebraic closure}}
Every counterexample of this section has incomparable vertices, as it must for real weights by \cref{cor:unilateral}; over $\Qzbar$ that corollary is not available, because it rests on \cref{prop:path-lemma}. The natural complex question is therefore the following.

\begin{problem}\label{prob:complex}
Characterize the digraphs that have a local weighted matching potential over $\Qzbar$. In particular:
\begin{enumerate}[label=\textup{(\alph*)}]
\item Does some strongly connected \textup{(}or unilateral\textup{)} digraph that is not symmetric have one?
\item Is there a counterexample on six vertices? None exists on five (\cref{thm:exact5}), and the triangle with three isolated vertices is not one (\cref{prop:c3-minimal}).
\item Does $C_m\sqcup\overline{K}_k$ have a potential for some $k\le m$ when $m\ge4$?
\item Is the Galois group of $Y_m$ over $\CC(z)$ equal to $S_{m+1}$ for every $m\ge3$?
\end{enumerate}
\end{problem}

\section{One-way arcs inside a strong component}\label{sec:oneway}

\Cref{thm:conj610} says where potentials exist; it says nothing about complexity, since directed cycles have no potential and trivial outcomes. This section asks how far matching data decide DVG when strong components carry one-way arcs. The answer depends on the tails of the arcs: arcs with a common tail are decided by a closed matching rule, with any number of arcs (\cref{thm:common-tail}); for every other geometry of two arcs, matching data at the key vertices do not suffice (\cref{prop:barrier}); and with $p$ tails a play crosses at most $p$ times (\cref{lem:p-tails}).

\subsection{Reduction to one component}\label{ssec:oneway-reduction}
By \cref{lem:entry} and the elementary proof of \cref{cor:dvg-poly}, a strong component $C$ is played from a fresh entry vertex. Exits to N-positions are losing moves and are deleted; exits to P-positions win at once and become one pendant leaf at their tail. So the game inside $C$ is $\game(H;\arcs)$: $H$ is the graph of the symmetric pairs of $C$ plus these pendant leaves, $\arcs$ is the set of one-way arcs of $C$ (each $u\to w$ with $u,w$ non-adjacent in $H$), and the token moves along edges of $H$, or along an arc of $\arcs$ whose head is unvisited. At a position with current graph $K$ (the subgraph of $H$ induced by the unvisited vertices) and token on $s$, an arc is \emph{live} if both ends are unvisited. DVG from a position $(W,v)$ is DVG on $D[W]$ started at $v$, so a rule for fresh starts decides every position. Write $\nu$ for the matching number; a vertex is \emph{essential} in a graph if every maximum matching covers it and \emph{inessential} otherwise, and by \cref{thm:matching-collapse} the player to move in undirected vertex geography wins exactly from essential vertices.

\begin{definition}\label{def:key-local}
At a position of $\game(H;\arcs)$ with current graph $K$ and token on $s$, the \emph{key points} are $s$ and the ends of the live arcs. The \emph{key-local data} are the coincidence pattern of the key points, the deficiencies $\nu(K)-\nu(K-X)$ for all sets $X$ of key points, and the adjacency in $K$ among the key points. The \emph{GE-local data} add, for every set $X$ of key points, the Gallai--Edmonds class ($\mathsf D$, $\mathsf A$ or $\mathsf C$) in $K-X$ of every key point outside $X$; here $\mathsf D(L)$ is the set of vertices missed by some maximum matching of $L$, $\mathsf A(L)$ the set of vertices outside $\mathsf D(L)$ with a neighbour in $\mathsf D(L)$, and $\mathsf C(L)$ the rest \cite{LovaszPlummer}.
\end{definition}

\subsection{Arcs with a common tail}

Let every arc of $\arcs$ leave the same vertex $u$. At a position in which $u$ is unvisited, let $W'$ be the set of live heads other than the token $s$, and let $K+uW'$ be $K$ with the edges $uw$, $w\in W'$, added. Call $K$ \emph{split} if $\nu(K+uW')=\nu(K)$, that is, if no maximum matching of $K$ misses $u$ and a live head, and \emph{joined} otherwise. Since all new edges share the vertex $u$,
\begin{equation}\label{eq:star-nu}
\nu(K+uW')=\max\Bigl(\nu(K),\ \max_{w\in W'}\nu(K-u-w)+1\Bigr).
\end{equation}

\begin{theorem}[arcs with a common tail]\label{thm:common-tail}
In the position just described the player to move wins if and only if
\begin{enumerate}[label=\textup{(\arabic*)}]
\item $s=u$, and $u$ is essential in $K+uW'$; or
\item $s\neq u$ is essential in $K$, and $K-s$ is split or $u$ is inessential in $K$; or
\item $s\neq u$ is inessential in $K$, $K$ is joined, and $\nu(K-s-u)<\nu(K)$.
\end{enumerate}
If $u$ has been visited, no arc is live and \cref{thm:matching-collapse} applies.
\end{theorem}

Two responder strategies drive the proof. Let the token be on $x$ in the current graph $K$, with the player to move called the mover.

\begin{enumerate}[label=(T\arabic*$'$),start=2]
\item If a maximum matching $M$ of $K$ misses both $x\neq u$ and $u$, the other player wins by answering each move to $t$ with $M(t)$. The mover stands only on $x$ and on vertices covered by $M$, never on $u$, so never uses an arc; and a move to a vertex missed by $M$ would close an $M$-augmenting path starting at $x$.
\item If $K$ is split and a maximum matching $M$ of $K$ misses $x$, the same answers win. $M$ is maximum in $K+uW'$, and every move, including a crossing $u\to w$, is an edge of $K+uW'$, so the augmenting-path argument applies. A split graph stays split when heads are visited.
\end{enumerate}

\begin{proof}[Proof of \cref{thm:common-tail}]
Induction on $|V(K)|$.

\emph{Item 1.} Every move from $u$ kills all arcs. Each option, a neighbour of $u$ in $K$ or a live head, is followed by undirected geography on $K-u=(K+uW')-u$, so \cref{thm:matching-collapse} for $K+uW'$ at $u$ gives the claim.

\emph{Item 2: $s$ essential in $K$.} If $K-s$ is split, move to a neighbour $t$ inessential in $K-s$ (it exists by \cref{thm:matching-collapse}). For $t=u$, a maximum matching of $K-s$ missing $u$ is maximum in $(K-s)+uW'$, so $u$ is inessential there and item 1 makes the opponent lose; otherwise (T3$'$) applies to $K-s$ with $x=t$. If $u$ is inessential in $K$, then $s$ is essential in $K-u$, and \cref{thm:matching-collapse} gives a neighbour $t\neq u$ of $s$ inessential in $K-s-u$; since $\nu(K-s-u)=\nu(K-s)$, some maximum matching of $K-s$ misses both $t$ and $u$, and after the move to $t$, (T2$'$) applies.

Otherwise $K-s$ is joined and $u$ is essential in $K$, and every move loses. Call a live head $w$ \emph{joinable} if some maximum matching $N$ of $K-s$ misses $u$ and $w$; then $s$ is not adjacent to $w$, since $N+sw$ would be a maximum matching of $K$ missing $u$. A move to $u$ loses by item 1, because $\nu((K-s)+uW')=\nu(K)>\nu(K-s-u)$. Let $t$ be any other neighbour and $K'=K-s$. Since $t$ is not a joinable head, $K'$ stays joined after $t$ is visited. If $t$ is inessential in $K'$, the move loses by item 3 for $K'$: no maximum matching of $K'$ misses $t$ and $u$, since adding $st$ would give a maximum matching of $K$ missing $u$. If $t$ is essential in $K'$, the move loses by item 2 for $K'$, because $u$ is inessential in the joined graph $K'$.

\emph{Item 3: $s$ inessential in $K$.} If $K$ is split, (T3$'$) applies with a maximum matching missing $s$, and the mover loses. If some maximum matching misses $s$ and $u$, (T2$'$) applies and the mover loses. Otherwise $K$ is joined, so some maximum matching misses $u$ and a live head $w$, and $\nu(K-s-u)=\nu(K)-1$. Then $\nu(K-s-u-w)=\nu(K)-1=\nu(K-u-w)-1$, so $s$ is essential in $L=K-u-w$, and \cref{thm:matching-collapse} gives a neighbour $t\notin\{u,w\}$ of $s$ inessential in $L-s$. Move to $t$. In $K'=K-s$ the vertex $t$ is essential, because $s$ is inessential in $K$; the graph $K'-t$ is joined, because a maximum matching of $L-s-t$ has $\nu(K)-1=\nu(K'-t)$ edges and misses $u$ and $w$; and $u$ is essential in $K'$, because $\nu(K'-u)=\nu(K)-1<\nu(K')$. By item 2 for $K'$ the opponent loses.
\end{proof}

With a single arc $u\to w$ the rule reads as follows; this settles the case $k=1$ of Open Problem~7 of the earlier draft.

\begin{corollary}[one one-way arc]\label{cor:one-arc}
Let $u,w$ be non-adjacent vertices of a graph $H$ and $s\in V(H)$. In $\game(H;u\to w)$ started at $s$ the player to move wins if and only if: $s=w$ and $s$ is essential in $H$; or $s=u$ and $u$ is essential in $H+uw$; or $s\notin\{u,w\}$ is essential in $H$ and $\nu(H-s-u-w)<\nu(H-s)$ or $u$ is inessential in $H$; or $s\notin\{u,w\}$ is inessential in $H$, $\nu(H-u-w)=\nu(H)$ and $\nu(H-s-u)<\nu(H)$.
\end{corollary}
\begin{proof}
For $s=w$ the set $W'$ is empty, so $K$ and $K-s$ are split and items 2--3 reduce to ``$s$ essential''. For $s\neq w$, $W'=\{w\}$, and by \eqref{eq:star-nu} a graph $L\ni u,w$ is split exactly when $\nu(L-u-w)<\nu(L)$.
\end{proof}

\begin{corollary}[common tails are polynomial]\label{cor:star-dvg}
Directed vertex geography is decidable in polynomial time on digraphs in which, inside every strong component, all one-way arcs share a tail, with any number of arcs per component. Processing the components in reverse topological order, every position is decided by at most six maximum-matching computations, and the proof supplies optimal moves.
\end{corollary}
\begin{proof}
By \S\ref{ssec:oneway-reduction} and \cref{thm:common-tail}; the rule uses $\nu(K)$, $\nu(K-s)$, $\nu(K-u)$, $\nu(K-s-u)$ and $\nu(K'+uW')$ for $K'\in\{K,K-s\}$, the last two through \eqref{eq:star-nu}. The package implements the rule as \code{star_tail.star_rule} and checks it against exhaustive search (\cref{app:verification}).
\end{proof}

\subsection{Two arcs: the barrier}

\begin{proposition}[two or more arcs: what survives]\label{prop:two-arcs}
Let $K^{+}$ be the current graph plus the edges $u_iw_i$ of all live one-way arcs. If $s$ is inessential in $K$ and $\nu(K^{+})=\nu(K)$, the mover loses; if $s$ is essential in $K$ and $\nu((K-s)^{+})=\nu(K-s)$, the mover wins; if a maximum matching of $K$ misses $s$ and every live tail, the mover loses. Outside these regimes the deficiencies of key-point deletions do not decide: on the bipartite graph $K$ with vertices $0,\dots,9$ and edges $0\text{--}8$, $0\text{--}9$, $1\text{--}8$, $1\text{--}9$, $2\text{--}6$, $2\text{--}7$, $2\text{--}8$, $3\text{--}7$, $3\text{--}8$, $4\text{--}6$, $4\text{--}8$, $5\text{--}6$, $5\text{--}7$, $5\text{--}9$ and the arcs $1\to4$, $2\to0$, the numbers $\nu(K-X)$ agree for the starts $s=5$ and $s=3$ over all $31$ nonempty $X\subseteq\{s,1,4,2,0\}$, yet the first player wins from $5$ and loses from $3$.
\end{proposition}
\begin{proof}
The three regimes are (T3$'$), the move of item 2 of \cref{thm:common-tail}, and (T2$'$), whose proofs use only that $M$ is maximum in $K^{+}$, respectively that the mover never stands on a tail. The example is an exhaustive computation.
\end{proof}

\begin{proposition}[the barrier for the other geometries]\label{prop:barrier}
For each geometry of two one-way arcs other than a common tail (chained: $u_1\to w_1=u_2\to w_2$; common head; disjoint) there are two positions with identical key-local data and opposite outcomes:
\begin{itemize}
\item \emph{Chained.} The first player wins on the graph with vertices $0,\dots,5$, edges $0\text{--}2$, $1\text{--}2$, $1\text{--}4$, $2\text{--}5$, $3\text{--}4$, $4\text{--}5$, arcs $0\to3\to1$ and start $5$. The first player loses on the graph with vertices $0,\dots,9$, edges $0\text{--}6$, $0\text{--}7$, $0\text{--}8$, $0\text{--}9$, $1\text{--}6$, $1\text{--}7$, $1\text{--}9$, $2\text{--}6$, $2\text{--}8$, $2\text{--}9$, $3\text{--}6$, $3\text{--}8$, $3\text{--}9$, $4\text{--}6$, $4\text{--}7$, $4\text{--}9$, $5\text{--}6$, $5\text{--}7$, $5\text{--}9$, arcs $1\to3\to4$ and start $0$.
\item \emph{Common head.} The first player wins on the graph with vertices $0,\dots,7$, edges $0\text{--}5$, $0\text{--}6$, $1\text{--}5$, $1\text{--}6$, $2\text{--}6$, $2\text{--}7$, $3\text{--}5$, $4\text{--}6$, $4\text{--}7$, arcs $1\to4$, $0\to4$ and start $3$. The first player loses on the graph with vertices $0,\dots,7$, edges $0\text{--}5$, $0\text{--}6$, $0\text{--}7$, $1\text{--}5$, $1\text{--}6$, $1\text{--}7$, $2\text{--}5$, $2\text{--}7$, $3\text{--}5$, $3\text{--}6$, $4\text{--}5$, $4\text{--}6$, $4\text{--}7$, arcs $4\to0$, $1\to0$ and start $2$.
\item \emph{Disjoint.} The graph and arcs of \cref{prop:two-arcs}, with starts $5$ (first player wins) and $3$ (first player loses).
\end{itemize}
\end{proposition}
\begin{proof}
Direct computation of both outcomes and of the key-local data (\code{star_tail.key_local_data}, \code{tests/test_star_tail.py}).
\end{proof}

\begin{corollary}[dichotomy for two arcs]\label{cor:dichotomy}
Key-local data decide every position of $\game(H;\{a_1,a_2\})$, for all graphs $H$, if and only if $a_1$ and $a_2$ share a tail.
\end{corollary}
\begin{proof}
For a common tail the rule of \cref{thm:common-tail} uses only key-local quantities, by \eqref{eq:star-nu}. The other geometries are \cref{prop:barrier}.
\end{proof}

A common tail is special because all its arcs can be used only while the token stands on $u$, so the arc system acts at one moment, and at that moment the game is undirected geography on $K+uW'$. With a common head, a chain or disjoint arcs, one play can use arcs at two different moments, and whether the first player can arrange a favourable second crossing depends on the play in between. A sampled search over $26{,}633$ starts with two arcs, classified by key-local data into $4{,}220$ classes, found conflicting classes only for chained and co-headed arcs (none for a common tail), and a five-vertex chained conflict already exists: on $K_{2,3}$ with parts $\{0,1\}$ and $\{2,3,4\}$ and arcs $3\to2\to4$ the first player wins from $0$, while after deleting the edge $0\text{--}3$ and reversing the chain to $4\to2\to3$ the first player loses from $0$, with the same matching numbers for every key-point deletion.

\subsection{Several tails}

\begin{lemma}[$p$ tails]\label{lem:p-tails}
Let the arcs of $\arcs$ have exactly $p$ distinct tails $u_1,\dots,u_p$, and let $\arcs_i$ be the set of arcs leaving $u_i$.
\begin{enumerate}[label=\textup{(\arabic*)}]
\item Every play uses at most one arc of each $\arcs_i$, hence at most $p$ arcs.
\item Let the token stand on $u_i$ in the current graph $K'$, let $H_i$ be the set of live heads of $\arcs_i$, and let $\arcs'$ be the set of live arcs of the other tails whose head is not $u_i$. This position has the outcome of the position $(K'+u_iH_i,u_i)$ of $\game(K'+u_iH_i;\arcs')$, which has $p-1$ tails.
\item Until the token first stands on a tail the play is undirected vertex geography on $K$. Consequently a $p$-tail game is undirected vertex geography with $p$ terminals $u_1,\dots,u_p$, where entering $u_i$ ends the play with the value of a $(p-1)$-tail game on the remaining graph; for $p=2$ that value is given by \cref{thm:common-tail}.
\end{enumerate}
\end{lemma}
\begin{proof}
An arc of $\arcs_i$ is used only by a move from $u_i$, and a play leaves $u_i$ at most once, which is (1). For (2), the moves from $u_i$ go to its neighbours in $K'$ and to the heads in $H_i$, which are exactly its neighbours in $K'+u_iH_i$. After any such move $u_i$ is visited, so the added edges, every arc of $\arcs_i$ and every arc with head $u_i$ are dead, and from then on both games are $\game(K'-u_i;\arcs')$ with the same live arcs; the two game trees coincide. The arcs of $\arcs'$ avoid $u_i$, and every added edge contains $u_i$, so they still join non-adjacent vertices of $K'+u_iH_i$, as the definition of $\game$ requires; and they have at most $p-1$ tails, none equal to $u_i$. For (3), no arc leaves a vertex that is not a tail.
\end{proof}

For $p=2$ every position at or after the first crossing is decided by at most six maximum-matching computations, and the difficulty is confined to the undirected phase before it. Gallai--Edmonds data at the key points do not decide that phase.

\begin{proposition}[Gallai--Edmonds data do not decide two tails]\label{prop:ge-barrier}
\begin{enumerate}[label=\textup{(\arabic*)}]
\item GE-local data separate the two positions of each pair in \cref{prop:barrier}.
\item GE-local data do not decide two-tail positions. Let $G_6$ have vertices $s,h,h',a,b,\ell$ and edges $sh$, $hh'$, $ha$, $hb$, $h'a$, $h'b$, $h'\ell$. With the arcs $a\to b\to\ell$ the first player wins from $s$; with the reversed chain $\ell\to b\to a$ the first player loses from $s$. The two positions have the same GE-local data.
\end{enumerate}
\end{proposition}
\begin{proof}
Item (1) and the equality of the data in (2) are direct computations over all sets of key points. For the outcomes in (2) the first move $s\to h$ is forced. With the chain $a\to b\to\ell$: if the second player moves to $h'$, the first moves to $\ell$ and the second is stuck; if the second moves to $a$, the first crosses to $b$, and the second must move to $h'$ or cross to $\ell$, which the first answers with $\ell$ or $h'$, leaving the second without a move; if the second moves to $b$, the first crosses to $\ell$, the second is forced to $h'$, and the first moves to $a$, where the second is stuck because the arc $a\to b$ leads to a visited vertex. With the chain $\ell\to b\to a$ the second player answers $h\to a$; no arc leaves $a$, so the first must move to $h'$; the second moves to $b$, and the first is stuck, because $h,h'$ are visited and the arc $b\to a$ leads to a visited vertex.
\end{proof}

A sampled search over $38{,}805$ two-tail starts on graphs with at most seven vertices found four such conflicts, two chained and two with a common head. In $G_6$ the matching structure around the key points is the same in both positions; what differs is the direction in which the chain can be traversed relative to where the play meets it.

\begin{remark}[where the complexity stands]\label{rem:tails-complexity}
DVG is $\PSPACE$-complete in general \cite{LichtensteinSipser} and polynomial when the one-way arcs of every strong component share a tail (\cref{cor:star-dvg}). Unless $\Pclass=\PSPACE$, hard instances need strong components with many tails. \Cref{lem:p-tails} confines the difficulty to undirected phases that end at matching-type terminals, and \cref{prop:barrier,prop:ge-barrier} show that no rule reading only the matching structure at the key points, through deficiencies of key-point deletions or Gallai--Edmonds classes, decides the two-tail phase. A polynomial algorithm would have to follow how the play approaches the tails; a hardness proof would have to encode a formula in an undirected phase that ends at one of two matching-type terminals. Whether DVG is fixed-parameter tractable in the number of tails is open already at two (\cref{prob:two-tails}). A bag invariant of a tree decomposition does not bear on this: the barrier graphs have small width, and DVG on bounded treewidth is linear \cite{Bodlaender}.
\end{remark}

\section{Set potentials}\label{sec:setpot}

A local weighted matching potential produces every resolvent from one vertex-deletion recurrence. Dropping the recurrence leaves an algebraic consistency condition on the game alone: the products of resolvents along plays must depend only on the set of visited vertices. This section classifies that condition on the families where it can be decided in closed form. A set potential carries no information about complexity; it is a structural invariant of the generating functions of the game.

\subsection{Three kinds of potential}\label{ssec:hierarchy}
Let $\mathcal W$ be the family of the sets $W$ and $W-v$ over the reachable positions $(W,v)$.
\begin{itemize}
\item A \emph{local weighted matching potential} is the object of \cref{thm:conj610}.
\item A \emph{set potential} is a map $\Mpot\colon\mathcal W\to\Qz^{\times}$ with $\Mpot(W-v)=R(W,v)\Mpot(W)$ at every reachable position (as in \cref{thm:cocycle}).
\item A \emph{valuation potential} is a map $\Lambda\colon\mathcal W\to\ZZ$ with $\Lambda(W-v)-\Lambda(W)=\ord_0R(W,v)$, which by \cref{thm:resolvent-outcome} is $+1$ at N-positions and $-1$ at P-positions.
\end{itemize}
Every set in $\mathcal W$ has the form $V\setminus V(P)$ for a directed path $P$, and $\Mpot(V\setminus V(P))=\Gamma(P)\Mpot(V)$; so a set potential exists exactly when $\Gamma(P)$ depends only on $V(P)$, and whether it does is decided exactly over $\Qz$ by propagating $\Mpot$ from $V$ (\code{potentials.set_potential}). By \cref{lem:field-independence} the answer does not change over any larger field.

\begin{theorem}[hierarchy]\label{thm:hierarchy}
A matching potential gives a set potential, $\Mpot=\mu_a(U[\cdot])$, and a set potential gives a valuation potential, $\Lambda=\ord_0\Mpot$. All three exist on SCC-symmetric digraphs. Both implications are strict: a directed cycle has a set potential but no matching potential, and the digraph with arcs $0\rightleftarrows1$, $0\to2$, $2\to1$, $2\to3$ has a valuation potential but no set potential. The digraph $T$ with arcs $0\rightleftarrows1$, $0\rightleftarrows2$, $1\to2$ has none of the three.
\end{theorem}
\begin{proof}
On a directed cycle every nonempty unvisited set other than $V$ is a cyclic interval reached by a unique play, and $\emptyset$ is reached from all $k$ starts, each play with product $1/\mu(P_k)$; so a set potential exists, and \cref{prop:fails}(1) excludes a matching potential. The four-vertex digraph has the Hamiltonian cycle $0\to2\to1\to0$ on $\{0,1,2\}$, where vertex $0$ has two out-neighbours and one in-neighbour, so \cref{thm:cocycle} excludes a set potential; a valuation potential exists by direct computation. For $T$ put $\Lambda(V)=0$. The position $(V,0)$ is an N-position (moving to $2$ leaves the opponent stuck), so $\Lambda(\{1,2\})=1$; $(\{1,2\},1)$ is N, so $\Lambda(\{2\})=2$; $(V,1)$ is P, since both options are N, so $\Lambda(\{0,2\})=-1$; and $(\{0,2\},0)$ is N, so $\Lambda(\{2\})=0$, a contradiction.
\end{proof}

\begin{center}\small
\begin{tabular}{@{}lrrr@{}}
\toprule
Non-SCC-symmetric digraphs, up to isomorphism & $3$ vertices & $4$ vertices & $5$ vertices\\ \midrule
no valuation potential & $2$ & $83$ & $6{,}566$\\
valuation potential only & $0$ & $25$ & $1{,}410$\\
set potential, no matching potential & $1$ & $6$ & $79$\\ \midrule
total & $3$ & $114$ & $8{,}055$\\
\bottomrule
\end{tabular}
\end{center}
Since $T$ has three vertices, the obstruction to valuation potentials is local and has nothing to do with width (\S\ref{sec:treewidth}).

\subsection{Components and symmetric branches}
Games with self-energies are as in \S\ref{sec:realstab}. A self-energy \emph{vanishes at infinity} if it lies in $\Qz$ (or $\Puis$) and tends to $0$ as $z\to\infty$; those produced by exits and by symmetric branches do.

\begin{lemma}[components]\label{lem:set-components}
For a strong component $C$ and $x\in C$ let $\eta_x$ be the sum of $\Rdown(t)$ over the out-neighbours $t$ of $x$ outside $C$. Then $D$ has a set potential if and only if every non-symmetric strong component $C$, played alone with the self-energies $\eta$, has one.
\end{lemma}
\begin{proof}
A directed path meets the strong components in topological order and never returns to one. While the token is in $C$, every vertex it can reach outside $C$ lies strictly below $C$ and is unvisited (\cref{lem:entry}), so an exit to $t$ contributes exactly $\Rdown(t)$. Hence $\Gamma(P)$ is the product, over the components met by $P$, of $\Gamma$ of its segment in the component game with self-energies $\eta$, and a segment inside $C$ is itself a directed path of $D$. Two paths with the same vertex set meet the same components in the same vertex sets, so $D$ passes if every component game does, and a failing pair inside one component is a failing pair of $D$. A symmetric component always passes: by Godsil's identity with vertex weights $z-\eta$, each factor is $\mu(G[W]-v)/\mu(G[W])$ and the product telescopes.
\end{proof}

So everything upstream of a component, and everything incomparable to it, is irrelevant to set potentials.

\begin{lemma}[symmetric branches act as exits]\label{lem:branches}
Let $C$ carry self-energies $\eta$, let $h\in C$, and let $B\subseteq C\setminus\{h\}$ be nonempty, with $D[B\cup\{h\}]$ symmetric and no arc between $B$ and $A\setminus\{h\}$, where $A=C\setminus B$. With vertex weights $z-\eta$ on $B$ put
\[
\sigma=\sum_{t\in N(h)\cap B}\frac{\mu(B-t)}{\mu(B)},
\]
the sum of the root resolvents of undirected geography on $B$ at the neighbours of $h$. Then $(C,\eta)$ has a set potential if and only if $(A,\eta')$ has one, where $\eta'_h=\eta_h+\sigma$ and $\eta'=\eta$ elsewhere.
\end{lemma}
\begin{proof}
Let $\Gamma'$ and $R'$ refer to the game $(A,\eta')$, and let $\tau=\eta_h+\sum_aR_{A-h}(a)$ over the out-neighbours $a$ of $h$ in $A$, where $R_{A-h}$ is the fresh resolvent of the game on $A-h$. Sort the directed paths $P$ of $(C,\eta)$ with vertex set $X$ into four kinds: (0) $X\cap B=\emptyset$; (1) $X\cap B\neq\emptyset$ and $X\subseteq B\cup\{h\}$; (2) $X$ meets $B$ and $A\setminus\{h\}$, and $P$ starts in $B$; (3) $X$ meets $B$ and $A\setminus\{h\}$, and $P$ starts in $A\setminus\{h\}$. A path of kind (2) or (3) passes through $h$ exactly once, because $B$ is joined to the rest only through $h$; it has the form $\beta\cdot h\cdot\alpha$ or $\alpha\cdot h\cdot\beta$ with $\beta$ in $B$ and $\alpha$ in $A\setminus\{h\}$.

Whenever the token stands in $A$ with $B$ unvisited, a move from $h$ into $B$ starts undirected geography on the fresh set $B$, so $B$ contributes exactly $\sigma$ to $1/R(W,h)$; by induction on the game tree every such position has the resolvent of $(A,\eta')$. Hence $\Gamma(P)=\Gamma'(P)$ for kind (0). While $A\setminus\{h\}$ is untouched, the token in $B\cup\{h\}$ plays undirected geography on $G=D[B\cup\{h\}]$ with weights $z-\eta$ on $B$ and $z-\tau$ on $h$, and once $h$ is visited it plays on $B$ alone. Therefore, by Godsil's identity,
\begin{itemize}
\item kind (1): $\Gamma(P)=\mu(G-X)/\mu(G)$;
\item kind (2): $\Gamma(\beta\cdot h\cdot\alpha)=\dfrac{\mu(B\setminus V(\beta))}{\mu(G)}\,\Gamma_{A-h}(\alpha)=\dfrac{\mu(B\setminus V(\beta))}{\mu(B)}\,\Gamma'(h\cdot\alpha)$, because $\mu(G)=\mu(B)(z-\tau-\sigma)$ and $R'(A,h)=1/(z-\tau-\sigma)$;
\item kind (3): $\Gamma(\alpha\cdot h\cdot\beta)=\Gamma'(\alpha\cdot h)\,\dfrac{\mu(B\setminus V(\beta))}{\mu(B)}$, because $B$ is fresh when the token enters it.
\end{itemize}
If $(C,\eta)$ has a set potential, so does $(A,\eta')$, since its paths are the kind-(0) paths. Conversely, suppose $(A,\eta')$ has one and let $P,P'$ have the same vertex set $X$. They have the same kind, except that kinds (2) and (3) can meet. In kinds (0) and (1) the formulas depend only on $X$. In kinds (2) and (3), $\Gamma$ is $\mu(B\setminus X)/\mu(B)$ times $\Gamma'$ of a path of $(A,\eta')$ with vertex set $X\cap A$, which contains $h$, and that value depends only on $X\cap A$.
\end{proof}

Repeated pruning removes every symmetric branch. A \emph{pruned core} is a strongly connected, non-symmetric digraph without symmetric branches.

\subsection{Directed cycles and books}\label{ssec:setpot-cycles}
For weights $w_1,\dots,w_n$ and couplings $\beta_2,\dots,\beta_n$ the \emph{continuant} $\cont(w_1,\dots,w_n)$ is defined by $\cont_0=1$, $\cont_1=w_1$ and $\cont_i=w_i\cont_{i-1}-\beta_i\cont_{i-2}$, where $\cont_i$ is the continuant of the first $i$ entries; unlisted couplings are $1$. It is the weighted matching polynomial of a path with edge weights $\beta$, it is unchanged when the sequence (with its couplings) is reversed, and two sequences $\sigma,\tau$ joined with coupling $\beta$ satisfy the \emph{junction rule}
\[
\cont(\sigma,\tau)=\cont(\sigma)\cont(\tau)-\beta\,\cont(\sigma^{-})\cont({}^{-}\tau),
\]
where $\sigma^{-}$ drops the last entry of $\sigma$ and ${}^{-}\tau$ the first entry of $\tau$. For a directed path with weights $w_i=z-\eta_i$ the product of resolvents along the whole path is $1/\cont(w_1,\dots,w_n)$. With $T_i=\bigl(\begin{smallmatrix}w_i&-1\\1&0\end{smallmatrix}\bigr)$,
\[
T_n\cdots T_1=\begin{pmatrix}\cont(w_1,\dots,w_n)&-\cont(w_2,\dots,w_n)\\ \cont(w_1,\dots,w_{n-1})&-\cont(w_2,\dots,w_{n-1})\end{pmatrix}.
\]

\begin{theorem}[directed cycles]\label{thm:set-cycles}
A chordless directed cycle $x_0\to x_1\to\dots\to x_{m-1}\to x_0$ with self-energies vanishing at infinity has a set potential if and only if $\eta(x_{i+2})=\eta(x_i)$ for every $i$, indices mod $m$. In particular $\eta$ is constant when $m$ is odd. Then all rotation continuants $\cont_i=\cont(w_i,\dots,w_{i-1})$ agree; write $\cont(C)$ for their value.
\end{theorem}
\begin{proof}
A proper vertex set carried by a path is a cyclic interval, carried by one path, so the only condition is that the $m$ rotations of the whole cycle have equal products $\Gamma=1/\cont_i$. The product $\Pi_i=T_{i-1}\cdots T_{i+1}T_i$ equals $\bigl(\begin{smallmatrix}\cont_i&-B_i\\C_i&-c_i\end{smallmatrix}\bigr)$, where $B_i,C_i,c_i$ are the continuants of $w_{i+1},\dots,w_{i-1}$, of $w_i,\dots,w_{i-2}$ and of $w_{i+1},\dots,w_{i-2}$. Since $\Pi_{i+1}T_i=T_i\Pi_i$, all $\Pi_i$ have the same trace $\cont_i-c_i$. If $\cont_i=A$ for all $i$, then $c_i=c$ for all $i$, and comparing entries of $\Pi_{i+1}T_i=T_i\Pi_i$ gives $B_{i+1}=C_i$, $C_{i+1}=B_i$ and $w_iB_i=A+c$. Because $\eta$ vanishes at infinity, $B_i$ and $A+c$ have degrees $m-1$ and $m$, so $w_i=(A+c)/B_i$ and $B_{i+2}=B_i$ give period $2$. Conversely, if $w$ has period $2$, reversing the sequence $w_{i+1},\dots,w_{i-2}$ gives $w_{i+2},\dots,w_{i-1}$, so $c_i=c_{i+1}$ and $\cont_i=\text{trace}+c_i$ is constant.
\end{proof}

\begin{theorem}[books]\label{thm:books}
The $r$-page book $B_r$ has a spine $x\to y$ and pages $p_1,\dots,p_r$ with arcs $y\to p_i\to x$; it is the directed triangle with one vertex blown up into $r$ twins. With self-energies vanishing at infinity it has a set potential if and only if
\[
\eta(y)=\eta(p_1)=\dots=\eta(p_r)=e,\qquad \eta(x)=e+\frac{r-1}{z-e}.
\]
By \cref{lem:branches} the extra term can be supplied by $r-1$ pendant vertices at $x$ with self-energy $e$, or, when $e=0$, by exits from $x$ to $r-1$ sinks.
\end{theorem}
\begin{proof}
Write $z_v=z-\eta(v)$, $S=\sum_j1/z_{p_j}$, $S_i=S-1/z_{p_i}$ and $T=\sum_jz_x/(z_{p_j}z_x-1)$. The vertex sets carried by several paths are the triangles $\{x,y,p_i\}$, with their three rotations, and the sets $\{p_i,x,y,p_j\}$, carried by $p_ixyp_j$ and $p_jxyp_i$. Direct computation gives
\begin{gather*}
\Gamma(xyp_i)=\frac{1}{(z_x(z_y-S)-1)\,z_{p_i}},\qquad \Gamma(yp_ix)=\frac{1}{(z_y-T)(z_{p_i}z_x-1)},\\
\Gamma(p_ixy)=\frac{1}{D_i},\qquad \Gamma(p_ixyp_j)=\frac{1}{D_i\,z_{p_j}},
\end{gather*}
with $D_i=z_{p_i}(z_x(z_y-S_i)-1)-(z_y-S_i)$. Equating the first two gives $z_{p_i}(z_x(T-S)-1)=T-z_y$. Since $\eta$ vanishes at infinity, $z_x(T-S)=O(z^{-2})$, so all pages have the same $z_{p_i}=w$, and then the equation reads $z_y=w$. Equating the first and third gives $z_x=w-(r-1)/w$. The sets $\{p_i,x,y,p_j\}$ agree because the pages are now interchangeable. Conversely these values satisfy all equations.
\end{proof}

\subsection{Pruned cores on at most five vertices}

\begin{theorem}[pruned cores on at most five vertices]\label{thm:cores5}
Among pruned cores with at most five vertices and self-energies vanishing at infinity, exactly the following admit a set potential, with exactly these self-energies.
\begin{center}\small
\begin{tabularx}{\linewidth}{@{}L{0.2}L{0.3}>{\raggedright\arraybackslash}X@{}}
\toprule
Pruned core & Arcs & Admissible self-energies\\ \midrule
$C_3$, $C_4$, $C_5$ & the cycle & period $2$ around the cycle (\cref{thm:set-cycles})\\
books $B_2$, $B_3$ & $x\to y\to p_i\to x$ & \cref{thm:books}\\
$C_4$ with one vertex doubled & $s_0\to s_1\to s_2\to\{s_3,s_3'\}\to s_0$ & $\eta(s_3)=\eta(s_3')=e$, $\eta(s_1)=e+1/(z-\eta(s_2))$, $\eta(s_0)=\eta(s_2)+1/(z-e)$\\
$C_3$ with two vertices doubled & $s_0\to\{s_1,s_1'\}\to\{s_2,s_2'\}\to s_0$ & $\eta=a$ on $s_0,s_2,s_2'$ and $\eta=a-1/(z-a)$ on $s_1,s_1'$\\
bowtie & two directed triangles sharing a vertex $h$ & $\eta=e$ off $h$ and $\eta(h)=e-1/((z-e)-1/(z-e))$\\
\bottomrule
\end{tabularx}
\end{center}
Every other strongly connected, non-symmetric digraph with at most five vertices either admits no self-energies vanishing at infinity, or has a symmetric branch and reduces by \cref{lem:branches} to a core in the table, with the self-energies the lemma predicts.
\end{theorem}
\begin{proof}[Method \textup{(computer-assisted)}]
The self-energies are indeterminates over $\Qz$. On three and four vertices every strongly connected, non-symmetric digraph was treated. On five vertices, self-energies vanishing at infinity enter \cref{lem:asymptotics} like out-degrees, so the back-arc test of \cref{thm:cocycle} is necessary; it leaves $14$ of the $5{,}027$ strongly connected, non-symmetric digraphs. For each digraph the equations $\Gamma(P)=\Gamma(P')$ over paths with equal vertex sets generate an ideal of $\Qz[\eta]$, which is saturated by every numerator and denominator of $\Gamma$ and decomposed into minimal primes; components on which some $\eta$ cannot vanish at infinity are discarded (one for the doubled triangle, with $\eta(s_1)=z$, and one for the bowtie, with $\eta(h)\sim z/2$).
\end{proof}

The cycles, the books and the doubled $4$-cycle are realized by exits; for instance, the doubled $4$-cycle with sinks at $s_0$ and $s_1$ has a set potential. For the doubled triangle and the bowtie, the admissible self-energy at one vertex class is smaller than at its neighbours by a resolvent, and no exit structure among the $8{,}844$ sums of at most three root resolvents of games on at most four vertices realizes either. With $e=0$ this is forced: a sum of root resolvents has positive leading coefficient.

\begin{corollary}[the census on at most five vertices]\label{cor:census5}
Every non-SCC-symmetric digraph with a set potential and at most five vertices has exactly one non-symmetric strong component, and its pruned core is a directed cycle or the $2$-page book. On five vertices the $79$ digraphs are: $52$ bare directed triangles (zero exit self-energies), $9$ triangles with self-energy $1/z$ at every vertex, $3$ with $z/(z^2-1)$ at every vertex, one each with $2/z$, $2z/(z^2-1)$ and $1/z+z/(z^2-1)$ at every vertex, one triangle with a pendant vertex and $1/z$ at the other two vertices; $6$ bare directed $4$-cycles, one with $1/z$ at every vertex and one with $1/z$ at two opposite vertices; one bare $5$-cycle; one $2$-page book with $1/z$ at the spine tail and one with a pendant vertex there. On three and four vertices the $1+6$ cases are the directed triangle, bare or with a common sink, a source or an in-pendant, and the directed $4$-cycle.
\end{corollary}
\begin{proof}
\Cref{lem:set-components,lem:branches}, \cref{thm:set-cycles,thm:books,thm:cores5}, applied to the census of \cref{thm:hierarchy}; every entry was confirmed by the exact test.
\end{proof}

\begin{remark}[the mechanism]\label{rem:mechanism}
There is no global Eulerian compensation. By \cref{lem:set-components} everything upstream of a component and everything incomparable to it is irrelevant, which is why $52$ of the $79$ are a bare triangle with arbitrary other attachments. Inside a component the mechanism is a balance of effective self-energies, to which exits and pruned branches contribute only through their resolvents. The hidden symmetry is the period-$2$ invariance of these self-energies around a cycle, together with the leaf compensation of twin pages in books; it need not come from an automorphism (in the $2$-page book with a sink at its spine tail no automorphism rotates a page triangle, yet its three rotations have equal $\Gamma$). The Eulerian condition of \cref{thm:cocycle} is the first-order shadow of this balance and is blind to exits, which enter \cref{lem:asymptotics} like out-degrees and cancel; a $4$-cycle with sinks at two adjacent vertices satisfies \cref{thm:cocycle} but has no set potential.
\end{remark}

\subsection{Blow-ups of directed cycles}
For sizes $s_0,\dots,s_{m-1}\ge1$ the \emph{blow-up} $C_m(s_0,\dots,s_{m-1})$ replaces $x_i$ by a class $S_i$ of $s_i$ twins, with an arc from every vertex of $S_i$ to every vertex of $S_{i+1}$, indices mod $m$. For $m\ge3$ it is a pruned core. The book $B_r$ is $C_3(r,1,1)$, the doubled $4$-cycle is $C_4(2,1,1,1)$, and the doubled triangle of \cref{thm:cores5} is $C_3(1,2,2)$.

\begin{theorem}[continuant criterion]\label{thm:blowup-criterion}
Let $D=C_m(s_0,\dots,s_{m-1})$ with $m\ge3$ carry a constant self-energy $e$, and put $y=z-e$. For $j\in\ZZ/m$ let
\[
b^{(j)}_k=s_{j+k+1}-\Bigl\lfloor\frac{k+1}{m}\Bigr\rfloor\quad(k\ge0),\qquad L_j=1+\min\{k:b^{(j)}_k=0\},
\]
and define the \emph{backward continuants} by $\cont^{(j)}_{L_j}=1$, $\cont^{(j)}_{L_j+1}=0$ and
\[
\cont^{(j)}_k=y\,\cont^{(j)}_{k+1}-b^{(j)}_k\cont^{(j)}_{k+2}\qquad(0\le k<L_j).
\]
\begin{enumerate}[label=\textup{(\arabic*)}]
\item Every directed path $P$ that starts in class $j$ and has $n$ vertices has $\Gamma(P)=\cont^{(j)}_n/\cont^{(j)}_0$.
\item $D$ has a set potential if and only if, for every $c$ with $1\le c\le\min_is_i$, the ratio $\cont^{(j)}_{cm}/\cont^{(j)}_0$ does not depend on $j$.
\end{enumerate}
\end{theorem}
\begin{proof}
(1) Let the play start at a vertex of $S_j$. At depth $k$ the token is in class $j+k$, and the vertices of $S_{j+k+1}$ visited so far are those visited at the depths $d\le k$ with $d\equiv k+1\pmod m$, of which there are $\lfloor(k+1)/m\rfloor$. Since the out-neighbours of every vertex of $S_i$ are exactly the vertices of $S_{i+1}$, every node at depth $k$ of the game tree has exactly $b^{(j)}_k$ children, all with isomorphic subtrees; the tree is spherically symmetric, and the play ends at depth $L_j-1$. With level resolvents $R_k$ this gives $R_k=1/(y-b^{(j)}_kR_{k+1})$ and $R_{L_j-1}=1/y$, so $R_k=\cont^{(j)}_{k+1}/\cont^{(j)}_k$. The $i$-th position of $P$ is a node at depth $i$, so $\Gamma(P)=\prod_{i<n}\cont^{(j)}_{i+1}/\cont^{(j)}_i=\cont^{(j)}_n/\cont^{(j)}_0$.

(2) By (1), $\Gamma(P)$ depends only on the start class and the number of vertices. A path with $n$ vertices from class $j$ meets class $i$ in as many vertices as there are $k<n$ with $j+k\equiv i$. If $m\nmid n$, these counts determine $j$, so two paths with the same vertex set have the same $\Gamma$. If $n=cm$, every class is met in $c$ vertices, which requires $c\le\min_is_i$; conversely, for such $c$ and any set $X$ with $c$ vertices in every class, and for every $j$, listing the vertices of $X$ class by class from $S_j$ gives a directed path with vertex set $X$ that starts in class $j$, because consecutive classes are completely joined. So the set potential exists exactly when $\cont^{(j)}_{cm}/\cont^{(j)}_0$ is independent of $j$ for these $c$.
\end{proof}

\begin{lemma}[twin identity]\label{lem:twin}
Let $m\ge4$ be even, $s>t\ge1$, and let $s_i=s$ for even $i$ and $s_i=t$ for odd $i$. Then $L_0=tm+1$, $L_1=tm$, and
\[
\cont^{(0)}_k=y\,\cont^{(1)}_k\qquad\text{for every even }k\text{ with }0\le k\le tm.
\]
\end{lemma}
\begin{proof}
Write $b=b^{(0)}$ and $b'=b^{(1)}$: $b_k=t-\lfloor(k+1)/m\rfloor$ for even $k$ and $s-\lfloor(k+1)/m\rfloor$ for odd $k$, while $b'_k=s-\lfloor(k+1)/m\rfloor$ for even $k$ and $t-\lfloor(k+1)/m\rfloor$ for odd $k$. Since $s>t$, $b_k\ge1$ for $k<tm$ and $b_{tm}=0$, and $b'_k\ge1$ for $k<tm-1$ and $b'_{tm-1}=0$; this gives $L_0,L_1$. Eliminating the odd levels from the backward recursion: for $2\le k\le L-1$,
\[
\cont_{k-2}=(y^{2}-b_{k-2}-b_{k-1})\,\cont_k-b_{k-1}b_k\,\cont_{k+2},
\]
using $y\cont_{k+1}=\cont_k+b_k\cont_{k+2}$. Put $E_j=\cont^{(0)}_{tm-2j}$ and $E'_j=\cont^{(1)}_{tm-2j}$. For even $k$ the sums $b_{k-2}+b_{k-1}$ and $b'_{k-2}+b'_{k-1}$ both equal $s+t-\lfloor(k-1)/m\rfloor-\lfloor k/m\rfloor$; and since $m$ is even and $k+1$ is odd, $\lfloor(k+1)/m\rfloor=\lfloor k/m\rfloor=:q$, so the products $b_{k-1}b_k$ and $b'_{k-1}b'_k$ both equal $(s-q)(t-q)$. Hence $E$ and $E'$ satisfy the same three-term recursion for $j\ge1$. At the start, $E_0=\cont^{(0)}_{tm}=y$ and $E'_0=\cont^{(1)}_{tm}=1$; and $E_1=(y^{2}-b_{tm-2}-b_{tm-1})\,y=(y^{2}-1-s+t)\,y$, while $\cont^{(1)}_{tm-1}=y$ and $E'_1=y^{2}-b'_{tm-2}=y^{2}-(s-t+1)$. So $E_j=yE'_j$ for $j=0,1$, hence for all $j\le tm/2$.
\end{proof}

\begin{corollary}[two-periodic blow-ups]\label{cor:two-periodic}
With a constant self-energy, every blow-up $C_m(s_0,\dots,s_{m-1})$ with $s_{i+2}=s_i$ for all $i$ has a set potential. For odd $m$ these are the balanced blow-ups $C_m(s,\dots,s)$.
\end{corollary}
\begin{proof}
If all sizes are equal, all $b^{(j)}$ coincide. Otherwise $m$ is even and, after a rotation, $s_0=s>t=s_1$; the sequences $b^{(j)}$ depend only on the parity of $j$. For $1\le c\le t$ the index $cm$ is even, so \cref{lem:twin} gives $\cont^{(0)}_{cm}/\cont^{(0)}_0=y\cont^{(1)}_{cm}/(y\cont^{(1)}_0)$, and \cref{thm:blowup-criterion} applies.
\end{proof}

The earlier draft proved the balanced case by an automorphism argument and found set potentials on $C_4(1,2,1,2)$, $C_4(1,3,1,3)$, $C_4(2,3,2,3)$ and $C_6(1,2,1,2,1,2)$ by exact computation, attributing them to a hidden symmetry; \cref{lem:twin} is that symmetry. As a check independent of the proofs, the formula of \cref{thm:blowup-criterion}(1) was compared with the products of resolvents along every directed path of seven blow-ups, the criterion with the direct propagation test on all $115$ blow-ups with at most nine vertices (up to rotation; exactly the $13$ two-periodic ones have a set potential), and \cref{lem:twin} was confirmed symbolically for $m\le10$ (\code{blowups.py}, \code{tests/test_blowups.py}).

\begin{theorem}[one blown-up vertex]\label{thm:one-blown}
Let $m\ge3$ and $r\ge2$, and let $D=C_m(r,1,\dots,1)$, with twins $t_1,\dots,t_r$ and arcs $t_j\to x_1\to x_2\to\dots\to x_{m-1}\to t_j$. With self-energies vanishing at infinity, $D$ has a set potential if and only if all twins have the same self-energy $e$ and, writing $\eta_i=\eta(x_i)$:
\begin{itemize}
\item for $m$ odd, $\eta_{m-1}=e$ and $\eta_i=e+(r-1)/(z-e)$ for $1\le i\le m-2$;
\item for $m$ even, $\eta_i=\eta_{m-1}+(r-1)/(z-e)$ for odd $i\le m-3$, and $\eta_i=e+(r-1)/(z-\eta_{m-1})$ for even $i$ with $2\le i\le m-2$.
\end{itemize}
For $m=3$ this is \cref{thm:books}, and for $m=4$, $r=2$ the doubled $4$-cycle of \cref{thm:cores5}. With $e=\eta_{m-1}=0$ each of $x_1,\dots,x_{m-2}$ needs self-energy $(r-1)/z$, which $r-1$ exits to sinks provide.
\end{theorem}
\begin{proof}
Write $w_v=z-\eta(v)$, $w_i=w(x_i)$ and $v_j=w(t_j)$; every continuant below has a power of $z$ as leading term, so it is nonzero.

\emph{Vertex sets.} A directed path moves forward through the positions $0,1,\dots,m-1$ of the cycle, position $0$ being the twin class; it meets each $x_i$ at most once, and meets two twins only as $t_jx_1\cdots x_{m-1}t_{j'}$. A proper cyclic interval of positions has a unique first position, so a vertex set carried by more than one path contains every position, and the only such sets are $Z_j=\{t_j,x_1,\dots,x_{m-1}\}$, carried by its $m$ rotations, and $Z_j\cup\{t_{j'}\}$, carried by $t_jx_1\cdots x_{m-1}t_{j'}$ and $t_{j'}x_1\cdots x_{m-1}t_j$.

\emph{Products.} Along each of these paths the path tree of every position is a path, except that at $x_{m-1}$ the unvisited twins hang as branches: a single vertex when $x_1$ has been visited, and otherwise, on a path that starts at $x_a$, the path $t_{j'},x_1,\dots,x_{a-1}$. Each branch enters the effective weight of $x_{m-1}$ through its root resolvent, and the products telescope. With $S_j=\sum_{j'\neq j}1/v_{j'}$ and $\Sigma=\sum_{j'}1/v_{j'}$,
\begin{gather*}
\Gamma(t_jx_1\cdots x_{m-1})=\frac{1}{\cont(v_j,w_1,\dots,w_{m-2},w_{m-1}-S_j)},\\
\Gamma(x_1\cdots x_{m-1}t_j)=\frac{1}{v_j\,\cont(w_1,\dots,w_{m-2},w_{m-1}-\Sigma)},
\end{gather*}
and $\Gamma(t_jx_1\cdots x_{m-1}t_{j'})=1/\bigl(v_{j'}\cont(v_j,w_1,\dots,w_{m-2},w_{m-1}-S_j)\bigr)$.

\emph{The twins agree.} Put $L=(w_1,\dots,w_{m-2})$, $L'=(w_2,\dots,w_{m-2})$ and $\psi=w_{m-1}-\Sigma$, so that $w_{m-1}-S_j=\psi+1/v_j$. The two rotations of $Z_j$ above agree exactly when $\cont(v_j,L,\psi+1/v_j)=v_j\cont(L,\psi)$; expanding the left side in its last entry and both sides in the first entry turns this into $\cont(L')/v_j=\cont(L)-\cont(L',\psi)$. The right side does not depend on $j$ and $\cont(L')\neq0$, so all $v_j$ are equal to some $w$, and every twin has self-energy $e=z-w$.

\emph{Transfer matrices.} With equal twins the two paths through two twins have equal products, and the remaining condition is that the $m$ rotations of $Z_j$ agree. Consider the cyclic sequence $u_0=w$, $u_i=w_i$ $(1\le i\le m-1)$ with coupling $\beta_0=r$ between $u_{m-1}$ and $u_0$ and $\beta_i=1$ otherwise. The rotation starting at $x_a$ has
\[
\Gamma=\frac{1}{\cont(w_a,\dots,w_{m-2},w_{m-1}-r\varrho_a)\,\cont(w,w_1,\dots,w_{a-1})},
\]
where $\varrho_a=\cont(w_1,\dots,w_{a-1})/\cont(w,w_1,\dots,w_{a-1})$; by the junction rule this is $1/\cont_a$, where $\cont_a$ is the continuant of the cyclic sequence read from position $a$. The rotation starting at $t_j$ has $\Gamma=1/\cont_t$ with $\cont_t=\cont(w,w_1,\dots,w_{m-2},w_{m-1}-(r-1)/w)$. So $D$ has a set potential exactly when $\cont_1=\dots=\cont_{m-1}=\cont_t$. Let $T_i=\bigl(\begin{smallmatrix}u_i&-\beta_i\\1&0\end{smallmatrix}\bigr)$ and $M_a=T_{a-1}\cdots T_{a+1}T_a$. As in \cref{thm:set-cycles},
\[
M_a=\begin{pmatrix}\cont_a&-\beta_aX_a\\Y_a&-\beta_ac_a\end{pmatrix},
\]
where $X_a,Y_a,c_a$ are the continuants of the cyclic sequence from $a+1$ to $a-1$, from $a$ to $a-2$, and from $a+1$ to $a-2$. The relation $M_{a+1}T_a=T_aM_a$ gives (i) $Y_{a+1}=X_a$, (ii) $\cont_{a+1}=u_aX_a-\beta_ac_a$, (iii) $u_a(\cont_{a+1}-\cont_a)=\beta_{a+1}X_{a+1}-\beta_aY_a$, and all $M_a$ have the same trace $\cont_a-\beta_ac_a$.

\emph{Necessity.} Suppose $\cont_a=\cont_t=:A$ for $1\le a\le m-1$. For these $a$, $\beta_a=1$, so $c_a=A-\mathrm{tr}=:c$, and (ii) gives $u_aX_a=A+c$ for $1\le a\le m-2$. For $1\le a\le m-2$, (iii) and (i) give $X_{a+1}=Y_a=X_{a-1}$, so $X_a=X_{a+2}$ for $0\le a\le m-3$. Expanding $\cont_t$ in its last entry gives $\cont_0=\cont(w,w_1,\dots,w_{m-1})=A+(r-1)Y_0/w$. Then: (iii) at $a=0$ reads $w(A-\cont_0)=X_1-rY_0$, so $X_1=Y_0=X_{m-1}$ by (i); (ii) at $a=m-1$ reads $\cont_0=w_{m-1}X_{m-1}-c$, so $(w_{m-1}-(r-1)/w)X_{m-1}=A+c$; (ii) and the trace at $a=0$ read $A=wX_0-rc_0$ and $rc_0=\cont_0-A+c$, so $wX_0-(r-1)X_{m-1}/w=A+c$; and (iii) at $a=m-1$ reads $w_{m-1}(r-1)X_{m-1}/w=rX_0-X_{m-2}$. For $m$ odd, $X_1=X_{m-1}$ joins the two parity classes, so $X$ is constant, and the equations give $w_a=w-(r-1)/w$ for $1\le a\le m-2$ and $w_{m-1}=w$. For $m$ even, $X_0=X_2=\dots=X_{m-2}$ and $X_1=X_3=\dots=X_{m-1}$; the last equation gives $w_{m-1}X_1=wX_0$, because $r\neq1$; hence $w_a=(A+c)/X_1=w_{m-1}-(r-1)/w$ for odd $a\le m-3$ and $w_a=(A+c)/X_0=w-(r-1)/w_{m-1}$ for even $a$. In terms of self-energies these are the stated conditions.

\emph{Sufficiency.} For $m$ odd put $q=w-(r-1)/w$. Then $T_i=T_q$ for $1\le i\le m-2$, $T_{m-1}=T_w$, and $T_0T_{m-1}=\bigl(\begin{smallmatrix}w^2-r&-w\\w&-1\end{smallmatrix}\bigr)=wT_q-I$ is a polynomial in $T_q$; hence $M_a=T_q^{m-2}(wT_q-I)$ for every $1\le a\le m-1$, and the $\cont_a$ agree. Moreover $\cont_t=\cont(w,q,\dots,q)$ with $m-1$ entries $q$, and expanding gives $\cont_t=(wq-1)\cont(q^{m-2})-w\cont(q^{m-3})$ and $\cont_1=(w^2-r)\cont(q^{m-2})-w\cont(q^{m-3})$, which agree because $wq=w^2-r+1$. For $m$ even put $p=w_{m-1}$, $u=p-(r-1)/w$ and $u'=w-(r-1)/p$, so that $wu=pu'=wp-r+1$. Let $S=T_{u'}T_u$, $S'=T_uT_{u'}$ and $\beta=w/u'$. Then $T_0T_{m-1}=\bigl(\begin{smallmatrix}wp-r&-w\\p&-1\end{smallmatrix}\bigr)=(\beta-1)I+\beta S$ and $T_1T_0T_{m-1}T_{m-2}=(\beta-1)S'+\beta S'^2$. Grouping the factors of $M_a$ in pairs, with $k=(m-2)/2$, gives $M_a=(\beta-1)S^{k}+\beta S^{k+1}$ for odd $a$ and $M_a=(\beta-1)S'^{k}+\beta S'^{k+1}$ for even $a$. Now $[S^{j}]_{11}=\cont((u,u')^{j})$ and $[S'^{j}]_{11}=\cont((u',u)^{j})$ are equal by reversal, so all $\cont_a$ agree. Finally, with $Q=\cont((u,u')^{k})$ and $Q'=\cont((u,u')^{k-1},u)$, expansion gives $\cont_t=\cont(w,(u,u')^{k},u)=(wu-1)Q-wQ'$ and $\cont_1=w\cont((u,u')^{k},p)-rQ=(wp-r)Q-wQ'$, which agree because $wu=wp-r+1$.
\end{proof}

With zero self-energies \cref{thm:one-blown} excludes $C_m(r,1,\dots,1)$ for every $r\ge2$, in accordance with \cref{conj:blowups}.

\subsection{Directed cacti}
A \emph{directed cactus} is a strongly connected digraph in which every block of the underlying graph is a chordless directed cycle of length at least $3$. Two cycles share at most one vertex, and the cycles and cut vertices form a tree. For a cycle $C$ and a vertex $v$ of $C$, every other cycle $C'$ through $v$ determines a \emph{branch} $B$ at $v$ off $C$: the component of $D-v$ that contains $C'-v$. The token enters $B$ only by the arc of $C'$ that leaves $v$, to a vertex $b$, and once $v$ is visited it cannot leave $B$. Write $R_B=R(B,b)$ for the root resolvent of the game on $D[B]$ started at $b$, and
\[
\varepsilon_C(v)=\eta(v)+\sum_{B\text{ a branch at }v\text{ off }C}R_B .
\]

\begin{theorem}[directed cacti]\label{thm:cacti}
A directed cactus with self-energies vanishing at infinity has a set potential if and only if, for every cycle $C$, the function $\varepsilon_C$ has period $2$ around $C$.
\end{theorem}
\begin{proof}
\emph{Vertex sets.} Consecutive vertices of a directed path $P$ are joined by an arc of a unique cycle. $P$ cannot return to a cycle it has left, since it would have to pass again through the cut vertex by which it left. So $P$ uses arcs of cycles $C_1,\dots,C_t$ that form a path in the block tree: it runs along a final segment of $C_1$ ending at the cut vertex $c_1=C_1\cap C_2$, along the segment of $C_s$ from $c_{s-1}$ to $c_s$ for $1<s<t$, and along an initial segment of $C_t$ starting at $c_{t-1}$. The cycles containing two vertices of $X=V(P)$ are exactly $C_1,\dots,C_t$, since another such cycle would close a circuit in the block tree. So a second path $P'$ with $V(P')=X$ uses cycles among $C_1,\dots,C_t$, and for $t\ge2$ it uses all of them, because $X\cap C_1$ contains a vertex other than $c_1$ that lies on no other $C_s$, and likewise for $C_t$; these cycles form a path in the block tree, so $P'$ uses them in the same or in the reverse order. In the same order, the segments of $P'$ are the sets $X\cap C_s$, so $P'=P$ unless $t=1$ and $X=C_1$, when $P,P'$ are rotations of $C_1$. In the reverse order, $P'$ runs through each middle cycle from $c_s$ to $c_{s-1}$; that segment has a different vertex set from the segment from $c_{s-1}$ to $c_s$, because the cycle has length at least $3$, so $t=2$. Then $X\cap C_1$ is both a segment ending at $c_1$ and a segment starting at $c_1$, which forces $X\cap C_1=C_1$, and likewise $X\cap C_2=C_2$. So the vertex sets carried by several paths are the cycles, carried by their rotations, and the unions $C\cup C'$ of two cycles sharing a vertex $v$, carried by $(C-v,v,C'-v)$ and $(C'-v,v,C-v)$, where $C-v$ is traversed from the successor of $v$.

\emph{Rotations.} Along a rotation of $C$ the token stays on $C$, so every branch off $C$ is unvisited while the token stands at its attaching vertex; as in the proof of \cref{lem:branches}, a branch $B$ at $v$ contributes exactly $R_B$ to $1/R(W,v)$. So the rotations of $C$ have the products $\Gamma$ of the rotations of the cycle game on $C$ with self-energies $\varepsilon_C$, which vanish at infinity, and by \cref{thm:set-cycles} they agree if and only if $\varepsilon_C$ has period $2$.

\emph{Two cycles.} Let $C,C'$ share $v$, and let $\sigma,\sigma'$ be the sequences of weights $z-\varepsilon_C$ along $C-v$ and $z-\varepsilon_{C'}$ along $C'-v$, each read from the successor of $v$. Let $\varepsilon^{*}(v)=\eta(v)+\sum R_B$ over the branches at $v$ off both $C$ and $C'$. Along $(C-v,v,C'-v)$ every branch hanging off the current cycle is unvisited when the token reaches it; the branch at $v$ through $C'$ has root resolvent $\cont({}^{-}\sigma')/\cont(\sigma')$, and the junction rule gives
\[
\Gamma(C-v,v,C'-v)=\frac{1}{\cont\bigl(\sigma,\,z-\varepsilon^{*}(v),\,\sigma'\bigr)},\qquad \Gamma(C'-v,v,C-v)=\frac{1}{\cont\bigl(\sigma',\,z-\varepsilon^{*}(v),\,\sigma\bigr)}.
\]
If $\varepsilon_C$ and $\varepsilon_{C'}$ have period $2$, then $\sigma$ and $\sigma'$ are palindromes: constant for an odd cycle, and for an even cycle of length $m$ the $i$-th and the $(m-i)$-th vertices after $v$ have the same parity. Reversing the continuant then shows that the two products agree. So the conditions on rotations imply those on pairs of cycles.
\end{proof}

\begin{corollary}[uniform cacti]\label{cor:uniform-cacti}
Let every cycle of a directed cactus have length $m$, let $e$ vanish at infinity, and let $d(v)$ be the number of cycles through $v$. With $y=z-e$ and $f(y)=\mu(P_{m-2})(y)/\mu(P_{m-1})(y)$, the self-energies $\eta(v)=e-(d(v)-1)f(y)$ give a set potential.
\end{corollary}
\begin{proof}
By induction on the number of vertices of a branch, $R_B=f(y)$ for every branch $B$: $B$ is entered through a cycle $C'$ and runs along $C'-v$; at each vertex $u$ of $C'-v$ the $d(u)-1$ branches off $C'$ are smaller than $B$, so they have resolvent $f(y)$ and reduce the weight $z-\eta(u)=y+(d(u)-1)f(y)$ to $y$; so $B$ plays as the path on $m-1$ vertices of weight $y$, whose root resolvent is $f(y)$. Hence $\varepsilon_C(v)=e$ on every cycle, and \cref{thm:cacti} applies.
\end{proof}

For two triangles this is the bowtie of \cref{thm:cores5}. With $e=0$ a vertex on several cycles has a self-energy with negative leading coefficient, so no set of exits produces it, and no cactus with more than one cycle arises from exits alone.

\subsection{Six and seven vertices}

\begin{proposition}[the census on six and seven vertices]\label{prop:census67}
\begin{enumerate}[label=\textup{(\arabic*)}]
\item Among six-vertex digraphs that are not SCC-symmetric, the back-arc test of \cref{thm:cocycle} leaves $9{,}113$, and exactly $2{,}291$ of them have a set potential. Twelve of these have two non-symmetric strong components, both directed triangles, and the others have one. The pruned cores of the $2{,}303$ non-symmetric components are $C_3$ ($2{,}030$ components), $C_4$ ($210$), $B_2=C_3(2,1,1)$ ($50$), $C_5$ ($9$), and $C_4(2,1,1,1)$, $C_6$, $C_3(2,2,2)$, $C_4(1,2,1,2)$ (one each).
\item Among strongly connected, non-symmetric seven-vertex digraphs the back-arc test leaves $469$, with $107$ pruned cores, $15$ of them blow-ups of cycles. Exactly three of the $469$ have a set potential with zero self-energies: $C_7$, the book $C_3(3,1,1)$ with two pendant vertices at its spine tail, and $C_4(2,1,1,1)$ with pendant vertices at $x_1$ and $x_2$; the pendant vertices supply, through \cref{lem:branches}, the self-energies $(r-1)/z$ of \cref{thm:one-blown}. On six vertices the same test finds $C_6$, $C_3(2,2,2)$ and $C_4(1,2,1,2)$ among the $24$ pruned cores.
\end{enumerate}
\end{proposition}
\begin{proof}[Method \textup{(computer-assisted)}]
Exact tests over $\Qz$ on all digraphs with six and seven vertices, generated with \code{nauty} (\code{geng}, \code{directg}) and filtered by the back-arc test (\code{verification/legacy}).
\end{proof}

Every pruned core in the census is a blow-up of a directed cycle, although the $24$ six-vertex pruned cores include $15$ that are not blow-ups, one of them a cactus.

\begin{conjecture}[blow-ups]\label{conj:blowups}
\begin{enumerate}[label=\textup{(\arabic*)}]
\item With zero self-energies, $C_m(s_0,\dots,s_{m-1})$ has a set potential only if $s_{i+2}=s_i$ for all $i$.
\item If a digraph $D$ has a set potential, then for every non-symmetric strong component $C$ of $D$ the pruned core of $C$, which carries the self-energies induced by the exits of $C$ (\cref{lem:set-components}) and by its pruned branches (\cref{lem:branches}), is a blow-up of a directed cycle.
\end{enumerate}
\end{conjecture}

The converse of (1) is \cref{cor:two-periodic}, so (1) asks for the ``only if'' half of the earlier Conjecture~11.27(1). It holds for one blown-up vertex (\cref{thm:one-blown}), and by \cref{thm:blowup-criterion} it is a statement about finitely many continuant identities for each size vector: the criterion confirms it on all $331$ size vectors with $3\le m\le8$ and sizes at most $4,4,3,3,2,2$ for $m=3,\dots,8$ (up to rotation). Part (2) holds on six vertices by \cref{prop:census67}, and it fails for arbitrary self-energies, because cacti are pruned cores with set potentials that are not blow-ups (\cref{cor:uniform-cacti}).

\section{Potentials versus bounded width}\label{sec:treewidth}

DVG takes linear time on digraphs of bounded treewidth \cite{Bodlaender}, and directed cycles have treewidth $2$; so Bodlaender's dynamic program does not come from a local weighted matching potential (\cref{thm:conj610}). This section asks what kind of potential a width-bounded dynamic program could carry, and answers with three barriers and one construction. Exact resolvents have exponential degree already at pathwidth $2$, and every exact multiplicative potential inherits at least half of that degree, however it is indexed (\cref{thm:intrinsic-degree}). The state on which the resolvent recursion closes is the canonical subgame, which is non-reversing and path-dependent, and at treewidth $3$ one vertex can carry $2^{\Omega(\sqrt{n/\log n})}$ distinct exact resolvents, so no bag-local state determines them (\cref{prop:memory}). What survives is a change of representation: resolvents as arithmetic circuits over component states have linear size when strong components are small (\cref{thm:component-circuits}), and the part of the spectral data that a bag carries is the valuation, the level at which Bodlaender's algorithm works.

\subsection{Degree}

\begin{theorem}[exponential degree on ladders]\label{thm:ladders}
Let $L_m$ have vertices $t_0,\dots,t_{m-1}$ and $b_0,\dots,b_{m-1}$, arcs $t_i\to t_{i+1}$ and $b_i\to b_{i+1}$, and rungs $t_i\rightleftarrows b_i$. Then $L_m$ is SCC-symmetric, its underlying graph has pathwidth $2$, and $R(V,t_0)$ in lowest terms has a denominator of degree $2^{m+1}-2$.
\end{theorem}
\begin{proof}
The strong components are the rungs. Let $F_i$ be the resolvent when the token enters column $i$ with columns $i,\dots,m-1$ unvisited; a vertex left behind in an earlier column can never be reached again, because the rails point right. From $t_i$ the token moves to $t_{i+1}$, worth $F_{i+1}$, or to $b_i$, from which the only move is to $b_{i+1}$; by the symmetry of the rails
\[
F_i=\cfrac{1}{z-F_{i+1}-\cfrac{1}{z-F_{i+1}}}=\frac{z-F_{i+1}}{(z-F_{i+1})^{2}-1},\qquad F_{m-1}=\frac{z}{z^{2}-1}.
\]
If $F_{i+1}=p/q$ in lowest terms with $\deg p<\deg q=d$, then $F_i=q(zq-p)/\bigl((zq-p-q)(zq-p+q)\bigr)$. A common irreducible factor of numerator and denominator would divide $q$ or $zq-p$, and in either case also the other one and then $p$; so this is in lowest terms, with denominator degree $2d+2$ and numerator degree $2d+1$. Hence $d_i+2=2(d_{i+1}+2)$, and $d_{m-1}=2$ gives $d_0=2^{m+1}-2$.
\end{proof}

The weights of \cref{thm:directed-collapse} inherit this: the weight at $t_0$ is $z-F_1$, whose denominator has degree $2^{m}-2$. The next theorem shows that no other indexing avoids it.

\begin{theorem}[the degree is intrinsic]\label{thm:intrinsic-degree}
Let $\mathcal S$ be any set of states, $\Mpot\colon\mathcal S\to\Qz^{\times}$ any map, and suppose that to the reachable position $(V,t_0)$ of $L_m$ are attached two states $s,s'$ with $R(V,t_0)=\Mpot(s')/\Mpot(s)$ \textup{(}for instance, states that record the unvisited set, a bag of a tree decomposition with the part of the play inside it, or the entire history\textup{)}. Then $\Mpot(s)$ or $\Mpot(s')$ has a numerator or denominator of degree at least $2^{m}-1$.
\end{theorem}
\begin{proof}
Write $\Mpot(s)=p_1/q_1$ and $\Mpot(s')=p_2/q_2$ with all four degrees at most $d$. Then $R(V,t_0)=p_2q_1/(q_2p_1)$ has, in lowest terms, a denominator of degree at most $2d$, so $2d\ge2^{m+1}-2$ by \cref{thm:ladders}.
\end{proof}

So a potential that avoids the degree explosion cannot consist of explicit rational functions, whatever it is indexed by; the escape is a change of representation (\S\ref{ssec:circuits}). Since $L_m$ has $2m$ vertices, the bound is $2^{n/2}-1$.

\subsection{States: the canonical subgame}\label{ssec:canonical}

\begin{definition}\label{def:canonical}
The \emph{canonical subgame} of a position $(W,v)$ is the pair $\operatorname{cs}(W,v)=(\mathrm{Reach}_W(v),v)$, where $\mathrm{Reach}_W(v)$ is the set of vertices reachable from $v$ inside $D[W]$.
\end{definition}

\begin{lemma}[the recursion closes on canonical subgames]\label{lem:canonical}
\begin{enumerate}[label=\textup{(\arabic*)}]
\item $R(W,v)$ depends only on $\operatorname{cs}(W,v)$, and with $S=\mathrm{Reach}_W(v)$,
\[
\frac1{R(W,v)}=z-\eta_v-\sum_{w\in\nplus(v)\cap(S-v)}R\bigl(\operatorname{cs}(S-v,w)\bigr).
\]
\item The canonical subgame is non-reversing, $\mathrm{Reach}_{W-v}(w)\subseteq\mathrm{Reach}_W(v)\setminus\{v\}$ for every option $(W-v,w)$, and path-dependent: it is a function of the history of play through $W$, not of the token alone.
\item For a position $(W,v)$ with $v$ in the strong component $C$, $\operatorname{cs}(W,v)$ is determined by $W\cap C$, $v$ and the set of vertices below $C$, which are all unvisited \textup{(}\cref{lem:entry}\textup{)}; so on reachable positions the component state $(W\cap C,v)$ determines the canonical subgame.
\end{enumerate}
\end{lemma}
\begin{proof}
The game tree of $(W,v)$ is the tree of directed paths from $v$ inside $D[W]$, and every such path stays inside $\mathrm{Reach}_W(v)$; the recursion is \eqref{eq:schur} with self-energies, restricted to this set. (2) is immediate, and (3) is \cref{lem:entry}: every vertex reachable from $v$ outside $C$ lies below $C$ and is unvisited, and inside $C$ the reachable set is computed in $D[W\cap C]$ together with the exits.
\end{proof}

A potential with states finer than canonical subgames records irrelevant history, and any exact recursion must distinguish canonical subgames with different resolvents. The next proposition shows how many there can be at one vertex, at treewidth $3$.

\begin{proposition}[memory]\label{prop:memory}
For $m\ge1$ let $Q_m$ be the digraph with a spine $p_0,\dots,p_m$, vertices $a_i,b_i$ $(0\le i<m)$ with arcs $p_i\to a_i\to p_{i+1}$ and $p_i\to b_i\to p_{i+1}$, a collector $q$ with arcs $p_m\to q$ and $q\to a_i$, $q\to b_i$, and at every $x\in\{a_i,b_i\}$ a private directed tail, so that the directed path from $x$ through its tail has $L(x)$ vertices, where the $2m$ numbers $L(x)+1$ are distinct primes. Then:
\begin{enumerate}[label=\textup{(\arabic*)}]
\item $Q_m$ has treewidth at most $3$, and $n=|V(Q_m)|=O(m^{2}\log m)$ when the primes are the first $2m$ odd primes;
\item the positions at $q$ reachable from $(V,p_0)$ have $2^{m}$ pairwise distinct resolvents;
\item consequently, for any set $B\ni q$ of at most $k+1$ vertices with $k<m$, some $2^{m-k}\ge2$ of these positions agree on $W\cap B$ and have pairwise distinct resolvents; no function of the token and of $W\cap B$ equals $R$, and an exact table of resolvents indexed by any states needs at least $2^{m}=2^{\Omega(\sqrt{n/\log n})}$ states at $q$.
\end{enumerate}
The outcome at every such position is the same \textup{(}the player to move at $q$ loses\textup{)}.
\end{proposition}
\begin{proof}
(1) $Q_m-q$ without its tails is a chain of diamonds, a series-parallel graph of treewidth $2$; adding $q$ to every bag of a width-$2$ decomposition, and attaching the tails as paths, gives width at most $3$. The tails have $\sum_x(L(x)-1)=O(m^{2}\log m)$ vertices by the prime number theorem.

(2) A play from $p_0$ reaches $q$ only as $p_0c_0p_1c_1\cdots c_{m-1}p_mq$ with $c_i\in\{a_i,b_i\}$, since entering a tail ends the play in it. At the resulting position $(W_c,q)$ the options are the unvisited $x_i\in\{a_i,b_i\}\setminus\{c_i\}$, and from $x_i$ the only continuation is its tail, because $p_{i+1}$ is visited; so the subgame from $x_i$ is a directed path with $L(x_i)$ vertices and
\[
R(W_c,q)=\frac{1}{z-S_c},\qquad S_c=\sum_{i<m}f_{L(x_i)},\qquad f_L=\frac{\mu(P_{L-1})}{\mu(P_L)} .
\]
The poles of $f_L$ are the zeros $2\cos(k\pi/(L+1))$, $1\le k\le L$, of $\mu(P_L)=U_L(z/2)$; they are simple, and they are not zeros of the coprime numerator $\mu(P_{L-1})$. For distinct primes $p=L+1$ and $p'=L'+1$ the pole sets are disjoint, since $k/p=k'/p'$ with $0<k<p$, $0<k'<p'$ is impossible. So the poles of $S_c$ are exactly the union of the pole sets of the chosen $f_{L(x_i)}$, which determines $c$; distinct $c$ give distinct $S_c$ and distinct $R(W_c,q)$.

(3) Since $q\in W_c$ always, $W_c\cap B$ takes at most $2^{k}$ values, so some value is shared by at least $2^{m-k}$ vectors $c$. The number of vertices is $n=O(m^{2}\log m)$, so $m=\Omega(\sqrt{n/\log n})$. Finally all $L(x)$ are even, since $L(x)+1$ is an odd prime, so every move from $q$ leads to a directed path with an even number of vertices, which is an N-position.
\end{proof}

The package reproduces (2) for $m\le6$ (\code{treewidth.comb}, \code{tests/test_treewidth.py}). So the ``path-dependent bag invariant'' that an exact potential on bounded treewidth would need is not bag-local: at treewidth $3$ it must remember $\Omega(\sqrt{n/\log n})$ bits of the history outside any bag, while the outcomes, which are what Bodlaender's algorithm computes, are constant on the same positions.

\subsection{Circuits: the component-state potential}\label{ssec:circuits}

\begin{theorem}[component-state circuits]\label{thm:component-circuits}
For every digraph $D$ with self-energies from exits there is an arithmetic circuit over $\QQ$ with input $z$ and gates $+,-,\times,\div$, with
\[
O\Bigl(\sum_{C}2^{|C|}\bigl(|C|+e(C)\bigr)\Bigr)
\]
gates, where $C$ runs over the strong components and $e(C)$ counts the arcs with tail in $C$, whose gates include $R(W,v)$ for every reachable position $(W,v)$. In particular it has $O(2^{k}(n+|A(D)|))$ gates when every strong component has at most $k$ vertices; for the ladders $L_m$ it has $O(m)$ gates and computes $R(V,t_0)$, of degree $2^{m+1}-2$.
\end{theorem}
\begin{proof}
Process the components in reverse topological order. For a component $C$ put $\eta_x=\sum\Rdown(y)$ over the exits $y$ of $x$, gates already built, and for $S\subseteq C$ and $v\in S$, in order of increasing $|S|$, the gate
\[
R_C(S,v)=\Bigl(z-\eta_v-\sum_{w\in\nplus(v)\cap(S-v)}R_C(S-v,w)\Bigr)^{-1}.
\]
Then $\Rdown(y)=R_C(C,y)$ for $y\in C$. By \cref{lem:entry,lem:canonical}(3), $R(W,v)=R_C(W\cap C,v)$ at every reachable position with $v\in C$. Each gate costs $O(1+|\nplus(v)|)$ operations, which gives the bound; for $L_m$ every component is a rung with two vertices.
\end{proof}

\begin{corollary}\label{cor:small-components}
DVG is decidable in time $2^{k}\poly(n)$, where $k$ is the largest number of vertices of a strong component that is not symmetric; in particular in polynomial time when every non-symmetric strong component has $O(\log n)$ vertices.
\end{corollary}
\begin{proof}
Symmetric components are decided by \cref{cor:dvg-poly} from the outcomes of their exits; a non-symmetric component $C$ is decided by the outcome version of the recursion of \cref{thm:component-circuits} over the $|C|2^{|C|}$ states $(S,v)$, exits to P-positions winning at once.
\end{proof}

Neither bound uses treewidth, and neither is small on long directed cycles with many chords. Canonical subgames can do better than component states: on the closed ladder, $L_m$ plus the arc $b_{m-1}\to t_0$, which is strongly connected and not symmetric and has pathwidth at most $3$, the recursion of \cref{lem:canonical} has few states although the component-state bound is $2m\cdot2^{2m}$ \textup{(}\code{treewidth.closed_ladder}, \code{canonical_subgame}\textup{)}:
\begin{center}\small
\begin{tabular}{@{}lrrrrrr@{}}
\toprule
$m$ & $2$ & $3$ & $4$ & $5$ & $6$ & $7$\\ \midrule
reachable positions & $22$ & $69$ & $188$ & $463$ & $1{,}092$ & $2{,}495$\\
distinct canonical subgames & $16$ & $48$ & $108$ & $212$ & $376$ & $616$\\
distinct exact resolvents & $9$ & $20$ & $46$ & $86$ & $150$ & $246$\\
\bottomrule
\end{tabular}
\end{center}
The subgame counts equal $(8m^{3}-30m^{2}+94m-84)/3$ for $2\le m\le7$; we have not proved cubic growth in general.

\begin{remark}[the bridge to Bodlaender's algorithm]\label{rem:bodlaender}
Combining \cref{thm:intrinsic-degree,prop:memory,thm:component-circuits}: an exact potential that a bounded-width dynamic program could carry must be (i) path-dependent, indexed at least by canonical subgames, which are not bag-local, and (ii) represented by circuits rather than by explicit rational functions. The part of the spectral data that is bag-local is the valuation $\ord_0R$, which is $\pm1$ by \cref{thm:resolvent-outcome}; that is the information Bodlaender's algorithm propagates, and on the comb of \cref{prop:memory} it is constant where the exact resolvents take $2^{m}$ values. Valuation potentials indexed by unvisited sets need not exist either (\cref{thm:hierarchy}: a three-vertex digraph has none), so even at the valuation level the index must record more than the unvisited set. Whether exact resolvents on bounded treewidth have circuits of polynomial size is \cref{prob:tw-circuits}.
\end{remark}

\section{Status and open problems}\label{sec:open}

The two conjectures that organized the earlier drafts are settled: \cref{thm:conj610} (formerly Conjecture~6.10) and \cref{thm:frozen-real} (formerly Conjecture~10.9$'$) hold for real weights and fail for complex ones (\cref{cor:complex-false}; \cref{thm:every-cycle}). Open Problem~7 of the first draft is settled for one arc and for common tails (\cref{cor:one-arc,thm:common-tail}) and reduced for $p$ tails to undirected phases (\cref{lem:p-tails}). The treewidth question is answered by three barriers and a circuit construction (\S\ref{sec:treewidth}). What remains:

\begin{problem}[Maker--Maker games of rank $3$]\label{prob:makermaker}
Are Maker--Maker games on hypergraphs of rank $3$ polynomial or $\PSPACE$-complete? Galliot lists this case as open \cite{Galliot25}.
\end{problem}

\begin{problem}[below $\Pi_2^{\Pclass}$]\label{prob:qcsp}
Classify $\mathrm{QCSP}(\Gamma)$ below $\Pi_2^{\Pclass}$ for every finite $\Gamma$. The examples of Zhuk and Martin show that $\mathsf{DP}$ and $\Theta_2^{\Pclass}$ must appear in the answer \cite{ZhukMartin}.
\end{problem}

\begin{problem}[Grundy values and two tokens]\label{prob:grundy}
Are Grundy values of undirected vertex geography fixed-parameter tractable in treewidth, and is two-token undirected geography on cubic graphs polynomial? (Open Questions~7 and~8 of \cite{BFT}.) Is there a natural partizan ruleset whose single instances are polynomial but whose sums of two instances are $\PSPACE$-hard? (Open Question~9 of \cite{BFT}.)
\end{problem}

\begin{problem}[real-number evaluation]\label{prob:ch}
Separate $\PSPACE$ from the counting hierarchy. By \cref{thm:bss} this would rule out polynomial-time real-number evaluation of $\val$ with algebraic constants.
\end{problem}

\begin{problem}[two tails]\label{prob:two-tails}
Is directed vertex geography polynomial, or fixed-parameter tractable in the number of distinct tails of one-way arcs per strong component, already when every strong component has two tails? By \cref{lem:p-tails} the question is undirected geography with two matching-type terminals, and by \cref{prop:barrier,prop:ge-barrier} no rule that reads only the matching structure at the key points decides it. For fixed $k\ge2$ arcs per component, is DVG polynomial, fixed-parameter tractable in $k$, or $\PSPACE$-complete?
\end{problem}

\begin{problem}[counting maximum matchings]\label{prob:fpt-deficiency}
On planar graphs of deficiency $d$, maximum matchings, and hence the residues of \cref{prop:residue}, can be counted in time $n^{O(d)}$ (\cref{thm:kasteleyn}). Is the count fixed-parameter tractable in $d$?
\end{problem}

\begin{problem}[collapsibility at polynomial horizon]\label{prob:sigma2}
For generators with polynomial horizon, is the set of collapsible generators $\Sigma^0_2$-complete when $\Pclass\neq\PSPACE$? The proof of \cref{thm:sigma2} needs a language that is hard at almost every input length, which $\Pclass\neq\PSPACE$ alone does not provide.
\end{problem}

\begin{problem}[the band]\label{prob:band}
\begin{enumerate}[label=\textup{(\alph*)}]
\item For a digraph $D$ that is not SCC-symmetric, give an explicit $z_D$ such that the path data of $D$ admit no real-stable completion at any real $z>z_D$; \cref{prop:triangle-band} gives $z_D=2$ for the directed triangle. Is $z_D=2$ for every directed cycle?
\item Determine the real locus of the family of \cref{thm:every-cycle} inside $|z|<2$; numerically it is all of $(-2,2)$ for $m=3$, a subinterval for $m=4,5$, and apparently empty for $m\ge6$.
\end{enumerate}
\end{problem}

\begin{problem}[exact resolvents at bounded width]\label{prob:tw-circuits}
Does every digraph of treewidth at most $k$ have an arithmetic circuit of size $f(k)\poly(n)$ that computes $R(V,s)$ for all $s$? By \cref{thm:intrinsic-degree} such a circuit cannot output explicit rational functions of polynomial size, and by \cref{prop:memory} it cannot tabulate resolvents by bag-local states; the comb $Q_m$ is the first test case. Is the number of canonical subgames of the closed ladder cubic in $m$?
\end{problem}

\begin{problem}[set and valuation potentials]\label{prob:setpot}
Which digraphs carry set potentials, and which carry valuation potentials? \Cref{thm:cocycle} gives a necessary condition, \cref{lem:set-components,lem:branches} reduce set potentials to pruned cores, and \cref{conj:blowups} predicts the pruned cores that occur. On five vertices only $79$ of the $8{,}055$ non-SCC-symmetric digraphs carry a set potential.
\end{problem}

Together with \cref{prob:complex} (complex potentials, in particular on strongly connected digraphs) and \cref{conj:minimal-env,conj:blowups}, these are the open questions on the boundaries drawn in this paper.

\appendix

\section{Computational verification}\label{app:verification}

Every computation is exact (rational arithmetic, or Gr\"obner bases over $\Qz$) unless marked as numerical. Evaluation at a rational point is used only to refute an identity. Scripts are in the accompanying archive: \code{src/dvgcollapse} (library), \code{tests} (\code{pytest}), \code{scripts} (one-off computations of this revision), \code{verification/legacy} (the verification package of the earlier drafts, unchanged). Results marked \emph{proof} in the paper do not depend on these checks; results marked \emph{computer-assisted} do.

\begin{center}\small
\begin{longtable}{@{}L{0.25}L{0.43}L{0.25}@{}}
\toprule
Claim & Check & Result\\ \midrule\endhead
\multicolumn{3}{@{}l}{\emph{Earlier drafts} (\code{verification/legacy}; reproduced for this revision)}\\
\cref{cor:index} & $600$ random game trees, random root player and payoffs; orders cross-checked by nullities on $200$ & no failures\\
\cref{thm:resolvent-outcome} & $300$ random rooted trees: outcome, order at $0$, image test, coverage & no failures\\
\cref{thm:matching-collapse} & $200$ random graphs: outcome, path-tree resolvent at random rational $z$ & no failures\\
\cref{thm:directed-collapse} & $22{,}187$ positions in $120$ random SCC-symmetric digraphs, up to $16$ vertices & no failures\\
\cref{prop:no-cancel}, \cref{cor:dvg-poly} & $150$ SCC-symmetric digraphs, all starts & no failures\\
\cref{cor:edge-geo} & $1{,}155$ starts and random positions & no failures\\
\cref{prop:residue} & $124$ P-positions & no failures\\
\cref{lem:two-point} & indeterminate weights, $1{,}680$ vertex pairs, graphs on $2$--$7$ vertices & identity holds\\
\cref{thm:exact5} & $16\times8$, $218\times64$ pairs $(D,U)$; five vertices: set-potential filter, then $80{,}798$ saturated ideals & exactly the SCC-symmetric digraphs have potentials\\
\cref{thm:hierarchy} (census) & all digraphs on $3$--$5$ vertices & table of \S\ref{ssec:hierarchy}\\
\cref{cor:one-arc} & $272{,}399$ positions against exhaustive search & no failures\\
\cref{prop:two-arcs,prop:barrier,prop:ge-barrier} & explicit examples; samples of $26{,}633$ and $38{,}805$ starts & as stated\\
\cref{prop:env-barrier}, \cref{thm:triangle-solved,thm:every-cycle} & Gr\"obner bases and factorization; symbolic identities & as stated\\
\cref{prop:smallest-env,prop:c5c6} & $276$, $4{,}496$, $6{,}560$, $216{,}288$; $9{,}168$, $14$, $724$, $59{,}976$ graphs $U$ & all unit ideals\\
\cref{prop:symmetric-env} & resultants for $m\le10$ & nonzero Laurent polynomials\\
\cref{thm:cores5}, \cref{cor:census5}, \cref{prop:census67} & self-energies as indeterminates; all digraphs on $\le7$ vertices via \code{nauty} & as stated\\
\midrule
\multicolumn{3}{@{}l}{\emph{This revision} (\code{src}, \code{tests}, \code{scripts})}\\
\cref{lem:first-order}, \cref{thm:main-real} & \code{real_stability.py}: first-order factorization with indeterminate $\ell_X$; $\kappa(\emptyset)=0$, $\kappa(R)=1$ on random digraphs with one-way arcs on cycles; \code{obstruction_certificate} returns a real point with negative Rayleigh difference & as stated\\
\cref{prop:triangle-band} & \code{triangle_band_value}: $\Delta_{20}G(-\tfrac12)$ symbolic & $(4-z^2)/(4(z^2-2)^2)$\\
\cref{cor:c4} & \code{scripts/check_c4.py}: the three Rayleigh expansions & as stated\\
numerical sanity check (not used in proofs) & \code{scripts/adversarial_*.py}: differential evolution searching real-stable completions of cycle path data at large $z$; a symmetric control digraph, where a completion exists, is found & no completion found for cycles; control succeeds\\
\cref{thm:complex-counterexample,prop:sinks} & \code{complex_counterexample.verify_symbolic}: all reachable positions, resolvents by the game recursion, $\mu_a$ in symmetric functions, for $(m,r)\in\{(3,0),(4,0),(5,0),(3,1),(3,2),(4,1)\}$ & all identities exact\\
same, numerical & \code{verify_numeric}: roots of $Y_m$ to $50$ digits, $m\le6$, real and complex $z$ & residuals $<10^{-50}$\\
\cref{prop:Ym-arithmetic}(3) & \code{scripts/ym_arithmetic.py}: irreducibility, discriminants, PARI \code{polgalois} at $z=3,5,7$ & $S_4,S_5,S_6,S_7$\\
\cref{prop:c3-minimal} & \code{scripts/search6_c3k3.py}: $1{,}928$ graphs $U$, saturated over $\Qz$ (Singular) & all unit ideals ($878$\,s)\\
\cref{thm:common-tail}, \cref{cor:star-dvg} & \code{tests/test_star_tail.py}: every position of random instances against exhaustive search; the DVG solver on $60$ random digraphs & no failures\\
\cref{thm:blowup-criterion}, \cref{lem:twin}, \cref{cor:two-periodic} & \code{tests/test_blowups.py}: formula against all paths of seven blow-ups; criterion against direct test on all $115$ blow-ups with $\le9$ vertices; twin identity for $m\le10$ & no failures; exactly the $13$ two-periodic ones pass\\
\cref{conj:blowups}(1) & \code{scripts/blowup_scan.py}: $331$ size vectors & no counterexample\\
\cref{thm:ladders,thm:component-circuits} & \code{tests/test_treewidth.py} & as stated\\
\cref{prop:memory} & \code{treewidth.comb}, $m\le6$ ($n\le232$) & $2^{m}$ distinct resolvents at $q$\\
closed ladders (\S\ref{ssec:circuits}) & $m\le7$ & table of \S\ref{ssec:circuits}\\
\bottomrule
\end{longtable}
\end{center}

\section{Notation}\label{app:notation}

\begin{center}\small
\begin{longtable}{@{}L{0.25}L{0.71}@{}}
\toprule
Symbol & Meaning\\ \midrule\endhead
$\arena$, $\SG$, $\val$, $\pay$ & arena, state graph, value, payoff (\S\ref{sec:setting})\\
$N$, $n$ & dimension of the cube; input size (\S\S\ref{sec:setting}--\ref{sec:extensions}, \ref{sec:boundary}) or number of vertices (from \S\ref{sec:spectral} on)\\
$D$, $V=V(D)$, $\nplus(v)$ & digraph, its vertex set, out-neighbourhood\\
$(W,v)$ & position of DVG: unvisited set $W$, token $v\in W$\\
$R(W,v)$, $\Rdown(w)$ & root resolvent of a position; fresh-entry resolvent $R(V,w)$\\
$\eta_v$ & self-energy; $1/R(W,v)=z-\eta_v-\sum R(W-v,w)$\\
$\mu(G)$, $\mu_a(G)$, $\mu(P_\ell)$ & matching polynomial; weighted by vertex weights $a$; of the path with $\ell$ vertices\\
$(U,a)$ & local weighted matching potential: graph $U$ on $V$, weights $a$\\
$\Mpot$ & potential values: $\mu_a(U[\cdot])$, a set potential, or a real-stable potential (\S\ref{sec:realstab})\\
$\Mpot_F$ & frozen potential with environment $F$ (\cref{def:frozen})\\
$c(X)$, $\gamma_X$ & $\Mpot(V\setminus X)/\Mpot(V)$; $z^{|X|}c(X)$\\
$\Gamma(P)$, $\back(P)$, $\ell(P)$ & product of resolvents along a play; back arcs; first-order coefficient (\cref{lem:asymptotics})\\
$\rho(S,s)$ & Green's function $\Mpot(S-s)/\Mpot(S)$ (\S\ref{sec:fields})\\
$\Delta^{+}(v)$, $\Delta^{-}(v)$ & exits and phantom neighbours of $v$\\
$\Delta_{ij}P$ & Rayleigh difference\\
$\Puis$ & real Puiseux series in $z^{-1}$ convergent at $\infty$; $\CC\{\!\{z^{-1}\}\!\}$ its complex analogue\\
$\mathbb K$ & a field containing $\Qz$\\
$\deg f$ & exponent of the leading term at $z=\infty$\\
$\kappa(S)$ & first-order coefficient combination in \cref{lem:first-order}\\
$e_j(a_F)$ & elementary symmetric functions of the environment weights (\S\ref{sec:complex})\\
$Y_m$, $\exp_m(w)$ & polynomial of the environment weights \eqref{eq:Ym}; truncated exponential $\sum_{k\le m}w^k/k!$\\
$\game(H;\arcs)$, $K$, $W'$ & DVG inside a component with one-way arcs $\arcs$; current graph; live heads (\S\ref{sec:oneway})\\
$\nu$ & matching number\\
$\cont(\cdot)$ & continuant (\S\ref{ssec:setpot-cycles})\\
$C_m(s_0,\dots,s_{m-1})$, $b^{(j)}_k$ & blow-up of a directed cycle; branching numbers (\cref{thm:blowup-criterion})\\
$L_m$, $Q_m$ & ladder; comb (\S\ref{sec:treewidth})\\
$\operatorname{cs}(W,v)$ & canonical subgame (\cref{def:canonical})\\
$\ord_0$, $\ord_{z=0}$ & order of vanishing at $z=0$\\
$\mm(G)$, $\mathrm{pm}(G)$ & numbers of maximum and of perfect matchings\\
\bottomrule
\end{longtable}
\end{center}

\section{Concordance with the earlier drafts}\label{app:concordance}

Numbers in \S\S\ref{sec:setting}--\ref{sec:boundary} are unchanged, except that Conjecture~6.10 is now \cref{thm:conj610}. The addendum's \S\S10--11 map as follows.

\begin{center}\small
\begin{longtable}{@{}L{0.29}L{0.33}L{0.31}@{}}
\toprule
Earlier draft & This version & Change\\ \midrule\endhead
Conjecture 6.10 & \cref{thm:conj610}, \cref{cor:complex-false} & proved over $\Qz$, $\Rz$, $\Puis$; false over $\Qzbar$\\
Open Problem 3 (\S9) & \cref{thm:conj610}, \cref{prob:complex} & settled, with a new complex problem\\
Open Problem 7 (\S9) & \cref{cor:one-arc,thm:common-tail}, \cref{prob:two-tails} & settled for one arc and for common tails\\
Lemma 10.A & \cref{lem:asymptotics} & with self-energies\\
Lemma 10.B, 10.C & \cref{lem:cavity,lem:two-point} & \\
Theorem 10.1, 10.2 & \cref{thm:normal,thm:no-phantom} & 10.1 extended to complex Puiseux weights\\
Theorem 10.3, Cor.\ 10.3$'$ & \cref{prop:path-lemma,cor:u-edges} & \\
Theorem 10.4, 10.5 & \cref{cor:unilateral,cor:sign-coherent} & now also consequences of \cref{thm:main-real}\\
Theorem 10.6, 10.7 & \cref{thm:cocycle,thm:locality} & \\
Theorem 10.7$'$, Prop.\ 10.7$''$ & \cref{cor:triangles,cor:c4} & 10.7$''$ was misnamed a theorem in places\\
Theorem 10.8 & \cref{thm:exact5} & \\
Theorem 10.9 & \cref{prop:shape} & restated over arbitrary fields\\
Conjecture 10.9$'$ & \cref{thm:frozen-real} & proved for real weights; false over $\Qzbar$ (\cref{thm:every-cycle})\\
Proposition 10.10 & \cref{prop:env-barrier} & \\
Theorem 10.11, Cor.\ 10.11$'$ & \cref{cor:one-arc}, \cref{cor:star-dvg} & special case of \cref{thm:common-tail}\\
Proposition 10.12 & \cref{prop:two-arcs} & \\
Theorem 10.13 & \cref{thm:hierarchy} & \\
Theorem 10.14 & \cref{thm:ladders} & with \cref{thm:intrinsic-degree}\\
Theorem 11.1, Cor.\ 11.1$'$ & \cref{thm:common-tail}, \cref{cor:star-dvg} & \\
Proposition 11.2, Cor.\ 11.3 & \cref{prop:barrier}, \cref{cor:dichotomy} & \\
Theorem 11.4, 11.5 & \cref{thm:triangle-solved,thm:every-cycle} & \\
Proposition 11.6 & \cref{prop:smallest-env} & \\
Lemma 11.8, 11.9 & \cref{lem:set-components,lem:branches} & 11.9: path kinds defined disjointly\\
Theorem 11.10, 11.11 & \cref{thm:set-cycles,thm:books} & \\
Theorem 11.12, Cor.\ 11.13 & \cref{thm:cores5}, \cref{cor:census5} & \\
Lemma 11.14 & \cref{lem:p-tails} & generalized to $p$ tails\\
Proposition 11.15 & \cref{prop:ge-barrier} & \\
Lemma 11.16, 11.17 & \cref{lem:frozen-path,lem:env-seen} & 11.17 with self-energies\\
Theorem 11.18 & \cref{thm:meanfield} & \\
Proposition 11.19, 11.20 & \cref{prop:symmetric-env,prop:c5c6} & \\
Conjecture 11.21 & \cref{conj:minimal-env} & \\
Theorem 11.22 & \cref{thm:one-blown} & \\
Theorem 11.23 & \cref{cor:two-periodic} & extended to all two-periodic blow-ups\\
Theorem 11.24, Cor.\ 11.25 & \cref{thm:cacti}, \cref{cor:uniform-cacti} & \\
Proposition 11.26 & \cref{prop:census67} & \\
Conjecture 11.27 & \cref{conj:blowups} & ``if'' half of (1) proved (\cref{cor:two-periodic})\\
\S11.11 items 1--5 & \cref{prob:two-tails}, \cref{thm:frozen-real}, \cref{conj:minimal-env}, \cref{conj:blowups}, \cref{prob:band} & item 2 settled\\
\bottomrule
\end{longtable}
\end{center}

\section*{Acknowledgments}
The author acts as the AI Research Architect, Lead Strategist, and Principal Investigator for this project, formulating objectives, defining priority research areas, and directing the programmatic verification trajectories. Generative artificial intelligence model, specifically Anthropic's Claude Opus 5.5, was orchestrated strictly under the author's prompt-engineering protocols, serving as high-performance formalization and computational assistants to accelerate algebraic derivations, draft technical expositions, and assemble the comprehensive search scripts. Every theoretical breakthrough, including the deterministic distinct-tails paramatrization, the cyclic Chebyshev stabilization, and the classification of pruned cores, was rigorously cross-checked, audited, and verified via exact symbolic computations (saturated Gr{\"o}bner bases via Singular) within the accompanying \code{limits-of-collapse} open-source repository.

\begin{thebibliography}{99}\small

\bibitem{ABKM} E.~Allender, P.~B\"urgisser, J.~Kjeldgaard-Pedersen, and P.~B.~Miltersen, \emph{On the complexity of numerical analysis}, SIAM J. Comput., 38 (2009), pp.~1987--2006.

\bibitem{ACRSZ} A.~Ambainis, A.~M.~Childs, B.~W.~Reichardt, R.~\v{S}palek, and S.~Zhang, \emph{Any AND-OR formula of size $N$ can be evaluated in time $N^{1/2+o(1)}$ on a quantum computer}, SIAM J. Comput., 39 (2010), pp.~2513--2530.

\bibitem{APT} B.~Aspvall, M.~F.~Plass, and R.~E.~Tarjan, \emph{A linear-time algorithm for testing the truth of certain quantified Boolean formulas}, Inform. Process. Lett., 8 (1979), pp.~121--123.

\bibitem{AtseriasOliva} A.~Atserias and S.~Oliva, \emph{Bounded-width QBF is PSPACE-complete}, J. Comput. System Sci., 80 (2014), pp.~1415--1429.

\bibitem{BFNW} L.~Babai, L.~Fortnow, N.~Nisan, and A.~Wigderson, \emph{BPP has subexponential time simulations unless EXPTIME has publishable proofs}, Comput. Complexity, 3 (1993), pp.~307--318.

\bibitem{BHN} R.~Bacher, P.~de~la~Harpe, and T.~Nagnibeda, \emph{The lattice of integral flows and the lattice of integral cuts on a finite graph}, Bull. Soc. Math. France, 125 (1997), pp.~167--198.

\bibitem{BakerNorine} M.~Baker and S.~Norine, \emph{Riemann--Roch and Abel--Jacobi theory on a finite graph}, Adv. Math., 215 (2007), pp.~766--788.

\bibitem{BakerShokrieh} M.~Baker and F.~Shokrieh, \emph{Chip-firing games, potential theory on graphs, and spanning trees}, J. Combin. Theory Ser. A, 120 (2013), pp.~164--182.

\bibitem{BarnumSaks} H.~Barnum and M.~Saks, \emph{A lower bound on the quantum query complexity of read-once functions}, J. Comput. System Sci., 69 (2004), pp.~244--258.

\bibitem{BeaverFeigenbaum} D.~Beaver and J.~Feigenbaum, \emph{Hiding instances in multioracle queries}, in STACS 90, Lecture Notes in Comput. Sci. 415, Springer, 1990.

\bibitem{WinningWays} E.~R.~Berlekamp, J.~H.~Conway, and R.~K.~Guy, \emph{Winning Ways for Your Mathematical Plays}, Academic Press, 1982; 2nd ed., A K Peters, 2001--2004.

\bibitem{Bodlaender} H.~L.~Bodlaender, \emph{Complexity of path-forming games}, Theoret. Comput. Sci., 110 (1993), pp.~215--245.

\bibitem{BodwinGrossman} G.~Bodwin and O.~Grossman, \emph{Strategy-stealing is non-constructive}, in ITCS 2020, LIPIcs 151, 2020, pp.~21:1--21:12.

\bibitem{BorceaBranden} J.~Borcea and P.~Br\"and\'en, \emph{Applications of stable polynomials to mixed determinants: Johnson's conjectures, unimodality, and symmetrized Fischer products}, Duke Math. J., 143 (2008), pp.~205--223.

\bibitem{BBCJK} F.~B\"orner, A.~Bulatov, H.~Chen, P.~Jeavons, and A.~Krokhin, \emph{The complexity of constraint satisfaction games and QCSP}, Inform. Comput., 207 (2009), pp.~923--944.

\bibitem{Bouton} C.~L.~Bouton, \emph{Nim, a game with a complete mathematical theory}, Ann. of Math., 3 (1901), pp.~35--39.

\bibitem{Branden} P.~Br\"and\'en, \emph{Polynomials with the half-plane property and matroid theory}, Adv. Math., 216 (2007), pp.~302--320.

\bibitem{BFT} K.~W.~Burke, M.~T.~Ferland, and S.-H.~Teng, \emph{Winning the war by (strategically) losing battles: settling the complexity of Grundy-values in undirected geography}, in FOCS 2021; arXiv:2106.02114.

\bibitem{CKS} A.~K.~Chandra, D.~C.~Kozen, and L.~J.~Stockmeyer, \emph{Alternation}, J. ACM, 28 (1981), pp.~114--133.

\bibitem{ChapuisKoiran} O.~Chapuis and P.~Koiran, \emph{Saturation and stability in the theory of computation over the reals}, Ann. Pure Appl. Logic, 99 (1999), pp.~1--49.

\bibitem{Chen04} H.~Chen, \emph{Quantified constraint satisfaction and bounded treewidth}, in ECAI 2004, IOS Press, 2004.

\bibitem{Chen08} H.~Chen, \emph{The complexity of quantified constraint satisfaction: collapsibility, sink algebras, and the three-element case}, SIAM J. Comput., 37 (2008), pp.~1674--1701.

\bibitem{Chen11} H.~Chen, \emph{Quantified constraint satisfaction and the polynomially generated powers property}, Algebra Universalis, 65 (2011), pp.~213--241.

\bibitem{Chen12} H.~Chen, \emph{Meditations on quantified constraint satisfaction}, in Logic and Program Semantics, Lecture Notes in Comput. Sci. 7230, Springer, 2012.

\bibitem{CCJY} A.~M.~Childs, R.~Cleve, S.~P.~Jordan, and D.~Yonge-Mallo, \emph{Discrete-query quantum algorithm for NAND trees}, Theory Comput., 5 (2009), pp.~119--123.

\bibitem{CKS01} N.~Creignou, S.~Khanna, and M.~Sudan, \emph{Complexity Classifications of Boolean Constraint Satisfaction Problems}, SIAM, 2001.

\bibitem{CvetkovicGutman} D.~M.~Cvetkovi\'c and I.~M.~Gutman, \emph{The algebraic multiplicity of the number zero in the spectrum of a bipartite graph}, Mat. Vesnik (Beograd), 9 (1972), pp.~141--150.

\bibitem{DemaineHearn} E.~D.~Demaine and R.~A.~Hearn, \emph{Playing games with algorithms: algorithmic combinatorial game theory}, in Games of No Chance 3, MSRI Publ. 56, Cambridge Univ. Press, 2009.

\bibitem{EdmondsPartition} J.~Edmonds, \emph{Minimum partition of a matroid into independent subsets}, J. Res. Nat. Bur. Standards Sect. B, 69B (1965), pp.~67--72.

\bibitem{EGO} E.~Eiben, R.~Ganian, and S.~Ordyniak, \emph{Using decomposition-parameters for QBF: mind the prefix!}, J. Comput. System Sci., 110 (2020), pp.~1--21.

\bibitem{EvenTarjan} S.~Even and R.~E.~Tarjan, \emph{A combinatorial problem which is complete in polynomial space}, J. ACM, 23 (1976), pp.~710--719.

\bibitem{FGG} E.~Farhi, J.~Goldstone, and S.~Gutmann, \emph{A quantum algorithm for the Hamiltonian NAND tree}, Theory Comput., 4 (2008), pp.~169--190.

\bibitem{FHP} J.~K.~Fichte, M.~Hecher, and A.~Pfandler, \emph{Lower bounds for QBFs of bounded treewidth}, in LICS 2020, ACM, 2020.

\bibitem{FraenkelBib} A.~S.~Fraenkel, \emph{Combinatorial games: selected bibliography with a succinct gourmet introduction}, Electron. J. Combin., Dynamic Survey DS2.

\bibitem{FraenkelLichtenstein} A.~S.~Fraenkel and D.~Lichtenstein, \emph{Computing a perfect strategy for $n\times n$ chess requires time exponential in $n$}, J. Combin. Theory Ser. A, 31 (1981), pp.~199--214.

\bibitem{FSU} A.~S.~Fraenkel, E.~R.~Scheinerman, and D.~Ullman, \emph{Undirected edge geography}, Theoret. Comput. Sci., 112 (1993), pp.~371--381.

\bibitem{FraenkelSimonson} A.~S.~Fraenkel and S.~Simonson, \emph{Geography}, Theoret. Comput. Sci., 110 (1993).

\bibitem{Galliot25} F.~Galliot, \emph{$4$-uniform Maker--Breaker and Maker--Maker games are PSPACE-complete}, preprint, arXiv:2509.13819 (v3, 2026).

\bibitem{GGS} F.~Galliot, S.~Gravier, and I.~Sivignon, \emph{Maker--Breaker is solved in polynomial time on hypergraphs of rank $3$}, preprint, arXiv:2209.12819.

\bibitem{GalliotSenizergues} F.~Galliot and J.~S\'enizergues, \emph{Maker--Maker games of rank $4$ are PSPACE-complete}, in STACS 2026, LIPIcs, 2026.

\bibitem{GalperinWigderson} H.~Galperin and A.~Wigderson, \emph{Succinct representations of graphs}, Inform. Control, 56 (1983), pp.~183--198.

\bibitem{Godsil81} C.~D.~Godsil, \emph{Matchings and walks in graphs}, J. Graph Theory, 5 (1981), pp.~285--297.

\bibitem{GodsilBook} C.~D.~Godsil, \emph{Algebraic Combinatorics}, Chapman \& Hall, 1993.

\bibitem{Godsil95} C.~D.~Godsil, \emph{Algebraic matching theory}, Electron. J. Combin., 2 (1995), R8.

\bibitem{GHR} R.~Greenlaw, H.~J.~Hoover, and W.~L.~Ruzzo, \emph{Limits to Parallel Computation: $P$-Completeness Theory}, Oxford University Press, 1995.

\bibitem{Grundy} P.~M.~Grundy, \emph{Mathematics and games}, Eureka, 2 (1939), pp.~6--8.

\bibitem{Hajek} P.~H\'ajek, \emph{Arithmetical hierarchy and complexity of computation}, Theoret. Comput. Sci., 8 (1979), pp.~227--237.

\bibitem{HartmanisSimon} J.~Hartmanis and J.~Simon, \emph{On the power of multiplication in random access machines}, in 15th Annual Symposium on Switching and Automata Theory, IEEE, 1974, pp.~13--23.

\bibitem{HartmanisStearns} J.~Hartmanis and R.~E.~Stearns, \emph{On the computational complexity of algorithms}, Trans. Amer. Math. Soc., 117 (1965), pp.~285--306.

\bibitem{HearnDemaine} R.~A.~Hearn and E.~D.~Demaine, \emph{Games, Puzzles, and Computation}, A K Peters, 2009.

\bibitem{HeilmannLieb} O.~J.~Heilmann and E.~H.~Lieb, \emph{Theory of monomer-dimer systems}, Comm. Math. Phys., 25 (1972), pp.~190--232.

\bibitem{Immerman} N.~Immerman, \emph{Number of quantifiers is better than number of tape cells}, J. Comput. System Sci., 22 (1981), pp.~384--406.

\bibitem{Jerrum} M.~Jerrum, \emph{Two-dimensional monomer-dimer systems are computationally intractable}, J. Stat. Phys., 48 (1987), pp.~121--134.

\bibitem{JLS} C.~R.~Johnson, A.~Leal Duarte, and C.~M.~Saiago, \emph{The Parter--Wiener theorem: refinement and generalization}, SIAM J. Matrix Anal. Appl., 25 (2004), pp.~352--361.

\bibitem{Jurdzinski} M.~Jurdzi\'nski, \emph{Small progress measures for solving parity games}, in STACS 2000, Lecture Notes in Comput. Sci. 1770, Springer, 2000, pp.~290--301.

\bibitem{KannanBachem} R.~Kannan and A.~Bachem, \emph{Polynomial algorithms for computing the Smith and Hermite normal forms of an integer matrix}, SIAM J. Comput., 8 (1979), pp.~499--507.

\bibitem{Kasteleyn} P.~W.~Kasteleyn, \emph{Graph theory and crystal physics}, in Graph Theory and Theoretical Physics, F.~Harary, ed., Academic Press, 1967, pp.~43--110.

\bibitem{KissToth} V.~Kiss and L.~T\'othm\'er\'esz, \emph{Chip-firing games on Eulerian digraphs and NP-hardness of computing the rank of a divisor on a graph}, Discrete Appl. Math., 193 (2015).

\bibitem{KBKF} H.~Kleine B\"uning, M.~Karpinski, and A.~Fl\"ogel, \emph{Resolution for quantified Boolean formulas}, Inform. Comput., 117 (1995), pp.~12--18.

\bibitem{Koepke} F.~O.~Koepke, \emph{Solving Maker--Breaker games on $5$-uniform hypergraphs is PSPACE-complete}, Electron. J. Combin., 32 (2025).

\bibitem{Kozen} D.~Kozen, \emph{Indexings of subrecursive classes}, Theoret. Comput. Sci., 11 (1980), pp.~277--301.

\bibitem{KuChen} C.~Y.~Ku and W.~Chen, \emph{An analogue of the Gallai--Edmonds structure theorem for nonzero roots of the matching polynomial}, J. Combin. Theory Ser. B, 100 (2010).

\bibitem{Lehman} A.~Lehman, \emph{A solution of the Shannon switching game}, J. Soc. Indust. Appl. Math., 12 (1964), pp.~687--725.

\bibitem{LichtensteinSipser} D.~Lichtenstein and M.~Sipser, \emph{GO is polynomial-space hard}, J. ACM, 27 (1980), pp.~393--401.

\bibitem{Lipton} R.~J.~Lipton, \emph{New directions in testing}, in Distributed Computing and Cryptography, DIMACS Ser. Discrete Math. Theoret. Comput. Sci. 2, AMS, 1991, pp.~191--202.

\bibitem{Lovasz} L.~Lov\'asz, \emph{On determinants, matchings, and random algorithms}, in Fundamentals of Computation Theory (FCT 1979), Akademie-Verlag, 1979, pp.~565--574.

\bibitem{LovaszPlummer} L.~Lov\'asz and M.~D.~Plummer, \emph{Matching Theory}, North-Holland, 1986; reprinted by AMS Chelsea, 2009.

\bibitem{LozanoBalcazar} A.~Lozano and J.~L.~Balc\'azar, \emph{The complexity of graph problems for succinctly represented graphs}, in Graph-Theoretic Concepts in Computer Science (WG 1989), Lecture Notes in Comput. Sci. 411, Springer, 1990.

\bibitem{LFKN} C.~Lund, L.~Fortnow, H.~Karloff, and N.~Nisan, \emph{Algebraic methods for interactive proof systems}, J. ACM, 39 (1992), pp.~859--868.

\bibitem{MetelmannClerk} A.~Metelmann and A.~A.~Clerk, \emph{Nonreciprocal photon transmission and amplification via reservoir engineering}, Phys. Rev. X, 5 (2015), 021025.

\bibitem{MikhalkinZharkov} G.~Mikhalkin and I.~Zharkov, \emph{Tropical curves, their Jacobians and theta functions}, in Curves and Abelian Varieties, Contemp. Math. 465, AMS, 2008, pp.~203--230.

\bibitem{MiklosKresz} I.~Mikl\'os and M.~Kr\'esz, \emph{Counting maximum matchings in planar graphs is hard}, arXiv:2001.01493 (2020), withdrawn 2021.

\bibitem{Morris} F.~L.~Morris, \emph{Playing disjunctive sums is polynomial space complete}, Int. J. Game Theory, 10 (1981), pp.~195--205.

\bibitem{PapadimitriouYannakakis} C.~H.~Papadimitriou and M.~Yannakakis, \emph{A note on succinct representations of graphs}, Inform. Control, 71 (1986), pp.~181--185.

\bibitem{PolyaSchur} G.~P\'olya and J.~Schur, \emph{\"Uber zwei Arten von Faktorenfolgen in der Theorie der algebraischen Gleichungen}, J. Reine Angew. Math., 144 (1914), pp.~89--113.

\bibitem{RahmanWatson} M.~L.~Rahman and T.~Watson, \emph{$6$-uniform Maker--Breaker game is PSPACE-complete}, in STACS 2021, LIPIcs 187, 2021.

\bibitem{Reichardt} B.~W.~Reichardt, \emph{Reflections for quantum query algorithms}, in SODA 2011, SIAM, 2011.

\bibitem{ReichardtSpalek} B.~W.~Reichardt and R.~\v{S}palek, \emph{Span-program-based quantum algorithm for evaluating formulas}, Theory Comput., 8 (2012), pp.~291--319.

\bibitem{Reingold} O.~Reingold, \emph{Undirected connectivity in log-space}, J. ACM, 55 (2008), Art.~17.

\bibitem{Reisch} S.~Reisch, \emph{Hex ist PSPACE-vollst\"andig}, Acta Inform., 15 (1981), pp.~167--191.

\bibitem{RenaultSchmidt} G.~Renault and S.~Schmidt, \emph{On the complexity of the mis\`ere version of three games played on graphs}, Theoret. Comput. Sci., 595 (2015), pp.~159--167.

\bibitem{Robson} J.~M.~Robson, \emph{$N$ by $N$ checkers is Exptime complete}, SIAM J. Comput., 13 (1984), pp.~252--267.

\bibitem{SaksWigderson} M.~Saks and A.~Wigderson, \emph{Probabilistic Boolean decision trees and the complexity of evaluating game trees}, in FOCS 1986, IEEE, 1986, pp.~29--38.

\bibitem{Santha} M.~Santha, \emph{On the Monte Carlo Boolean decision tree complexity of read-once formulae}, Random Structures Algorithms, 6 (1995), pp.~75--87.

\bibitem{Savitch} W.~J.~Savitch, \emph{Relationships between nondeterministic and deterministic tape complexities}, J. Comput. System Sci., 4 (1970), pp.~177--192.

\bibitem{SchaeferSAT} T.~J.~Schaefer, \emph{The complexity of satisfiability problems}, in STOC 1978, ACM, 1978, pp.~216--226.

\bibitem{SchaeferGames} T.~J.~Schaefer, \emph{On the complexity of some two-person perfect-information games}, J. Comput. System Sci., 16 (1978), pp.~185--225.

\bibitem{Shamir} A.~Shamir, \emph{IP = PSPACE}, J. ACM, 39 (1992), pp.~869--877.

\bibitem{Siegel} A.~N.~Siegel, \emph{Combinatorial Game Theory}, Grad. Stud. Math. 146, AMS, 2013.

\bibitem{Snir} M.~Snir, \emph{Lower bounds on probabilistic linear decision trees}, Theoret. Comput. Sci., 38 (1985), pp.~69--82.

\bibitem{Sprague} R.~P.~Sprague, \emph{\"Uber mathematische Kampfspiele}, T\^ohoku Math. J., 41 (1935--36), pp.~438--444.

\bibitem{StockmeyerChandra} L.~J.~Stockmeyer and A.~K.~Chandra, \emph{Provably difficult combinatorial games}, SIAM J. Comput., 8 (1979), pp.~151--174.

\bibitem{STV} M.~Sudan, L.~Trevisan, and S.~Vadhan, \emph{Pseudorandom generators without the XOR lemma}, J. Comput. System Sci., 62 (2001), pp.~236--266.

\bibitem{TemperleyFisher} H.~N.~V.~Temperley and M.~E.~Fisher, \emph{Dimer problem in statistical mechanics: an exact result}, Philos. Mag., 6 (1961), pp.~1061--1063.

\bibitem{TrevisanVadhan} L.~Trevisan and S.~Vadhan, \emph{Pseudorandomness and average-case complexity via uniform reductions}, Comput. Complexity, 16 (2007), pp.~331--364.

\bibitem{Tutte} W.~T.~Tutte, \emph{The factorization of linear graphs}, J. London Math. Soc., 22 (1947), pp.~107--111.

\bibitem{Vadhan} S.~P.~Vadhan, \emph{The complexity of counting in sparse, regular, and planar graphs}, SIAM J. Comput., 31 (2001), pp.~398--427.

\bibitem{Valiant} L.~G.~Valiant, \emph{The complexity of computing the permanent}, Theoret. Comput. Sci., 8 (1979), pp.~189--201.

\bibitem{Wythoff} W.~A.~Wythoff, \emph{A modification of the game of Nim}, Nieuw Arch. Wisk., 7 (1907), pp.~199--202.

\bibitem{Zhuk24} D.~Zhuk, \emph{$\Pi_2^P$ vs PSpace dichotomy for the quantified constraint satisfaction problem}, in FOCS 2024, IEEE, 2024, pp.~560--572.

\bibitem{ZhukMartin} D.~Zhuk and B.~Martin, \emph{QCSP monsters and the demise of the Chen conjecture}, J. ACM, 69 (2022), Art.~35.

\end{thebibliography}

\end{document}
